\documentclass[11pt,a4paper,oneside]{book}
\usepackage[margin=20mm,top=22mm,bottom=22mm]{geometry}
\usepackage{fontspec}
\usepackage{amsmath}
\usepackage[math-style=ISO]{unicode-math}
\usepackage{polyglossia}
\setdefaultlanguage{ukrainian}
\setotherlanguage{english}
\directlua{
  luaotfload.add_fallback("fallbacks", {"DejaVu Sans:mode=harf;", "FreeSerif:mode=harf;", "DejaVu Math TeX Gyre:mode=harf;"})
}
\setmainfont{DejaVu Serif}[RawFeature={fallback=fallbacks}]
\setsansfont{DejaVu Sans}[RawFeature={fallback=fallbacks}]
\setmonofont{DejaVu Sans Mono}[Scale=0.85,RawFeature={fallback=fallbacks}]
\setmathfont{DejaVu Math TeX Gyre}
\usepackage{microtype}
\usepackage{needspace}
\usepackage{booktabs,longtable,array}
\usepackage{enumitem}
\setlist{nosep,leftmargin=*}
\usepackage{xcolor}
\definecolor{accent}{HTML}{1F4E79}
\definecolor{warncol}{HTML}{B45309}
\usepackage{titlesec}
\titleformat{\chapter}[hang]{\Large\bfseries\sffamily\color{accent}}{\thechapter.}{0.6em}{}
\titlespacing*{\chapter}{0pt}{-10pt}{10pt}
\titleformat{\section}{\normalsize\bfseries\sffamily\color{accent}}{}{0pt}{}
\titlespacing*{\section}{0pt}{8pt}{3pt}
\titleformat{\part}[display]{\centering\Huge\bfseries\sffamily\color{accent}}{\partname~\thepart}{12pt}{}
\usepackage{fancyhdr}
\pagestyle{fancy}\fancyhf{}
\fancyhead[L]{\small\sffamily\leftmark}\fancyhead[R]{\small\sffamily\thepage}
\renewcommand{\headrulewidth}{0.3pt}
\usepackage{fancyvrb}
\DefineVerbatimEnvironment{layout}{Verbatim}{fontsize=\footnotesize,frame=single,framesep=2mm,rulecolor=\color{accent!40}}
\usepackage[hidelinks,unicode]{hyperref}
\newcommand{\pen}{{\sffamily\color{accent}✎}}
\newcommand{\warn}{{\sffamily\color{warncol}⚠}}
\newcommand{\ok}{{\sffamily\color{green!45!black}✔}}
\newcommand{\spark}{{\sffamily\color{warncol}⚡}}
\newcommand{\book}{}
\newcommand{\srcline}[1]{\par\noindent{\small\color{black!70}\textit{Джерело:} #1}\par\medskip}
\setlength{\parindent}{0pt}\setlength{\parskip}{3pt}
% Mixed Ukrainian/English text: Ukrainian hyphenation patterns can't break English words,
% so give the line breaker room instead of producing overfull lines.
\setlength{\emergencystretch}{3em}\tolerance=2000
\setlength{\tabcolsep}{3.5pt}
\raggedbottom


\begin{document}

\frontmatter


\begin{titlepage}
\centering\sffamily
\vspace*{40mm}
{\Huge\bfseries\color{accent} Загальний курс фізики\par}
\vspace{6mm}
{\LARGE Каркаси конспектів лекцій\par}
\vspace{12mm}
{\large Електростатика і постійний струм\\[2mm] Молекулярна фізика і термодинаміка\par}
\vfill
{\normalsize За підручником: Кучерук І.\,М., Горбачук І.\,Т., Луцик П.\,П.\\
\textit{Загальний курс фізики}, т.~1 (1999), т.~2 (2001). — К.: Техніка.\par}
\vspace{8mm}
{\small Згенеровано з вікі проєкту, \today\par}
\end{titlepage}

\chapter*{Як користуватися}
\addcontentsline{toc}{chapter}{Як користуватися}\markboth{Як користуватися}{}
Кожен розділ — це \textbf{каркас} однієї рукописної сторінки конспекту, а не готовий конспект. Він містить:
\begin{itemize}
\item \textbf{План сторінки} — розкладка сторінки на рамки [A], [B], …;
\item точні терміни й означення з підручника (у «лапках»), формули з номерами рівнянь книги, сторінки;
\item \pen{} — місця, де треба намалювати схему чи сформулювати думку власними словами;
\item \warn{} — розбіжності між підручником і програмою, помилки в книзі, неперевірені місця;
\item \ok{} — формули звірено зі сканом сторінки; \spark{} — важливий або парадоксальний висновок;
\item \textbf{Самоперевірка} — 2–4 запитання для самоперевірки.
\end{itemize}
Усі пояснення — українською; поруч із ключовими термінами в дужках подано англійські відповідники «(англ. …)». Терміни й цитати — дослівно з підручника. Посилання «т.~2, с.~38» — друковані сторінки підручника.


\tableofcontents

\mainmatter

\part{Електростатика і постійний струм}

\chapter*{Програма: 23 питання}\addcontentsline{toc}{chapter}{Програма: 23 питання}\markboth{Програма: 23 питання}{Програма: 23 питання}


Джерело: т. 2, розд. 1–2.

Завдання — рукописний конспект кожного питання. Статус:  каркас готовий · \ok{} формули звірено зі сканом ·  написано від руки (позначає автор конспекту).

{\footnotesize
\begin{longtable}{>{\raggedright\arraybackslash}p{0.040\linewidth}>{\raggedright\arraybackslash}p{0.597\linewidth}>{\raggedright\arraybackslash}p{0.092\linewidth}>{\raggedright\arraybackslash}p{0.122\linewidth}>{\raggedright\arraybackslash}p{0.069\linewidth}}
\toprule
\textbf{№} & \textbf{Питання (як у програмі)} & \textbf{Сторінка} & \textbf{§} & \textbf{Статус} \\
\midrule
\endhead
1 & Електричний заряд та його властивості & 01 & 1.1 &  \\
2 & Закон Кулона. Диполь в неоднорідному полі & 02 & 1.2–1.3, 1.6 &  \\
3 & Напруженість електричного поля E & 03 & 1.4–1.5 &  \\
4 & Потенціал електричного поля & 04 & 1.10–1.11 &  \\
5 & Диполь в однорідному електричному полі & 05 & 1.6 &  \\
6 & Потік вектора напруженості електричного поля E. Теорема Гауса & 06 & 1.7 &  \\
7 & Дивергенція вектора E & 07 & 1.8 &  \\
8 & Зв'язок між напруженістю та потенціалом & 08 & 1.11 &  \\
9 & Теорема Стокса & 09 & 1.10 &  \\
10 & Поле зарядженої площини & 10–12 & 1.7 &  \\
11 & Поле двох заряджених площин (теорема Гауса) & 10–12 & 1.7 &  \\
12 & Поле сферичної поверхні & 10–12 & 1.7, 1.11 &  \\
13 & Теорема Остроградського-Гауса & 13 & 1.8 &  \\
14 & Провідник у зовнішньому електричному полі & 14 & 1.12–1.13 &  \\
15 & Конденсатори та їх ємність. Ємність плоского, циліндричного та сферичного конденсаторів & 15 & 1.14 &  \\
16 & З'єднання конденсаторів & 16 & 1.14 &  \\
17 & Енергія зарядженого провідника & 17 & 1.25–1.26 &  \\
18 & Енергія електростатичного поля & 18 & 1.26 &  \\
19 & Електричний струм. Постійний електричний струм. Закони Ома & 19 & 2.1–2.3 &  \\
20 & Сила струму, густина струму. Електрорушійна сила & 20 & 2.1, 2.3 &  \\
21 & Поняття опору. Потужність & 21 & 2.2, 2.4 &  \\
22 & Перший та другий закон Кірхгофа & 22 & 2.6 &  \\
23 & Закон Джоуля-Ленца & 23 & 2.5 &  \\
\bottomrule
\end{longtable}
}

\needspace{5\baselineskip}\section*{Що перетинається}
\begin{itemize}
\item 2 і 5 мають спільний диполь. Означте його один раз (у 5) і посилайтеся з 2.
\item 6, 7, 13 утворюють один ланцюжок: Гаусс (інтегральна форма) → div (диференціальна) → Остроградський–Гаусс (математичний місток).
\item 19, 20, 21 перетинаються. Пропонований поділ: 19 — закон Ома в усіх формах, 20 — I, j, ЕРС, 21 — R, ρ, P.
\item 17 і 18 дають те саме W — один раз через заряди, другий — через поле.
\end{itemize}

\chapter*{1. Електричний заряд та його властивості}\addcontentsline{toc}{chapter}{1. Електричний заряд та його властивості}\markboth{1. Електричний заряд та його властивості}{1. Електричний заряд та його властивості}

\srcline{\book{} Кучерук т. 2, \textbf{§ 1.1 «Електричний заряд», с. 14–21} (розвороти 8–11). \ok{} Усі формули звірено зі сканом.}
Це каркас рукописної сторінки: точні терміни, формули й структура. Формулювання на сторінці — ваші.

\par\noindent{\color{accent!30}\rule{\linewidth}{0.4pt}}\par

\needspace{5\baselineskip}\section*{План сторінки (одна сторона A4)}

\begin{layout}
┌──────────────────────────────────────────────────────┐
│ 1. ЕЛЕКТРИЧНИЙ ЗАРЯД                                  │
│ [0] Що це? (означення + 1.1)                          │
├──────────────────────────┬───────────────────────────┤
│ [1] Два види (±)         │ [2] Квантування            │
│                          │     + схема Міллікена      │
├──────────────────────────┼───────────────────────────┤
│ [3] Інваріантність       │ [4] Закон збереження (1.2)│
├──────────────────────────┴───────────────────────────┤
│ [5] Неперервний розподіл: ρ, σ, τ  (1.3)–(1.5)        │
│ [низ] Носії заряду: e⁻, p⁺, іони · електризація        │
└──────────────────────────────────────────────────────┘
\end{layout}

\par\noindent{\color{accent!30}\rule{\linewidth}{0.4pt}}\par

\needspace{5\baselineskip}\section*{[0] Що таке заряд (с. 15)}

\begin{itemize}
\item Означення книги: «Електричний заряд є \textbf{внутрішньою характеристикою} деяких елементарних матеріальних частинок, яка проявляється в їхній взаємодії на відстані або під час дії на них зовнішніх електромагнітних полів.»
\item Головна думка, яку варто записати своїми словами: заряд — це \textbf{і джерело, і об'єкт дії} поля («як джерело, так і об'єкт дії»).
\item Як вимірюють заряд (с. 14): почергово вміщують дві частинки в ту саму точку M поля:
\end{itemize}

\[\frac{F_1}{F_2} = \frac{q_1}{q_2} \qquad (1.1)\]

Відношення не залежить ні від поля, ні від точки M — лише від самих частинок.
\begin{itemize}
\item Порівняння з масою (с. 15): маса має один знак і лише притягує; заряд має два знаки, тож заряди можуть \textbf{компенсувати} один одного.
\end{itemize}

\needspace{5\baselineskip}\section*{[1] Два види зарядів (с. 15)}

\begin{itemize}
\item Однойменні заряди відштовхуються, різнойменні — притягуються (англ. like / unlike charges).
\item Угода про знак: \textbf{позитивний} — скло, натерте шовком; \textbf{негативний} — ебоніт (бурштин), натертий хутром.
\item Угода довільна. Суть у тому, що видів саме два.
\item \pen{} Ваш рядок: чому тіла зазвичай нейтральні?
\end{itemize}

\needspace{5\baselineskip}\section*{[2] Дискретність: закон квантування електричного заряду (с. 16–19)}

\begin{itemize}
\item Термін книги (курсивом): \textbf{\textit{закон квантування електричного заряду}} (англ. charge quantization). Вільні заряди складаються з однакових дискретних порцій, що дорівнюють заряду електрона.
\end{itemize}

\[q = n|e|,\quad n = \pm1, \pm2, \ldots \qquad |e| = 1{,}602\cdot10^{-19}\ \text{\text{Кл}}\]

\begin{itemize}
\item \textbf{Дослід Міллікена (рис. 1.1, с. 18)}: намалюйте дві горизонтальні пластини, краплину олії та стрілки qE↑, $F_{A}$↑, mg↓, av↕.
\begin{itemize}
\item Без поля: краплина падає рівномірно, mg = av + $F_{A}$, звідси знаходять \textit{a}.
\item З полем E: краплина рівномірно піднімається, \textbf{qE + $F_{A}$ = mg + av′}, звідси знаходять \textit{q}.
\item Рентгенівське опромінення змінює заряд: Δq = q′ − q = n|e| \textbf{завжди}.
\end{itemize}
\item Саме квантування — це \textbf{постулат}: теорія його не пояснює (с. 19).
\item Додаткові факти (необов'язково): e/m = 1,7588·10¹¹ Кл/кг (Томсон, 1897); $m_{e}$ = 9,109·10⁻³¹ кг; кварки ±e/3, 2e/3, у вільному стані не виявлені.
\end{itemize}

\needspace{5\baselineskip}\section*{[3] Інваріантність заряду (с. 20)}

\begin{itemize}
\item Значення заряду \textbf{не залежить} від інерціальної системи відліку чи швидкості носія (англ. charge invariance).
\item Доказ — нейтральність атомів. Електрони в He рухаються приблизно вдвічі швидше, ніж у H₂, проте обидва нейтральні, отже заряд не залежить від швидкості. Перевірено до v ≈ 0,5c.
\item \pen{} Порівняйте самі: маса \textit{змінюється} зі швидкістю (теорія відносності), а заряд — ні.
\end{itemize}

\needspace{5\baselineskip}\section*{[4] Закон збереження заряду (с. 20–21)}

\begin{itemize}
\item \textbf{\textit{Замкнена система}} (англ. closed system) — така, крізь межі якої не проникають заряджені частинки.
\item «Повний електричний заряд замкненої системи є величиною постійною і дорівнює алгебричній сумі її позитивних та негативних зарядів»:
\end{itemize}

\[q = q_1 + q_2 + \ldots + q_n = \sum_{i=1}^{n} q_i = \text{const} \qquad (1.2)\]

\begin{itemize}
\item Наслідок: заряди народжуються й зникають \textbf{лише парами ±} (народження пари e⁻e⁺ з γ-кванта, анігіляція).
\item Електризація тертям лише \textbf{перерозподіляє} заряд, сумарний заряд не змінюється (с. 16). Гарний приклад для малюнка.
\end{itemize}

\needspace{5\baselineskip}\section*{[5] Неперервний розподіл: густини заряду (с. 21)}

\begin{itemize}
\item Навіщо: при 1 А крізь переріз проходить ≈ 6·10¹⁸ елементарних зарядів за секунду, тож дискретність непомітна і заряд вважають неперервним.
\end{itemize}

{\small
\begin{longtable}{>{\raggedright\arraybackslash}p{0.341\linewidth}>{\raggedright\arraybackslash}p{0.520\linewidth}>{\raggedright\arraybackslash}p{0.091\linewidth}}
\toprule
\textbf{Термін (курсив книги)} & \textbf{Формула} & \textbf{Одиниця} \\
\midrule
\endhead
\textbf{\textit{об'ємна густина заряду}} (англ. volume charge density) & $\rho = \lim\limits_{\Delta V\to0}\dfrac{\Delta q}{\Delta V} = \dfrac{dq}{dV}$ (1.3) & Кл/м³ \\
\textbf{\textit{поверхнева густина заряду}} (англ. surface charge density) & $\sigma = \lim\limits_{\Delta S\to0}\dfrac{\Delta q}{\Delta S} = \dfrac{dq}{dS}$ (1.4) & Кл/м² \\
\textbf{\textit{лінійна густина}} (англ. linear charge density) & $\tau = \lim\limits_{\Delta l\to0}\dfrac{\Delta q}{\Delta l} = \dfrac{dq}{dl}$ (1.5) & Кл/м \\
\bottomrule
\end{longtable}
}

\begin{itemize}
\item Застереження книги: ΔV, ΔS, Δl малі, але кожен містить багато елементарних зарядів.
\item σ знадобиться в темах 10–12 (заряджена площина, дві площини, сфера), τ — для зарядженої нитки.
\end{itemize}

\needspace{5\baselineskip}\section*{[низ] Носії заряду (с. 16)}

\begin{itemize}
\item Стабільні носії: \textbf{електрон} (−e), \textbf{протон} (+e, за модулем дорівнює заряду електрона з дуже високою відносною точністю, с. 19). Іони утворюються, коли атом втрачає або набуває електрони.
\end{itemize}

\par\noindent{\color{accent!30}\rule{\linewidth}{0.4pt}}\par

\needspace{5\baselineskip}\section*{Пропустити (відступи, що не стосуються теми)}

Історія (Гільберт, Франклін, Дюфе, Стоні), класичний радіус електрона r₀ = e²/(4πε₀m₀c²) = 2,82·10⁻¹⁵ м, спін, парадокс «нескінченної енергії» електрона, час життя нейтрона. Додавайте лише якщо лектор цікавиться історією.

\needspace{5\baselineskip}\section*{Самоперевірка}

\begin{enumerate}
\item Чому нейтральність атома доводить інваріантність заряду?
\item Що доводить дослід Міллікена і що він вимірює?
\item Чому тертям електризуються діелектрики, а метали — ні? (с. 16)
\item Яка властивість є \textit{постулатом}, а які — дослідні факти?
\end{enumerate}

\chapter*{2. Закон Кулона. Диполь в неоднорідному полі}\addcontentsline{toc}{chapter}{2. Закон Кулона. Диполь в неоднорідному полі}\markboth{2. Закон Кулона. Диполь в неоднорідному полі}{2. Закон Кулона. Диполь в неоднорідному полі}

\srcline{\book{} т. 2, \textbf{§ 1.2} (с. 21–26), \textbf{§ 1.3} (с. 26), \textbf{§ 1.6} (с. 32–35). \ok{} Формули звірено зі сканом.}

\needspace{5\baselineskip}\section*{План сторінки}
\begin{layout}
┌ [A] Закон Кулона: формулювання + (1.6)/(1.7) + рис. 1.2 крутильні ваги ┐
├ [B] Векторна форма (1.8)  │ [C] k, ε₀, одиниці                         ┤
├ [D] Принцип суперпозиції (1.10) + рис. 1.4                              ┤
├ [E] Середовище: ε (§1.3)  │ [F] Диполь у неоднорідному полі (1.26)       ┤
└ низ: межі застосування закону                                            ┘
\end{layout}

\needspace{5\baselineskip}\section*{[A] Закон Кулона (с. 22)}
\begin{itemize}
\item Формулювання книги (варто вивчити дослівно): «сила взаємодії F двох \textbf{нерухомих точкових зарядів у вакуумі} прямо пропорційна добуткові кількостей електрики цих зарядів q₁ і q₂, обернено пропорційна квадратові відстані r між ними і напрямлена вздовж прямої, яка сполучає ці заряди» (англ. Coulomb's law)
\item $F = k\dfrac{q_1q_2}{r^2}$ (1.6) → у СІ: $F = \dfrac{1}{4\pi\varepsilon_0}\dfrac{q_1q_2}{r^2}$ (1.7)
\item \textbf{Точковий заряд} (англ. point charge): розміри тіл ≪ відстані між ними (с. 24).
\item Угода про знак: відштовхування F > 0, притягання F < 0.
\item \pen{} Намалюйте крутильні ваги Кулона (рис. 1.2): нитка, коромисло, кульки a і d.
\end{itemize}

\needspace{5\baselineskip}\section*{[B] Векторна форма (с. 24)}
\begin{itemize}
\item $\vec F_{12} = \dfrac{1}{4\pi\varepsilon_0}\dfrac{q_1q_2}{r_{12}^3}\vec r_{12}$ (1.8), причому $\vec F_{12} = -\vec F_{21}$ (1.9)
\item Сили \textbf{центральні} (англ. central forces): діють уздовж прямої, що сполучає заряди.
\end{itemize}

\needspace{5\baselineskip}\section*{[C] Коефіцієнт і одиниці (с. 23–24)}
\begin{itemize}
\item $k = \dfrac{1}{4\pi\varepsilon_0} = 9\cdot10^9\ \dfrac{\text{\text{Н}·\text{м}}^2}{\text{\text{Кл}}^2}$; \textbf{електрична стала} (англ. electric constant) $\varepsilon_0 = 8{,}85\cdot10^{-12}$ Ф/м.
\item 1 Кл = 1 А·1 с; основна одиниця — ампер, тож кулон — похідна одиниця.
\item «Раціоналізована форма» вносить 4π у k, щоб він зник із формул, які використовують найчастіше.
\item Необов'язково: у СГСЕ k = 1, а 1 Кл = 3·10⁹ одиниць СГСЕ.
\end{itemize}

\needspace{5\baselineskip}\section*{[D] Принцип суперпозиції (с. 25)}
\begin{itemize}
\item Сила взаємодії двох зарядів \textbf{не залежить від наявності інших зарядів}; результуюча сила — векторна сума (англ. superposition principle):
$\vec F = \sum_i \vec F_i = \sum_i \dfrac{1}{4\pi\varepsilon_0}\dfrac{q_iq_0}{r_{i0}^3}\vec r_{i0}$ (1.10)
\item Протяжні тіла подумки ділять на точкові заряди.
\end{itemize}

\needspace{5\baselineskip}\section*{[E] Вплив середовища (§ 1.3, с. 26)}
\begin{itemize}
\item \textbf{Діелектрична проникність} (англ. permittivity) $\varepsilon = F_0/F$ — у скільки разів сила в середовищі слабша, ніж у вакуумі (безрозмірна).
\item $F = \dfrac{1}{4\pi\varepsilon_0}\dfrac{q_1q_2}{\varepsilon r^2}$
\item Причина: середовище зазнає \textbf{поляризації} (англ. polarization) і частково нейтралізує заряди.
\end{itemize}

\needspace{5\baselineskip}\section*{[F] Диполь у неоднорідному полі (с. 35)}
\begin{itemize}
\item У неоднорідному полі дві сили різні, тож крім моменту сил з'являється ще й рівнодійна:
$F = p\,\dfrac{\partial E}{\partial x}$ (1.26), диполь уздовж силової лінії.
\item Напрям — \textbf{у бік зростання напруженості}: диполь утягується в область сильнішого поля.
\item \pen{} Приклад для запису: чому наелектризований гребінець притягує клаптики паперу (індуковані диполі).
\item Необов'язково: приклад 1.1, два диполі на одній прямій, $F = -\dfrac{6p_1p_2}{4\pi\varepsilon_0\varepsilon r^4}$.
\item Саме означення диполя — у темі 5 (запишіть його там або тут, але один раз).
\end{itemize}

\needspace{5\baselineskip}\section*{Низ: межі застосування (с. 26)}
\begin{itemize}
\item Перевірено від ≈ 10⁻¹⁷ м (розсіяння електронів) до ≈ 10⁷ м (супутники).
\end{itemize}

\needspace{5\baselineskip}\section*{Самоперевірка}
\begin{enumerate}
\item Чому закон сформульовано для \textit{точкових} і \textit{нерухомих} зарядів?
\item Що додає принцип суперпозиції до самого закону Кулона?
\item Чому диполь в однорідному полі повертається, але не рухається поступально?
\end{enumerate}

\chapter*{3. Напруженість електричного поля E}\addcontentsline{toc}{chapter}{3. Напруженість електричного поля E}\markboth{3. Напруженість електричного поля E}{3. Напруженість електричного поля E}

\srcline{\book{} т. 2, \textbf{§ 1.4} (с. 27), \textbf{§ 1.5} (с. 28–31). \ok{} Формули звірено зі сканом.}

\needspace{5\baselineskip}\section*{План сторінки}
\begin{layout}
┌ [A] Що таке поле (§1.4): близькодія vs далекодія ┐
├ [B] Означення E (1.11) + пробний заряд             ┤
├ [C] Поле точкового заряду (1.12) │ [D] Суперпозиція (1.13)–(1.14) ┤
├ [E] Неперервні заряди (1.15)–(1.18)                 ┤
└ [F] Силові лінії: рис. 1.6–1.8; однорідне поле       ┘
\end{layout}

\needspace{5\baselineskip}\section*{[A] Електричне поле (с. 27)}
\begin{itemize}
\item \textbf{Електричне поле} (англ. electric field): «матеріальна складова електромагнітного поля, яке діє на заряд, зумовлене зарядами або змінним у часі магнітним полем».
\item Основна ознака поля: \textbf{на будь-який внесений у нього заряд діє сила}.
\item Дві концепції: \textbf{далекодія} (англ. action at a distance: миттєво, без носія) — відкинута; \textbf{близькодія} (англ. short-range action: через поле, зі скінченною швидкістю) — прийнята.
\item В електростатиці обидві дають однакові результати; різняться лише для рухомих зарядів.
\end{itemize}

\needspace{5\baselineskip}\section*{[B] Означення напруженості (с. 28)}
\begin{itemize}
\item \textbf{Пробний заряд} (англ. test charge): точковий, позитивний, досить малий, щоб не спотворювати поле.
\item «Напруженість E у даній точці електростатичного поля чисельно дорівнює силі, з якою поле діє на \textbf{одиничний позитивний пробний заряд}, вміщений у цю точку» (англ. electric field strength)
\item $\vec E = \vec F / q_0$ (1.11) — \textbf{вектор} і \textbf{силова характеристика поля} (англ. force characteristic of the field).
\item Одиниці: Н/Кл = В/м.
\item E — \textbf{локальна} (англ. local) величина: кожна точка має своє значення.
\end{itemize}

\needspace{5\baselineskip}\section*{[C] Поле точкового заряду (с. 28)}
\begin{itemize}
\item $\vec E = \dfrac{1}{4\pi\varepsilon_0}\dfrac{q}{r^3}\vec r$ (1.12), тобто $E = \dfrac{q}{4\pi\varepsilon_0 r^2}$
\item Поле \textbf{центрально-симетричне} (англ. centrally symmetric).
\item Взаємодія у два етапи (с. 29): q₁ створює поле E₁₂, а поле діє на q₂ із силою $\vec F_{12} = q_2\vec E_{12}$.
\end{itemize}

\needspace{5\baselineskip}\section*{[D] Принцип суперпозиції полів (с. 29–30)}
\begin{itemize}
\item $\vec E = \sum_{i=1}^n \vec E_i$ (1.13), $\vec E = \sum_i \dfrac{1}{4\pi\varepsilon_0}\dfrac{q_i}{r_i^3}\vec r_i$ (1.14)
\item \textbf{Польове формулювання} (англ. field formulation): поле системи дорівнює сумі полів кожного заряду «за відсутності всіх інших».
\item «Основне завдання електростатики»: знайти \textbf{E} у кожній точці за заданим розподілом зарядів (с. 29).
\end{itemize}

\needspace{5\baselineskip}\section*{[E] Неперервний розподіл (с. 30–31)}
\begin{itemize}
\item $dq = \rho\,dV$ → $\vec E = \dfrac{1}{4\pi\varepsilon_0}\int_V \dfrac{\rho\,\vec r}{r^3}dV$ (1.15); аналогічно з σdS і τdl (1.18).
\item Зв'язок із темою 1: ρ, σ, τ.
\end{itemize}

\needspace{5\baselineskip}\section*{[F] Силові лінії (с. 31)}
\begin{itemize}
\item \textbf{Силова лінія} (англ. field line) — лінія, дотична до якої в кожній точці напрямлена вздовж E.
\item Починаються на + зарядах і закінчуються на − (або йдуть у нескінченність); \textbf{ніколи не перетинаються}.
\item Густота ліній (на одиницю площі) ∝ E.
\item \textbf{Однорідне поле} (англ. uniform field): E скрізь однакове, лінії паралельні й рівновіддалені.
\item \pen{} Намалюйте три картинки: окремий +, окремий −, пара +/− (рис. 1.6–1.8).
\end{itemize}

\needspace{5\baselineskip}\section*{Самоперевірка}
\begin{enumerate}
\item Чому пробний заряд має бути малим?
\item Чому силові лінії не можуть перетинатися?
\item Що первинне в концепції близькодії: сила чи поле?
\end{enumerate}

\chapter*{4. Потенціал електричного поля}\addcontentsline{toc}{chapter}{4. Потенціал електричного поля}\markboth{4. Потенціал електричного поля}{4. Потенціал електричного поля}

\srcline{\book{} т. 2, \textbf{§ 1.10} (с. 44–45, робота), \textbf{§ 1.11} (с. 48–56). \ok{} Формули звірено зі сканом.}

\needspace{5\baselineskip}\section*{План сторінки}
\begin{layout}
┌ [A] Робота поля не залежить від шляху (1.44)–(1.46) ┐
├ [B] Означення φ (1.54)  │ [C] φ точкового заряду (1.55)  ┤
├ [D] Суперпозиція (1.56) │ [E] Різниця потенціалів, напруга (1.59), еВ ┤
└ [F] Еквіпотенціальні поверхні: рис., E ⟂ поверхні     ┘
\end{layout}

\needspace{5\baselineskip}\section*{[A] Робота сил поля (с. 44–45)}
\begin{itemize}
\item Переміщення q₀ у полі заряду q з точки 1 у точку 2: $A_{12} = \dfrac{qq_0}{4\pi\varepsilon_0}\left(\dfrac1{r_1}-\dfrac1{r_2}\right)$ (1.45)
\item Робота \textbf{не залежить від форми шляху}, отже поле \textbf{потенціальне (консервативне)} (англ. conservative field).
\item По замкненому контуру робота дорівнює 0: $\oint_L \vec E\,d\vec l = 0$ (1.46) — \textbf{циркуляція вектора E} (англ. circulation). Докладніше — у темі 9.
\end{itemize}

\needspace{5\baselineskip}\section*{[B] Означення потенціалу (с. 48–49)}
\begin{itemize}
\item $A_{12} = W_1 - W_2$: робота дорівнює \textbf{зменшенню} потенціальної енергії (1.53).
\item $\varphi_1 = \dfrac{A_{1\infty}}{q_0} = \int_1^\infty E\,dl\cos\alpha$ (1.54)
\item Означення книги: «фізична величина, що чисельно дорівнює роботі, яку виконують сили електричного поля при переміщенні \textbf{пробного одиничного заряду} із заданої точки \textbf{в нескінченність}» (англ. electric potential)
\item Рівнозначне означення: \textbf{відношення потенціальної енергії заряду до значення цього заряду}, $\varphi = W/q_0$.
\item φ — \textbf{скалярна} (англ. scalar) величина й \textbf{енергетична характеристика} поля (англ. energy characteristic). 1 В = 1 Дж/1 Кл.
\item Нуль відліку: φ(∞) = 0 у теорії, φ(Землі) = 0 на практиці.
\end{itemize}

\needspace{5\baselineskip}\section*{[C] Потенціал точкового заряду (с. 50)}
\begin{itemize}
\item $\varphi = \dfrac{1}{4\pi\varepsilon_0}\dfrac{q}{r}$ (1.55); додатний для +q, від'ємний для −q.
\end{itemize}

\needspace{5\baselineskip}\section*{[D] Суперпозиція (с. 50)}
\begin{itemize}
\item $\varphi = \dfrac{1}{4\pi\varepsilon_0}\sum_i \dfrac{q_i}{r_i}$ (1.56) — \textbf{алгебрична}, а не векторна сума. У цьому головна практична перевага φ над E.
\item Неперервний розподіл: $\varphi = \dfrac{1}{4\pi\varepsilon_0}\int \dfrac{\rho\,dV}{r}$ (1.57).
\end{itemize}

\needspace{5\baselineskip}\section*{[E] Різниця потенціалів = напруга (с. 50–51)}
\begin{itemize}
\item φ визначений лише з точністю до адитивної сталої, тож фізичний зміст має \textbf{Δφ}:
$\varphi_1 - \varphi_2 = \dfrac{A_{12}}{q_0}$ (1.59), позначають \textbf{U} (англ. voltage), одиниця — вольт.
\item \textbf{Електрон-вольт}: 1 еВ = 1,6·10⁻¹⁹ Дж.
\end{itemize}

\needspace{5\baselineskip}\section*{[F] Еквіпотенціальні поверхні (с. 51)}
\begin{itemize}
\item Поверхня, на якій φ = const (англ. equipotential surface). Переміщення заряду вздовж неї не потребує роботи, отже \textbf{E ⟂ поверхні}.
\item Поверхня провідника еквіпотенціальна.
\item \pen{} Малюнок: точковий заряд, концентричні кола (φ) і радіальні стрілки (E).
\end{itemize}

\needspace{5\baselineskip}\section*{Самоперевірка}
\begin{enumerate}
\item Чому взагалі можна ввести φ? (Підказка: див. [A].)
\item Чому φ обчислювати простіше, ніж E?
\item Чому Δφ «фізичніше», ніж φ?
\end{enumerate}

\chapter*{5. Диполь в однорідному електричному полі}\addcontentsline{toc}{chapter}{5. Диполь в однорідному електричному полі}\markboth{5. Диполь в однорідному електричному полі}{5. Диполь в однорідному електричному полі}

\srcline{\book{} т. 2, \textbf{§ 1.6} (с. 32–35); потенціал і енергія — у \textbf{§ 1.11} (с. 53–54). \ok{} Формули звірено зі сканом.}

\needspace{5\baselineskip}\section*{План сторінки}
\begin{layout}
┌ [A] Означення диполя, p (1.19) + рис. 1.9  ┐
├ [B] Поле диполя: на осі (1.20), на перпендикулярі (1.21) ┤
├ [C] Момент сил у однорідному полі (1.23)–(1.25) + рис. 1.13 ┤
└ [D] Енергія, рівновага  │ [E] Потенціал диполя (1.66)       ┘
\end{layout}

\needspace{5\baselineskip}\section*{[A] Електричний диполь (с. 32)}
\begin{itemize}
\item «сукупність двох однакових за абсолютним значенням і протилежних за знаком точкових зарядів +q і −q, розміщених на деякій відстані l» (англ. electric dipole)
\item \textbf{Плече диполя} l (англ. dipole arm). \textbf{Жорсткий} диполь (англ. rigid dipole): l = const. \textbf{Точковий} диполь (англ. point dipole): l ≪ r. Приклад — полярні молекули.
\item \textbf{Електричний дипольний момент} (англ. dipole moment): $\vec p = q\vec l$ (1.19), напрямлений \textbf{від − до +}; одиниця Кл·м.
\end{itemize}

\needspace{5\baselineskip}\section*{[B] Поле точкового диполя (с. 32–34)}
\begin{itemize}
\item На осі: $\vec E = \dfrac{1}{4\pi\varepsilon_0\varepsilon}\dfrac{2\vec p}{r^3}$ (1.20)
\item На перпендикулярі до середини плеча: $\vec E = -\dfrac{1}{4\pi\varepsilon_0\varepsilon}\dfrac{\vec p}{r^3}$ (1.21), антипаралельне p.
\item У загальному випадку (1.22) спадає як \textbf{1/r³} (для точкового заряду — 1/r²).
\end{itemize}

\needspace{5\baselineskip}\section*{[C] Дія однорідного поля (с. 34)}
\begin{itemize}
\item F₊ = −F₋, рівнодійна 0, але сили утворюють \textbf{пару сил} (англ. couple):
$M = pE\sin\alpha$ (1.24), $\vec M = [\vec p\times\vec E]$ (1.25)
\item Пара сил повертає \textbf{p} у напрямі \textbf{E}.
\item \pen{} Намалюйте рис. 1.13: диполь під кутом α до E і дві сили qE.
\end{itemize}

\needspace{5\baselineskip}\section*{[D] Енергія і рівновага (с. 34, 54)}
\begin{itemize}
\item $W = -pE\cos\alpha = -\vec p\cdot\vec E$ (1.67); у книзі на с. 54 енергію позначено U.
\item p ∥ E — \textbf{стійка рівновага} (англ. stable equilibrium; W мінімальна); p антипаралельний до E — \textbf{нестійка} (англ. unstable).
\end{itemize}

\needspace{5\baselineskip}\section*{[E] Потенціал диполя (с. 53)}
\begin{itemize}
\item $\varphi = \dfrac{1}{4\pi\varepsilon_0}\dfrac{\vec p\cdot\vec r}{r^3} = \dfrac{p\cos\theta}{4\pi\varepsilon_0 r^2}$ (1.66), спадає як 1/r².
\end{itemize}

\needspace{5\baselineskip}\section*{Порівняння з темою 2}
{\small
\begin{longtable}{>{\raggedright\arraybackslash}p{0.214\linewidth}>{\raggedright\arraybackslash}p{0.583\linewidth}>{\raggedright\arraybackslash}p{0.155\linewidth}}
\toprule
\textbf{Поле} & \textbf{Сила} & \textbf{Момент} \\
\midrule
\endhead
однорідне & 0 & pE sin α \\
неоднорідне & p ∂E/∂x, у бік сильнішого поля & теж є \\
\bottomrule
\end{longtable}
}

\needspace{5\baselineskip}\section*{Самоперевірка}
\begin{enumerate}
\item Чому в однорідному полі рівнодійна дорівнює нулю?
\item Яка рівновага стійка і чому? (Скористайтеся W.)
\item Чому поле диполя спадає швидше, ніж поле заряду?
\end{enumerate}

\chapter*{6. Потік вектора напруженості. Теорема Гаусса}\addcontentsline{toc}{chapter}{6. Потік вектора напруженості. Теорема Гаусса}\markboth{6. Потік вектора напруженості. Теорема Гаусса}{6. Потік вектора напруженості. Теорема Гаусса}

\srcline{\book{} т. 2, \textbf{§ 1.7} (с. 36–38). Застосування — у темах 10–12. \ok{} Формули звірено зі сканом.}
Правопис: у книзі \textbf{Гаусса} (з подвоєним «с»), у програмі — «Гауса».

\needspace{5\baselineskip}\section*{План сторінки}
\begin{layout}
┌ [A] Потік: означення (1.27) + рис. 1.16 (площадка, нормаль, кут α) ┐
├ [B] Ідея доведення: тілесний кут (1.29), рис. 1.17 а/б             ┤
├ [C] Теорема Гаусса: формулювання + (1.31), (1.33)                  ┤
└ [D] Фізичний зміст: джерела/стоки; Гаусс = Кулон в інтегральній формі ┘
\end{layout}

\needspace{5\baselineskip}\section*{[A] Потік вектора E (с. 36)}
\begin{itemize}
\item Однорідне поле крізь плоску площадку: $\Delta\Phi_E = E\,\Delta S\cos\alpha = E_n\Delta S$ (1.27) (англ. electric flux)
\item Загалом: $\Phi_E = \int_S E_n\,dS$
\item Знак визначає cos α; додатною вважають \textbf{зовнішню} нормаль.
\item Наочно: кількість силових ліній, що перетинають поверхню.
\end{itemize}

\needspace{5\baselineskip}\section*{[B] Ідея доведення (с. 36–38)}
\begin{itemize}
\item Для точкового заряду $d\Phi_E = \dfrac{q}{4\pi\varepsilon_0}d\omega$ (1.29), де dω — \textbf{тілесний кут} (англ. solid angle).
\item Заряд \textbf{усередині}: $\oint d\omega = 4\pi$, отже $\Phi_E = q/\varepsilon_0$ (1.30).
\item Заряд \textbf{зовні}: кожен конус перетинає поверхню парну кількість разів, отже $\Phi_E = 0$ (1.32).
\item \pen{} Намалюйте рис. 1.17 а/б: конус від q перетинає замкнену поверхню один раз і двічі.
\end{itemize}

\needspace{5\baselineskip}\section*{[C] Теорема Гаусса (с. 38)}
\begin{itemize}
\item «потік вектора напруженості електростатичного поля у вакуумі Φ\_E крізь \textbf{довільну замкнену поверхню} S зсередини назовні дорівнює \textbf{алгебричній сумі} тих точкових зарядів $q_{i}$, які охоплюються поверхнею S, поділеній на електричну сталу ε₀» (англ. Gauss's theorem)
\item $\Phi_E = \oint_S E_n\,dS = \dfrac{1}{\varepsilon_0}\sum_i q_i$ (1.31)
\item Неперервний розподіл: $\Phi_E = \dfrac{1}{\varepsilon_0}\left(\int_V\rho\,dV + \int_S\sigma\,dS\right)$ (1.33)
\end{itemize}

\needspace{5\baselineskip}\section*{[D] Фізичний зміст (с. 38)}
\begin{itemize}
\item Позитивні заряди — \textbf{джерела} поля (англ. sources), негативні — \textbf{стоки} (англ. sinks).
\item Теорема справджується \textit{саме тому}, що E ∝ 1/r². Тому «теорема Гаусса є \textbf{інтегральною формою запису закону Кулона}».
\item Чим корисна: замінює громіздкі інтеграли (1.15) коротким міркуванням про симетрію (теми 10–12) і веде до диференціальної форми (тема 7).
\end{itemize}

\needspace{5\baselineskip}\section*{Самоперевірка}
\begin{enumerate}
\item Чому потік не залежить від форми поверхні?
\item Чому заряд зовні дає нульовий потік?
\item Чому теорема не справджувалася б, якби сила спадала як 1/r³?
\end{enumerate}

\chapter*{7. Дивергенція вектора E}\addcontentsline{toc}{chapter}{7. Дивергенція вектора E}\markboth{7. Дивергенція вектора E}{7. Дивергенція вектора E}

\srcline{\book{} т. 2, \textbf{§ 1.8 «Теорема Гаусса в диференціальній формі»} (с. 41–43). \ok{} Формули звірено зі сканом.}
Перетинається з темою 13 (Остроградський–Гаусс). Пропонований поділ: \textbf{ця сторінка — div і її зміст}, \textbf{тема 13 — місток між інтегральною і диференціальною формами}.

\needspace{5\baselineskip}\section*{План сторінки}
\begin{layout}
┌ [A] Навіщо: інтегральна vs диференціальна форма ┐
├ [B] Виведення: паралелепіпед dx dy dz, рис. 1.22 ┤
├ [C] Означення div (через границю) + декартова форма ┤
└ [D] div E = ρ/ε₀ (1.42) + фізичний зміст        ┘
\end{layout}

\needspace{5\baselineskip}\section*{[A] Мотивація (с. 41)}
\begin{itemize}
\item Інтегральна форма пов'язує \textbf{різні} точки: E на поверхні й q усередині об'єму.
\item Диференціальна форма пов'язує величини \textbf{в одній точці} — це відповідає ідеї близькодії.
\end{itemize}

\needspace{5\baselineskip}\section*{[B] Виведення (с. 41–42)}
\begin{itemize}
\item Малий паралелепіпед dx·dy·dz із зарядом ρ dV. Потік крізь пару граней ⟂ Ox дорівнює $\dfrac{\partial E_x}{\partial x}dx\,dy\,dz$; так само для y і z.
\item За теоремою Гаусса: $\left(\dfrac{\partial E_x}{\partial x}+\dfrac{\partial E_y}{\partial y}+\dfrac{\partial E_z}{\partial z}\right)dV = \dfrac{\rho\,dV}{\varepsilon_0}$ (1.38)
\item \pen{} Намалюйте паралелепіпед: $E_{x}$ входить зліва, $E_{x}$ + (∂$E_{x}$/∂x)dx виходить справа.
\end{itemize}

\needspace{5\baselineskip}\section*{[C] Означення дивергенції (с. 42)}
\begin{itemize}
\item \textbf{Розходження, або дивергенція} (англ. divergence) вектора A в точці — границя відношення потоку до об'єму:
$\operatorname{div}\vec A = \lim\limits_{V\to0}\dfrac{\oint_S A_n\,dS}{V}$; не залежить від форми S.
\item Декартова форма: $\operatorname{div}\vec E = \dfrac{\partial E_x}{\partial x}+\dfrac{\partial E_y}{\partial y}+\dfrac{\partial E_z}{\partial z}$ (1.43)
\end{itemize}

\needspace{5\baselineskip}\section*{[D] Результат (с. 42–43)}
\begin{itemize}
\item $\operatorname{div}\vec E = \dfrac{\rho}{\varepsilon_0}$ (1.42) — теорема Гаусса в диференціальній формі = \textbf{диференціальна форма закону Кулона}.
\item Зміст: div характеризує \textbf{потужність джерел або стоків} (англ. strength of sources or sinks).
\begin{itemize}
\item div E > 0: тут лінії починаються (+ заряд).
\item div E < 0: тут лінії закінчуються (− заряд).
\item Де ρ = 0: div E = 0, лінії проходять крізь точку.
\end{itemize}
\end{itemize}

\needspace{5\baselineskip}\section*{Самоперевірка}
\begin{enumerate}
\item Що означає div E = 0 у точці порожнього простору?
\item Чому (1.42) «локальна», а (1.31) — ні?
\end{enumerate}

\chapter*{8. Зв'язок між напруженістю та потенціалом}\addcontentsline{toc}{chapter}{8. Зв'язок між напруженістю та потенціалом}\markboth{8. Зв'язок між напруженістю та потенціалом}{8. Зв'язок між напруженістю та потенціалом}

\srcline{\book{} т. 2, \textbf{§ 1.11} (с. 51–53). \ok{} Формули звірено зі сканом.}

\needspace{5\baselineskip}\section*{План сторінки}
\begin{layout}
┌ [A] Інтегральний зв'язок: φ₁ − φ₂ = ∫E dl (1.60)        ┐
├ [B] Диференціальний: E_x = −∂φ/∂x → E = −grad φ (1.62)   ┤
├ [C] Однорідне поле: E = U/d   │ [D] Градієнт: що це        ┤
└ [E] Рівняння Пуассона (1.63) і Лапласа (1.64)            ┘
\end{layout}

\needspace{5\baselineskip}\section*{[A] φ через E (с. 51)}
\begin{itemize}
\item $\varphi_1 - \varphi_2 = \int_1^2 \vec E\,d\vec l$ (1.60) — так знаходять φ, коли відоме E.
\begin{itemize}
\item \warn{} \textbf{Друкарська помилка в книзі} (с. 51): (1.60) надруковано з мінусом, φ₁ − φ₂ = −∫E dl. Це суперечить її ж (1.53) і (1.59): A₁₂ = q₀∫E dl = q₀(φ₁ − φ₂). Правильна форма — без мінуса: уздовж E потенціал спадає.
\end{itemize}
\end{itemize}

\needspace{5\baselineskip}\section*{[B] E через φ (с. 51–52)}
\begin{itemize}
\item На малому dx: $dA = q_0E_x\,dx = -q_0\,d\varphi$, звідси $E_x = -\dfrac{\partial\varphi}{\partial x}$, так само $E_y$, $E_z$.
\item $\vec E = -\left(\dfrac{\partial\varphi}{\partial x}\vec i + \dfrac{\partial\varphi}{\partial y}\vec j + \dfrac{\partial\varphi}{\partial z}\vec k\right) = -\operatorname{grad}\varphi$ (1.61)–(1.62)
\item \textbf{Мінус}: E напрямлена в бік \textbf{спадання} φ.
\end{itemize}

\needspace{5\baselineskip}\section*{[C] Однорідне поле (с. 52)}
\begin{itemize}
\item $E = \dfrac{\varphi_1-\varphi_2}{d} = \dfrac{U}{d}$ — тому E вимірюють у \textbf{В/м}.
\end{itemize}

\needspace{5\baselineskip}\section*{[D] Градієнт (с. 52)}
\begin{itemize}
\item \textbf{grad φ} (англ. gradient) — вектор, напрямлений у бік \textbf{найшвидшого зростання} φ; його довжина — похідна φ у цьому напрямі.
\item Тому E ⟂ еквіпотенціалям (тема 4, [F]) і напрямлена «згори вниз».
\item \pen{} Корисний образ: φ — висота пагорба, E — напрям, куди скочується кулька.
\end{itemize}

\needspace{5\baselineskip}\section*{[E] Рівняння Пуассона і Лапласа (с. 52–53)}
\begin{itemize}
\item Підставляємо E = −grad φ у div E = ρ/ε₀:
$\nabla^2\varphi = -\dfrac{\rho}{\varepsilon_0}$ (1.63) — \textbf{рівняння Пуассона} (англ. Poisson's equation)
\item Де ρ = 0: $\nabla^2\varphi = 0$ (1.64) — \textbf{рівняння Лапласа} (англ. Laplace's equation)
\item $\nabla^2 = \dfrac{\partial^2}{\partial x^2}+\dfrac{\partial^2}{\partial y^2}+\dfrac{\partial^2}{\partial z^2}$ — \textbf{лапласіан}
\item А rot grad φ ≡ 0, тож друге рівняння (rot E = 0) виконується автоматично (зв'язок із темою 9).
\end{itemize}

\needspace{5\baselineskip}\section*{Самоперевірка}
\begin{enumerate}
\item Де поле сильніше: там, де еквіпотенціалі густіші, чи де рідші?
\item Чи може E = 0 там, де φ ≠ 0? Наведіть приклад. (Усередині провідника.)
\end{enumerate}

\chapter*{9. Теорема Стокса}\addcontentsline{toc}{chapter}{9. Теорема Стокса}\markboth{9. Теорема Стокса}{9. Теорема Стокса}

\srcline{\book{} т. 2, \textbf{§ 1.10} (с. 45–48). Окремого параграфа про теорему Стокса в книзі немає: вона входить до § «Робота сил електричного поля. Потенціальний характер електростатичного поля». \ok{} Формули звірено зі сканом.}

\needspace{5\baselineskip}\section*{План сторінки}
\begin{layout}
┌ [A] Циркуляція E = 0 (1.46): електростатичне поле потенціальне ┐
├ [B] Ротор: означення (1.47) + рис. 1.25, 1.26                  ┤
├ [C] Декартова форма rot A (1.49)                                ┤
├ [D] Теорема Стокса: ∮A dl = ∫rot A dS                           ┤
└ [E] rot E = 0 (1.51); система div E = ρ/ε₀, rot E = 0 (1.52)    ┘
\end{layout}

\needspace{5\baselineskip}\section*{[A] Циркуляція (с. 45)}
\begin{itemize}
\item \textbf{Циркуляція вектора E} (англ. circulation): $\oint_L \vec E\,d\vec l$ по замкненому контуру L.
\item В електростатиці $\oint_L \vec E\,d\vec l = 0$ (1.46) — інтегральна форма \textbf{потенціальності} поля (англ. conservative field).
\item Наслідок: силові лінії електростатичного поля \textbf{незамкнені} — починаються й закінчуються на зарядах.
\end{itemize}

\needspace{5\baselineskip}\section*{[B] Ротор (с. 46)}
\begin{itemize}
\item Від лат. \textit{roto} — «обертаю». Характеризує «завихрення» векторного поля (англ. curl).
\item Проєкція на напрям n: $\operatorname{rot}_n\vec A = \lim\limits_{\Delta S\to0}\dfrac{\oint_L \vec A\,d\vec l}{\Delta S}$ (1.47); напрям обходу контуру — за правилом \textbf{правого гвинта} (англ. right-hand screw rule).
\end{itemize}

\needspace{5\baselineskip}\section*{[C] Декартова форма (с. 47)}
$\operatorname{rot}\vec A = \left(\dfrac{\partial A_z}{\partial y}-\dfrac{\partial A_y}{\partial z}\right)\vec i + \left(\dfrac{\partial A_x}{\partial z}-\dfrac{\partial A_z}{\partial x}\right)\vec j + \left(\dfrac{\partial A_y}{\partial x}-\dfrac{\partial A_x}{\partial y}\right)\vec k$ (1.49)
\begin{itemize}
\item \pen{} Як запам'ятати: визначник 3×3 з рядками (i j k), (∂x ∂y ∂z), ($A_{x}$ $A_{y}$ $A_{z}$).
\end{itemize}

\needspace{5\baselineskip}\section*{[D] Теорема Стокса (с. 47)}
\begin{itemize}
\item $\oint_L \vec A\,d\vec l = \int_S \operatorname{rot}\vec A\,d\vec S$, де S — \textbf{будь-яка} поверхня, що спирається на контур L (англ. Stokes' theorem).
\item Словами: циркуляція по контуру = потік ротора крізь поверхню, натягнуту на цей контур.
\item \pen{} Малюнок: контур L, «мильна плівка» S на ньому та нормалі n.
\end{itemize}

\needspace{5\baselineskip}\section*{[E] Застосування до E (с. 47–48)}
\begin{itemize}
\item $\int_S \operatorname{rot}\vec E\,d\vec S = \oint \vec E\,d\vec l = 0$ (1.50) для будь-якої S, отже $\operatorname{rot}\vec E = 0$ (1.51).
\item Два \textbf{основні диференціальні рівняння електростатики} (1.52):
$\operatorname{div}\vec E = \rho/\varepsilon_0,\qquad \operatorname{rot}\vec E = 0$
\item Вони замінюють закон Кулона разом із принципом суперпозиції і є окремим випадком \textbf{рівнянь Максвелла}.
\end{itemize}

\needspace{5\baselineskip}\section*{Таблиця-пара (добре підходить для низу сторінки)}
{\small
\begin{longtable}{>{\raggedright\arraybackslash}p{0.091\linewidth}>{\raggedright\arraybackslash}p{0.208\linewidth}>{\raggedright\arraybackslash}p{0.245\linewidth}>{\raggedright\arraybackslash}p{0.392\linewidth}}
\toprule
\textbf{} & \textbf{інтегральна форма} & \textbf{диференціальна форма} & \textbf{теорема-міст} \\
\midrule
\endhead
джерела & ∮E dS = q/ε₀ & div E = ρ/ε₀ & Остроградського–Гаусса (тема 13) \\
вихори & ∮E dl = 0 & rot E = 0 & Стокса (ця тема) \\
\bottomrule
\end{longtable}
}

\needspace{5\baselineskip}\section*{Самоперевірка}
\begin{enumerate}
\item Чому замкнені силові лінії суперечили б rot E = 0?
\item Чим відрізняються «потік» і «циркуляція»?
\end{enumerate}

\chapter*{10–12. Поле зарядженої площини · двох площин · сферичної поверхні}\addcontentsline{toc}{chapter}{10–12. Поле зарядженої площини · двох площин · сферичної поверхні}\markboth{10–12. Поле зарядженої площини}{10–12. Поле зарядженої площини}

\srcline{\book{} т. 2, застосування \textbf{§ 1.7} (с. 38–40); потенціали — у \textbf{§ 1.11} (с. 54–56). \ok{} Формули звірено зі сканом.}
Один метод — три форми, тож вони добре вміщуються на одному розвороті зі спільною рамкою-«рецептом».

\needspace{5\baselineskip}\section*{Рецепт (запишіть один раз, угорі)}
\begin{enumerate}
\item Із симетрії визначте \textbf{напрям} E.
\item Оберіть замкнену \textbf{гауссову поверхню} (англ. Gaussian surface), на якій E стала або потік дорівнює 0.
\item Обчисліть потік за означенням (E·площа) і прирівняйте до $q_{\text{\text{всередині}}}/\varepsilon_0$.
\item Знайдіть E. Потім φ — із $\varphi_1-\varphi_2 = \int E\,dl$.
\end{enumerate}

\par\noindent{\color{accent!30}\rule{\linewidth}{0.4pt}}\par

\needspace{5\baselineskip}\section*{10. Поле рівномірно зарядженої нескінченної площини (с. 38–39, 54)}
\begin{itemize}
\item «Нескінченна» означає: відстань до площини ≪ її розмірів.
\item Симетрія: E ⟂ площині. Поверхня — \textbf{прямий циліндр} (англ. right cylinder), що перетинає площину, з основами S.
\item Потік: крізь бічну поверхню — 0, крізь дві основи — 2ES. Охоплений заряд σS.
\item $2ES = \dfrac{\sigma S}{\varepsilon_0}$ ⟹ $E = \dfrac{\sigma}{2\varepsilon_0}$ (1.34)
\item E не залежить від відстані, тож поле з кожного боку \textbf{однорідне} (англ. uniform).
\item Потенціал: $\varphi(x) = \varphi(0) - \dfrac{\sigma}{2\varepsilon_0}x$ — спадає лінійно.
\item \pen{} Намалюйте рис. 1.18: площина, циліндр крізь неї, стрілки з обох основ.
\end{itemize}

\needspace{5\baselineskip}\section*{11. Поле двох різнойменно заряджених паралельних площин (с. 39)}
\begin{itemize}
\item +σ і −σ. За \textbf{принципом суперпозиції} між площинами поля додаються, зовні — компенсуються.
\item Між площинами: $E = \dfrac{\sigma}{\varepsilon_0}$ (1.35). Зовні: E = 0.
\item Дві \textbf{однойменно} заряджені площини: між ними E = 0.
\item Потенціал між площинами: $\varphi(x) = \varphi(0) - \dfrac{\sigma}{\varepsilon_0}x$, звідси $U = \dfrac{\sigma d}{\varepsilon_0}$ — плоский конденсатор (тема 15).
\item \pen{} Намалюйте рис. 1.19: дві пластини, три області, стрілки поля від кожної пластини та їхня сума.
\item Пов'язане: поле біля самої поверхні провідника $E = \sigma/\varepsilon_0$ (1.36), с. 40 (тема 14).
\end{itemize}

\needspace{5\baselineskip}\section*{12. Поле рівномірно зарядженої сферичної поверхні (кулі) (с. 40, 55)}
\begin{itemize}
\item Симетрія: E радіальна. Поверхня — \textbf{концентрична сфера} радіуса r.
\item Зовні (r > R): $E\cdot4\pi r^2 = q/\varepsilon_0$ ⟹ $E = \dfrac{1}{4\pi\varepsilon_0}\dfrac{q}{r^2}$ — як поле точкового заряду в центрі.
\item Усередині \textbf{сфери} (провідна куля, заряд на поверхні): E = 0, φ = const = $\dfrac{q}{4\pi\varepsilon_0R}$.
\item Потенціал зовні: $\varphi(r) = \dfrac{q}{4\pi\varepsilon_0 r}$.
\item Необов'язково: куля, заряджена \textbf{по об'єму} з густиною ρ, — усередині $E = \dfrac{\rho r}{3\varepsilon_0\varepsilon}$, зростає лінійно.
\item \pen{} Побудуйте графіки E(r) і φ(r) для сфери: E стрибає на R, далі спадає як 1/r²; φ стала всередині, далі спадає як 1/r.
\end{itemize}

\needspace{5\baselineskip}\section*{Додатковий рядок (с. 40, циліндр)}
\begin{itemize}
\item Нескінченний заряджений циліндр або нитка з лінійною густиною τ: $E = \dfrac{\tau}{2\pi\varepsilon_0 r}$ (1.37) зовні, 0 усередині циліндричної поверхні.
\begin{itemize}
\item \warn{} \textbf{Друкарська помилка в книзі} (с. 40): (1.37) надруковано як E = η/(4πε₀r), але рядок вище, E·2πrh = ηh/ε₀, дає \textbf{η/(2πε₀r)}. Сама книга використовує η/(2πε₀r) для циліндричного конденсатора (с. 62). Лінійну густину тут позначено \textbf{η}, а не τ.
\end{itemize}
\end{itemize}

\needspace{5\baselineskip}\section*{Підсумкова таблиця (добре підходить для низу сторінки)}
{\small
\begin{longtable}{>{\raggedright\arraybackslash}p{0.180\linewidth}>{\raggedright\arraybackslash}p{0.240\linewidth}>{\raggedright\arraybackslash}p{0.324\linewidth}>{\raggedright\arraybackslash}p{0.192\linewidth}}
\toprule
\textbf{Розподіл} & \textbf{Гауссова поверхня} & \textbf{E} & \textbf{залежність від r} \\
\midrule
\endhead
площина σ & циліндр ⟂ & σ/2ε₀ & const \\
дві площини ±σ & — (суперпозиція) & σ/ε₀ між, 0 зовні & const \\
нитка/циліндр τ & коаксіальний циліндр & τ/2πε₀r & 1/r \\
сфера q & концентрична сфера & q/4πε₀r² зовні, 0 всередині & 1/r² \\
\bottomrule
\end{longtable}
}

\needspace{5\baselineskip}\section*{Самоперевірка}
\begin{enumerate}
\item Чому для площини виходить 1/(2ε₀), а біля поверхні провідника — 1/ε₀?
\item Чому для площини E не залежить від відстані, а для сфери — залежить?
\end{enumerate}

\chapter*{13. Теорема Остроградського–Гаусса}\addcontentsline{toc}{chapter}{13. Теорема Остроградського–Гаусса}\markboth{13. Теорема Остроградського–Гаусса}{13. Теорема Остроградського–Гаусса}

\srcline{\book{} т. 2, \textbf{§ 1.8} (с. 42); також у § 2.1 (с. 109) для рівняння неперервності. \ok{} Формули звірено зі сканом.}
Зверніть увагу на назви. У програмі — «Теорема Остроградського-Гауса». Книга називає так \textbf{математичну формулу} (англ. divergence theorem), а фізичний закон (тема 6) у неї зветься «теорема Гаусса». Не змішуйте їх на сторінці.

\needspace{5\baselineskip}\section*{План сторінки}
\begin{layout}
┌ [A] Математична теорема: ∮A dS = ∫div A dV     ┐
├ [B] Застосування до E: (1.39) → (1.40) → (1.41) ┤
├ [C] Висновок: div E = ρ/ε₀ (1.42)              ┤
└ [D] Ще одне застосування: рівняння неперервності (2.12) ┘
\end{layout}

\needspace{5\baselineskip}\section*{[A] Формула Остроградського–Гаусса (с. 42)}
\begin{itemize}
\item Для будь-якого векторного поля A, замкненої поверхні S та обмеженого нею об'єму V:
$\oint_S A_n\,dS = \int_V \operatorname{div}\vec A\,dV$
\item Словами: потік назовні крізь межу = сума «джерел» усередині.
\item \pen{} Малюнок: об'єм, заповнений малими кубиками; потоки крізь спільні грані взаємно знищуються, лишається лише зовнішня поверхня.
\end{itemize}

\needspace{5\baselineskip}\section*{[B] Застосування до поля E (с. 42)}
\begin{enumerate}
\item Фізика (Гаусс, тема 6): $\oint_S E_n\,dS = \dfrac1{\varepsilon_0}\int_V\rho\,dV$ (1.39)
\item Математика (Остроградський–Гаусс): $\oint_S E_n\,dS = \int_V\operatorname{div}\vec E\,dV$ (1.40)
\item Отже: $\int_V\operatorname{div}\vec E\,dV = \dfrac1{\varepsilon_0}\int_V\rho\,dV$ (1.41)
\end{enumerate}

\needspace{5\baselineskip}\section*{[C] Висновок (с. 42)}
\begin{itemize}
\item Об'єм V \textbf{довільний}, тож підінтегральні вирази рівні: $\operatorname{div}\vec E = \dfrac{\rho}{\varepsilon_0}$ (1.42)
\item Той самий результат, що в темі 7, але іншим шляхом: там — через малий паралелепіпед, тут — через загальну теорему.
\end{itemize}

\needspace{5\baselineskip}\section*{[D] Додатково: рівняння неперервності (§ 2.1, с. 109)}
\begin{itemize}
\item Той самий прийом для густини струму j: $\dfrac{\partial\rho}{\partial t} = -\operatorname{div}\vec j$ (2.12) — закон збереження заряду в локальній формі. Зв'язок із темами 1 і 20.
\end{itemize}

\needspace{5\baselineskip}\section*{Самоперевірка}
\begin{enumerate}
\item На якому кроці потрібне «V довільний» і чому?
\item Хто що довів? (Остроградський — математику, Гаусс — електростатику.)
\end{enumerate}

\chapter*{14. Провідник у зовнішньому електричному полі}\addcontentsline{toc}{chapter}{14. Провідник у зовнішньому електричному полі}\markboth{14. Провідник у зовнішньому електричному полі}{14. Провідник у зовнішньому електричному полі}

\srcline{\book{} т. 2, \textbf{§ 1.12} «Розподіл зарядів на поверхні провідника» (с. 56–59), \textbf{§ 1.13} «Провідники в електричному полі» (с. 59–60). \ok{} Формули звірено зі сканом.}

\needspace{5\baselineskip}\section*{План сторінки}
\begin{layout}
┌ [A] Провідник: вільні носії; рівновага → E = 0 всередині ┐
├ [B] Наслідки: ρ = 0, заряд на поверхні, поверхня еквіпотенціальна, E ⟂ поверхні ┤
├ [C] Електростатична індукція (§1.13) + малюнок        ┤
├ [D] Порожнина, клітка Фарадея │ [E] Кривизна і σ, вістря ┤
└ низ: E біля поверхні = σ/ε₀ (1.36); Ван де Грааф (необов'язково) ┘
\end{layout}

\needspace{5\baselineskip}\section*{[A] Провідник у рівновазі (с. 56)}
\begin{itemize}
\item \textbf{Провідники} (англ. conductors) мають \textbf{вільні носії заряду} (англ. free charge carriers): електрони в металах, іони в електролітах.
\item У рівновазі носії не дрейфують, отже \textbf{усередині E = 0}.
\end{itemize}

\needspace{5\baselineskip}\section*{[B] Наслідки (с. 56–57)}
\begin{itemize}
\item З div E = ρ/ε₀ при E = 0 маємо \textbf{ρ = 0 усередині}: надлишковий заряд розміщується \textbf{лише на зовнішній поверхні} й описується σ.
\item Поверхня \textbf{еквіпотенціальна}, і в кожній точці \textbf{E ⟂ поверхні}.
\item Безпосередньо біля поверхні: $E = \sigma/\varepsilon_0$ (1.36), с. 40.
\end{itemize}

\needspace{5\baselineskip}\section*{[C] Електростатична індукція (с. 59–60)}
\begin{itemize}
\item Термін книги: «Явище перерозподілу електричних зарядів у провіднику під дією зовнішнього електричного поля … називають \textbf{електростатичною індукцією} або \textbf{електризацією через вплив}» (англ. electrostatic induction)
\item Електрони зміщуються проти E, тож ближня до поля грань заряджається −, дальня — +.
\item Індуковані заряди створюють усередині поле \textbf{E′, протилежне E}. Зміщення припиняється, коли E′ = E, і сумарне поле всередині дорівнює 0.
\item Заряджання через вплив: два дотичні провідники вносять у поле, роз'єднують, потім прибирають поле — одне тіло лишається +, інше −.
\item \pen{} Малюнок: металевий брусок в однорідному E, − на лівій грані, + на правій, E′ усередині напрямлена назустріч.
\end{itemize}

\needspace{5\baselineskip}\section*{[D] Порожнина, захист (с. 57)}
\begin{itemize}
\item Поле всередині замкненої провідної порожнини дорівнює 0 (досліди Фарадея). На цьому ґрунтується \textbf{електростатичний захист} (англ. electrostatic shielding).
\item Заряджене тіло всередині порожнини віддає \textbf{весь} свій заряд на зовнішню поверхню. Необов'язково: так працює \textbf{електростатичний генератор} (Ван де Граафа), с. 57–58.
\end{itemize}

\needspace{5\baselineskip}\section*{[E] Кривизна поверхні (с. 58–59)}
\begin{itemize}
\item Дві кулі, з'єднані дротом, мають $\varphi_1=\varphi_2$, тож $\dfrac{\sigma_1}{\sigma_2} = \dfrac{R_2}{R_1}$.
\item σ, а отже й E, більші там, де більша кривизна, тому заряд стікає з \textbf{вістря} (англ. sharp points). На цьому працюють блискавковідводи; через це ж втрачають енергію лінії електропередачі.
\end{itemize}

\needspace{5\baselineskip}\section*{Самоперевірка}
\begin{enumerate}
\item Чому E = 0 усередині — з погляду носіїв заряду?
\item Куди дівається надлишковий заряд і чому не в об'єм?
\item Чому заряд стікає з вістря?
\end{enumerate}

\chapter*{15. Конденсатори та їх ємність. Ємність плоского, циліндричного та сферичного конденсаторів}\addcontentsline{toc}{chapter}{15. Конденсатори та їх ємність. Ємність плоского, циліндричного та сферичного конденсаторів}\markboth{15. Конденсатори та їх ємність. Ємність плоского,…}{15. Конденсатори та їх ємність. Ємність плоского,…}

\srcline{\book{} т. 2, \textbf{§ 1.14} (с. 60–63). \ok{} Формули звірено зі сканом.}

\needspace{5\baselineskip}\section*{План сторінки}
\begin{layout}
┌ [A] Ємність відокремленого провідника (1.68); фарад; куля C = 4πε₀εR ┐
├ [B] Конденсатор: означення, обкладки, (1.70); вплив діелектрика ε     ┤
├ [C] Плоский │ [D] Циліндричний │ [E] Сферичний  (три малюнки в ряд)  ┤
└ низ: робоча напруга / пробій                                           ┘
\end{layout}

\needspace{5\baselineskip}\section*{[A] Електрична ємність провідника (с. 60)}
\begin{itemize}
\item Для відокремленого провідника q/φ = const: $C = \dfrac{q}{\varphi}$ (1.68), $\Delta q = C\Delta\varphi$ (1.69) (англ. capacitance)
\item Зміст: «електрична ємність провідника чисельно дорівнює зарядові, який потрібний для зміни потенціалу провідника на одиницю»
\item Одиниця — \textbf{фарад}, 1 Ф = 1 Кл/В. Це дуже велика одиниця: ємність 1 Ф мала б куля радіуса R = 9·10⁶ км. Тому вживають мкФ = 10⁻⁶ Ф і пФ = 10⁻¹² Ф.
\item Куля: $C = 4\pi\varepsilon_0\varepsilon R$.
\item C залежить від \textbf{розмірів, форми та навколишнього середовища}, але не від матеріалу чи порожнин усередині.
\end{itemize}

\needspace{5\baselineskip}\section*{[B] Конденсатор (с. 61)}
\begin{itemize}
\item «система з двох металевих електродів, розміщених на близькій відстані один від одного і розділених шаром діелектрика» (англ. capacitor); електроди називають \textbf{обкладками} (англ. plates).
\item Поле зосереджене майже повністю між обкладками, тож C майже не залежить від сусідніх тіл (на відміну від відокремленого провідника).
\item $C = \dfrac{q}{\varphi_1-\varphi_2} = \dfrac{q}{U}$ (1.70)
\item Із діелектриком: $C/C_0 = \varepsilon$.
\end{itemize}

\needspace{5\baselineskip}\section*{[C] Плоский (с. 61)}
\begin{itemize}
\item Дві паралельні пластини площею S на відстані d ≪ розмірів пластин (крайовими ефектами нехтуємо).
\item $U = Ed = \dfrac{\sigma}{\varepsilon_0}d$ (тема 11) ⟹ $C = \dfrac{\varepsilon_0\varepsilon S}{d}$
\end{itemize}

\needspace{5\baselineskip}\section*{[D] Циліндричний (с. 61–62)}
\begin{itemize}
\item Два коаксіальні циліндри r₁ < r₂ висотою h ≫ r.
\item $E = \dfrac{\tau}{2\pi\varepsilon_0\varepsilon r}$ ⟹ $U = \dfrac{\tau}{2\pi\varepsilon_0\varepsilon}\ln\dfrac{r_2}{r_1}$ ⟹ $C = \dfrac{2\pi\varepsilon_0\varepsilon h}{\ln(r_2/r_1)}$
\item Малий зазор d = r₂ − r₁ ≪ r₁ дає $C\approx\dfrac{2\pi\varepsilon_0\varepsilon r_1h}{d}$ — плоский конденсатор. \pen{} Варто записати, чому.
\end{itemize}

\needspace{5\baselineskip}\section*{[E] Сферичний (с. 62)}
\begin{itemize}
\item Дві концентричні сфери r₁ < r₂.
\item $U = \dfrac{q}{4\pi\varepsilon_0\varepsilon}\left(\dfrac1{r_1}-\dfrac1{r_2}\right)$ ⟹ $C = \dfrac{4\pi\varepsilon_0\varepsilon r_1r_2}{r_2-r_1}$
\item При r₂ → ∞ виходить відокремлена куля $4\pi\varepsilon_0\varepsilon r_1$ — гарна перевірка.
\end{itemize}

\needspace{5\baselineskip}\section*{Та сама схема тричі}
Знайти E (Гаусс) → проінтегрувати й отримати U → C = q/U. \pen{} Запишіть цей рядок над трьома малюнками.

\needspace{5\baselineskip}\section*{Низ (с. 63)}
\begin{itemize}
\item \textbf{Робоча напруга} (англ. rated voltage) $U_{p}$: вище неї настає \textbf{пробій} (англ. breakdown).
\item Необов'язково: двопровідна лінія, $C = \pi\varepsilon_0\varepsilon l/\ln(d/r)$ (приклад 1.3).
\end{itemize}

\needspace{5\baselineskip}\section*{Самоперевірка}
\begin{enumerate}
\item Чому конденсатори кращі за відокремлені провідники?
\item Як три формули C показують, що «C залежить лише від геометрії та ε»?
\end{enumerate}

\chapter*{16. З'єднання конденсаторів}\addcontentsline{toc}{chapter}{16. З'єднання конденсаторів}\markboth{16. З'єднання конденсаторів}{16. З'єднання конденсаторів}

\srcline{\book{} т. 2, \textbf{§ 1.14} (с. 63–64). \ok{} Формули звірено зі сканом. Коротка тема: пів сторінки.}

\needspace{5\baselineskip}\section*{План сторінки (дві колонки)}
\begin{layout}
┌ Паралельне (рис. 1.31)        │ Послідовне (рис. 1.32)          ┐
│ U однакова                    │ q однаковий (електростатична індукція) │
│ q = q₁ + q₂ + …               │ U = U₁ + U₂ + …                  │
│ C = C₁ + C₂ + … (1.71)        │ 1/C = 1/C₁ + 1/C₂ + … (1.72)     │
│ n однакових: C = nC₀          │ n однакових: C = C₀/n            │
└ низ: навіщо (батареї); імпульс nU (рис. 1.33)                    ┘
\end{layout}

\needspace{5\baselineskip}\section*{Паралельне з'єднання (с. 63)}
\begin{itemize}
\item На всіх конденсаторах однакова напруга U, заряди додаються: $q = (C_1+C_2+\ldots+C_n)U$ (англ. parallel connection).
\item $C = C_1 + C_2 + \ldots + C_n$ (1.71)
\end{itemize}

\needspace{5\baselineskip}\section*{Послідовне з'єднання (с. 63–64)}
\begin{itemize}
\item Завдяки \textbf{електростатичній індукції} заряд q на всіх конденсаторах однаковий, напруги додаються (англ. series connection).
\item $\dfrac1C = \dfrac1{C_1}+\dfrac1{C_2}+\ldots+\dfrac1{C_n}$ (1.72)
\item C завжди \textbf{менша за найменшу} Cᵢ.
\end{itemize}

\needspace{5\baselineskip}\section*{Навіщо з'єднувати (с. 63)}
\begin{itemize}
\item Щоб отримати потрібну ємність при потрібній напрузі — \textbf{батарея конденсаторів} (англ. capacitor bank). Послідовне з'єднання ділить напругу, тож батарея витримує більшу.
\item Прийом (с. 64, рис. 1.33): заряджають n конденсаторів паралельно до U, потім перемикають послідовно й отримують nU. Так одержують імпульси високої напруги.
\end{itemize}

\needspace{5\baselineskip}\section*{\pen{} Ваш рядок}
\begin{itemize}
\item Порівняйте з резисторами (тема 22): ємності \textbf{додаються при паралельному}, а опори — \textbf{при послідовному} з'єднанні. Запишіть чому: C ∝ q/U, а R ∝ U/I.
\end{itemize}

\needspace{5\baselineskip}\section*{Самоперевірка}
\begin{enumerate}
\item Чому заряд на послідовно з'єднаних конденсаторах однаковий?
\item Два конденсатори, 2 мкФ і 3 мкФ: знайдіть C для обох з'єднань. (5 і 1,2 мкФ.)
\end{enumerate}

\chapter*{17. Енергія зарядженого провідника}\addcontentsline{toc}{chapter}{17. Енергія зарядженого провідника}\markboth{17. Енергія зарядженого провідника}{17. Енергія зарядженого провідника}

\srcline{\book{} т. 2, \textbf{§ 1.25} «Енергія взаємодії електричних зарядів» (с. 101–103), початок \textbf{§ 1.26} (с. 103–104). \ok{} Формули звірено зі сканом.}

\needspace{5\baselineskip}\section*{План сторінки}
\begin{layout}
┌ [A] Енергія двох зарядів (1.128)–(1.129) ┐
├ [B] Система n зарядів (1.130): чому ½    ┤
├ [C] Заряджання провідника: dA = φ dq → (1.132) ┤
└ [D] Те саме для конденсатора (1.133)       ┘
\end{layout}

\needspace{5\baselineskip}\section*{[A] Два точкові заряди (с. 101–102)}
\begin{itemize}
\item Зближуємо заряди з нескінченності: $W = \dfrac{1}{4\pi\varepsilon_0}\dfrac{q_1q_2}{r_{12}}$ (1.128) (англ. interaction energy)
\item У симетричній формі: $W = \tfrac12(q_1\varphi_1 + q_2\varphi_2)$ (1.129), де φ₁ — потенціал, який створює q₂ у точці q₁.
\end{itemize}

\needspace{5\baselineskip}\section*{[B] Система n зарядів (с. 102)}
\begin{itemize}
\item $W = \dfrac12\sum_{i=1}^n q_i\varphi_i$ (1.130), де φᵢ — потенціал у точці qᵢ від \textbf{усіх інших} зарядів.
\item \textbf{Чому ½}: кожна пара входить у суму двічі.
\item Неперервний розподіл: $W = \tfrac12\left(\int\rho\varphi\,dV + \int\sigma\varphi\,dS\right)$.
\end{itemize}

\needspace{5\baselineskip}\section*{[C] Енергія зарядженого провідника (с. 103–104)}
\begin{itemize}
\item Заряджаємо провідник порціями dq з нескінченності. Кожна потребує роботи $dA = \varphi\,dq = \dfrac{q}{C}dq$.
\item Від 0 до q: $A = \dfrac{q^2}{2C}$ — вона й стає запасеною енергією.
\item $W = \dfrac{q\varphi}{2} = \dfrac{C\varphi^2}{2} = \dfrac{q^2}{2C}$ (1.132)
\item \pen{} Намалюйте графік: φ від q — пряма через 0, а \textbf{W — площа під нею}, трикутник. Звідси й ½.
\end{itemize}

\needspace{5\baselineskip}\section*{[D] Конденсатор (с. 104)}
\begin{itemize}
\item Та сама формула, де φ замінено на U: $W = \dfrac{CU^2}{2} = \dfrac{qU}{2} = \dfrac{q^2}{2C}$ (1.133)
\item Наочно: заряд переносять з однієї обкладки на іншу.
\end{itemize}

\needspace{5\baselineskip}\section*{Самоперевірка}
\begin{enumerate}
\item Звідки береться ½? Дайте два пояснення: подвійний рахунок і трикутник.
\item Як зміниться W, якщо q подвоїти? (У 4 рази.)
\end{enumerate}

\chapter*{18. Енергія електростатичного поля}\addcontentsline{toc}{chapter}{18. Енергія електростатичного поля}\markboth{18. Енергія електростатичного поля}{18. Енергія електростатичного поля}

\srcline{\book{} т. 2, \textbf{§ 1.26} (с. 103–105). \ok{} Формули звірено зі сканом.}

\needspace{5\baselineskip}\section*{План сторінки}
\begin{layout}
┌ [A] Від енергії конденсатора до енергії поля (1.133)→(1.134) ┐
├ [B] Неоднорідне поле (1.135)                                 ┤
├ [C] Об'ємна густина енергії w (1.136)                         ┤
└ [D] Де «живе» енергія: заряди чи поле? (близькодія vs далекодія) ┘
\end{layout}

\needspace{5\baselineskip}\section*{[A] Перетворення для плоского конденсатора (с. 104)}
\begin{itemize}
\item Беремо $W = \tfrac12 C(\Delta\varphi)^2$ і підставляємо $\Delta\varphi = Ed$, $C = \dfrac{\varepsilon_0\varepsilon S}{d}$:
\item $W = \dfrac{\varepsilon_0\varepsilon E^2}{2}Sd = \dfrac{\varepsilon_0\varepsilon E^2}{2}V$ (1.134), де V = Sd — об'єм між пластинами.
\item Тепер енергію виражено через \textbf{поле E} і \textbf{об'єм}, без зарядів.
\end{itemize}

\needspace{5\baselineskip}\section*{[B] Загальний випадок (с. 104)}
\begin{itemize}
\item Неоднорідне поле: $W = \dfrac12\int_V \varepsilon_0\varepsilon E^2\,dV$ (1.135)
\end{itemize}

\needspace{5\baselineskip}\section*{[C] Густина енергії (с. 105)}
\begin{itemize}
\item \textbf{Об'ємна густина енергії електричного поля} (англ. energy density):
$w = \dfrac{W}{V} = \dfrac{\varepsilon_0\varepsilon E^2}{2} = \dfrac{ED}{2}$ (1.136)
\item Тут $D = \varepsilon_0\varepsilon E$ — \textbf{вектор електричної індукції} (англ. electric displacement), означений у параграфах про діелектрики (§ 1.16 і далі). На сторінці досить одного рядка.
\item Одиниця: Дж/м³.
\end{itemize}

\needspace{5\baselineskip}\section*{[D] Фізичне тлумачення (с. 105) — варто виділити рамкою}
\begin{itemize}
\item \textbf{Близькодія}: енергію несе \textbf{поле}; вона розподілена в просторі з густиною w.
\item \textbf{Далекодія}: енергія зосереджена на заряджених провідниках.
\item \textbf{В електростатиці жоден дослід не розрізнить} ці погляди, бо поле невіддільне від своїх зарядів.
\item Розв'язують питання \textbf{змінні поля}: електромагнітні хвилі переносять енергію без зарядів, отже енергія — у полі.
\end{itemize}

\needspace{5\baselineskip}\section*{Самоперевірка}
\begin{enumerate}
\item Теми 17 і 18 дають однакове W. У чому різниця між ними?
\item Чому електромагнітна хвиля доводить, що енергія зосереджена в полі?
\end{enumerate}

\chapter*{19. Електричний струм. Постійний електричний струм. Закони Ома}\addcontentsline{toc}{chapter}{19. Електричний струм. Постійний електричний струм. Закони Ома}\markboth{19. Електричний струм. Постійний електричний струм.…}{19. Електричний струм. Постійний електричний струм.…}

\srcline{\book{} т. 2, \textbf{§ 2.1} (с. 105–110), \textbf{§ 2.2} (с. 110–115), \textbf{§ 2.3} (с. 118–120). \ok{} Формули звірено зі сканом.}
Перетинається з темами 20–21. Пропонований поділ: \textbf{ця сторінка — що таке струм і всі три форми закону Ома}; тема 20 — величини I, j, ЕРС; тема 21 — R, ρ, потужність.

\needspace{5\baselineskip}\section*{План сторінки}
\begin{layout}
┌ [A] Струм: означення, носії, напрям, види ┐
├ [B] Постійний струм: умови (2.3), (2.13)–(2.14) ┤
├ [C] Закон Ома для однорідної ділянки (2.16)–(2.17) ┤
├ [D] Диференціальна форма j = σE (2.23)              ┤
└ [E] Неоднорідна ділянка (2.33) і повне коло (2.34) ┘
\end{layout}

\needspace{5\baselineskip}\section*{[A] Електричний струм (с. 105–106)}
\begin{itemize}
\item «\textbf{упорядкований рух} електричних зарядів» (англ. electric current)
\item \textbf{Носії заряду} (англ. charge carriers): електрони в металах і напівпровідниках; іони в електролітах; електрони та іони в газах.
\item \textbf{Напрям струму} (англ. current direction) — напрям руху \textbf{позитивних} зарядів; у металах він протилежний рухові електронів.
\item Види струму: \textbf{струм провідності} (англ. conduction current, усередині провідника), \textbf{конвекційний} (англ. convection current, рухається саме заряджене тіло), \textbf{струм поляризації} (англ. polarization current, у діелектриках).
\end{itemize}

\needspace{5\baselineskip}\section*{[B] Постійний струм (с. 107, 110)}
\begin{itemize}
\item \textbf{Стаціонарний, або постійний} (англ. steady current): за однакові проміжки часу проходять однакові заряди, напрям не змінюється. $I = \Delta q/\Delta t$ (2.3)
\item Неперервність постійного струму: $\oint \vec j\,d\vec S = 0$ (2.13), $\operatorname{div}\vec j = 0$ (2.14). Лінії струму \textbf{замкнені}, а I однакова вздовж нерозгалуженого провідника.
\end{itemize}

\needspace{5\baselineskip}\section*{[C] Закон Ома (однорідна ділянка) (с. 111–112)}
\begin{itemize}
\item «сила струму в однорідному провіднику прямо пропорційна напрузі на кінцях цього провідника» (англ. Ohm's law)
\item $I = kU$ (2.16), де k — \textbf{провідність} (англ. conductance, одиниця См) → $I = \dfrac{U}{R}$ (2.17)
\item \textbf{Спад напруги} (англ. voltage drop) $U = \varphi_1-\varphi_2$ уздовж провідника зі струмом.
\end{itemize}

\needspace{5\baselineskip}\section*{[D] Диференціальна форма (с. 114)}
\begin{itemize}
\item Малий циліндр dl·dS: $j\,dS = \dfrac{E\,dl}{\rho\,dl/dS}$ ⟹ $\vec j = \sigma\vec E$ (2.23)
\item «густина електричного струму в будь-якій точці фізично однорідного провідника пропорційна напруженості електричного поля»
\item σ залежить \textbf{лише від матеріалу}, а не від форми чи розмірів.
\item Закон Ома \textbf{не фундаментальний}: порушується в сильних полях, на контактах, у діодах (нелінійні елементи), с. 115.
\end{itemize}

\needspace{5\baselineskip}\section*{[E] Неоднорідна ділянка і повне коло (с. 119–120)}
\begin{itemize}
\item Ділянка з джерелом: $IR_{12} = (\varphi_1-\varphi_2) + \mathcal E_{12}$ (2.33). Правило знаків: \textbf{ЕРС > 0}, якщо, рухаючись за струмом, проходимо джерело від − до +.
\item Замкнене коло (φ₁ = φ₂): $I = \dfrac{\mathcal E}{R+r}$ (2.34), де r — \textbf{внутрішній опір} (англ. internal resistance).
\item $U = IR = \mathcal E - Ir$ (2.35). Розімкнене коло: U = ℰ. Коротке замикання: $I_{max} = \mathcal E/r$ (2.37).
\end{itemize}

\needspace{5\baselineskip}\section*{Три форми закону Ома (добре підходить для низу сторінки)}
{\small
\begin{longtable}{>{\raggedright\arraybackslash}p{0.394\linewidth}>{\raggedright\arraybackslash}p{0.574\linewidth}}
\toprule
\textbf{ділянка} & \textbf{закон} \\
\midrule
\endhead
однорідна & I = U/R \\
неоднорідна & IR = φ₁ − φ₂ + ℰ \\
повне коло & I = ℰ/(R + r) \\
точка & j = σE \\
\bottomrule
\end{longtable}
}

\needspace{5\baselineskip}\section*{Самоперевірка}
\begin{enumerate}
\item Чому постійний струм неможливий лише з кулонівськими силами? (Див. тему 20.)
\item Яка форма закону Ома локальна і чому це важливо?
\end{enumerate}

\chapter*{20. Сила струму, густина струму. Електрорушійна сила}\addcontentsline{toc}{chapter}{20. Сила струму, густина струму. Електрорушійна сила}\markboth{20. Сила струму, густина струму. Електрорушійна сила}{20. Сила струму, густина струму. Електрорушійна сила}

\srcline{\book{} т. 2, \textbf{§ 2.1} (с. 106–110), \textbf{§ 2.3} (с. 115–118). \ok{} Формули звірено зі сканом.}

\needspace{5\baselineskip}\section*{План сторінки}
\begin{layout}
┌ [A] Сила струму I (2.1); ампер, кулон              ┐
├ [B] Густина струму j (2.4)–(2.10) + рис. 2.1 (циліндр v·dt) ┤
├ [C] Рівняння неперервності (2.11)–(2.12)          ┤
├ [D] Сторонні сили: навіщо + рис. 2.4 (кругообіг зарядів) ┤
└ [E] ЕРС: означення (2.28), циркуляція (2.29)–(2.30) ┘
\end{layout}

\needspace{5\baselineskip}\section*{[A] Сила струму (с. 106)}
\begin{itemize}
\item «скалярна величина, яка чисельно дорівнює зарядові, що переноситься крізь поперечний переріз провідника за одиницю часу» (англ. current)
\item $I = \dfrac{dq}{dt}$ (2.1). Одиниця — \textbf{ампер}, основна одиниця СІ. 1 Кл = 1 А·1 с.
\end{itemize}

\needspace{5\baselineskip}\section*{[B] Густина струму (с. 107–109)}
\begin{itemize}
\item $j = \dfrac{dI}{dS_\perp}$ (2.4) — \textbf{вектор}, напрямлений за дрейфом позитивних носіїв (англ. current density).
\item Через носії: $\vec j = n_0 e\vec v_d$ (2.6), де n₀ — концентрація, $v_{d}$ — \textbf{дрейфова швидкість} (англ. drift velocity).
\item Для електронів: $\vec j = -n_0|e|\vec v_d$ (2.7).
\item $I = \int_S j_n\,dS$ (2.9). Рівномірний струм: $I = n_0ev_dS = jS$ (2.10). Одиниця А/м².
\item Зауважте: середнє значення вектора швидкості хаотичного руху дорівнює 0, тож струм переносить лише дрейф.
\item \pen{} Намалюйте рис. 2.1а: циліндр з основою dS і довжиною $v_{d}$·dt; усе, що в ньому, проходить крізь dS за час dt.
\end{itemize}

\needspace{5\baselineskip}\section*{[C] Рівняння неперервності (с. 109–110)}
\begin{itemize}
\item $\oint_S j_n\,dS = -\dfrac{\partial q}{\partial t}$ (2.11) ⟹ $\dfrac{\partial\rho}{\partial t} = -\operatorname{div}\vec j$ (2.12) за формулою Остроградського–Гаусса (тема 13) (англ. continuity equation).
\item Це \textbf{закон збереження заряду} в макроскопічній електродинаміці (тема 1).
\end{itemize}

\needspace{5\baselineskip}\section*{[D] Сторонні сили (с. 115–116)}
\begin{itemize}
\item Самі лише кулонівські сили вирівняли б потенціали (E → 0) і струм припинився б.
\item Усередині джерела заряди треба переносити \textbf{проти} E, а для цього потрібні \textbf{сторонні сили} (англ. extraneous / non-electrostatic forces): «Будь-які сили неелектростатичної природи, які діють на заряджені частинки». Вони бувають хімічні, механічні, магнітні, світлові тощо.
\item \textbf{Джерело струму} (англ. current source) — пристрій, у якому діють сторонні сили (гальванічний елемент, генератор, термопара, сонячна батарея).
\item Їхнє поле $E_{\text{ст}}$ дає $\vec j = \sigma(\vec E + \vec E_{\text{ст}})$ (2.26).
\item \pen{} Намалюйте рис. 2.4: контур, де в провіднику заряди йдуть φ₁→φ₂, а джерело переносить їх φ₂→φ₁.
\end{itemize}

\needspace{5\baselineskip}\section*{[E] Електрорушійна сила (с. 118)}
\begin{itemize}
\item «Скалярну величину ℰ, яка дорівнює \textbf{роботі сторонніх сил, віднесеній до одиниці позитивного заряду}, називають електрорушійною силою (ЕРС)» (англ. electromotive force, emf)
\item $\mathcal E = \dfrac{A_{\text{ст}}}{q}$ (2.28); у замкненому колі $\mathcal E = \oint \vec E_{\text{ст}}d\vec l$ (2.29) — \textbf{циркуляція поля сторонніх сил}. На ділянці 1–2: $\mathcal E_{12} = \int_1^2\vec E_{\text{ст}}d\vec l$ (2.30).
\item Попри назву, це \textbf{не сила}, а енергетична характеристика джерела; вимірюється у вольтах.
\item Порівняйте з темою 9: циркуляція E дорівнює 0, а циркуляція $E_{\text{ст}}$ дорівнює ℰ ≠ 0. Саме ця ненульова циркуляція й підтримує струм.
\end{itemize}

\needspace{5\baselineskip}\section*{Самоперевірка}
\begin{enumerate}
\item Чому дрейфова швидкість мізерна, а лампа спалахує миттєво?
\item Чому назва «електрорушійна \textit{сила}» оманлива?
\end{enumerate}

\chapter*{21. Поняття опору. Потужність}\addcontentsline{toc}{chapter}{21. Поняття опору. Потужність}\markboth{21. Поняття опору. Потужність}{21. Поняття опору. Потужність}

\srcline{\book{} т. 2, \textbf{§ 2.2} (с. 111–114), \textbf{§ 2.4} (с. 121–123). \ok{} Формули звірено зі сканом.}

\needspace{5\baselineskip}\section*{План сторінки}
\begin{layout}
┌ [A] Опір R: означення, Ом, звідки береться        ┐
├ [B] R = ρl/S (2.18); питомий опір; σ = 1/ρ (2.21) ┤
├ [C] Залежність від T (2.19)–(2.20) + графік ρ(T)  ┤
├ [D] Робота струму (2.38)–(2.41)                   ┤
└ [E] Потужність (2.43)–(2.46); ККД; макс. при R = r + графік P(R) ┘
\end{layout}

\needspace{5\baselineskip}\section*{[A] Електричний опір (с. 111–112)}
\begin{itemize}
\item $R = 1/k$ — величина, обернена до провідності (англ. resistance). 1 Ом: при 1 В тече 1 А.
\item Звідки береться (метали): електрони \textbf{розсіюються} на дефектах ґратки та її \textbf{теплових коливаннях}.
\item Не залежить від I (при помірних струмах); залежить від геометрії, матеріалу та стану (T, p).
\end{itemize}

\needspace{5\baselineskip}\section*{[B] Опір провідника (с. 112–113)}
\begin{itemize}
\item $R = \rho\dfrac{l}{S}$ (2.18), де ρ — \textbf{питомий опір} (англ. resistivity), Ом·м.
\item $\sigma = 1/\rho$ (2.21) — \textbf{питома електропровідність} (англ. conductivity), См/м.
\item Найменший ρ: срібло, мідь, золото, алюміній (з них роблять дроти). Великий ρ: ніхром (нагрівачі), константан і нікелін (еталонні резистори).
\item Поділ за σ: провідники σ > 10⁶, діелектрики σ < 10⁻¹⁰, напівпровідники — проміжні значення (См/м); звірено зі сканом.
\end{itemize}

\needspace{5\baselineskip}\section*{[C] Температурна залежність (с. 112–113)}
\begin{itemize}
\item \textbf{Температурний коефіцієнт опору} (англ. temperature coefficient of resistance) $\alpha = \dfrac1\rho\dfrac{d\rho}{dT}$ (2.19). Метали: α > 0. Електроліти: α < 0.
\item Для металів поблизу 0 °C $\alpha\approx 1/273$ К⁻¹ і $\rho = \rho_0(1+\alpha t)$, або $\rho = \rho_0\alpha T$ (2.20).
\item Застосування: термометри опору, термостати, болометри.
\end{itemize}

\needspace{5\baselineskip}\section*{[D] Робота струму (с. 121)}
\begin{itemize}
\item $A = IUt$ (2.38) $= I^2Rt = \dfrac{U^2}{R}t$ (2.39)
\item Повна робота джерела по всьому колу: $A_{\text{заг}} = \mathcal E It$ (2.41).
\end{itemize}

\needspace{5\baselineskip}\section*{[E] Потужність (с. 121–123)}
\begin{itemize}
\item \textbf{Потужність} (англ. power) — робота за одиницю часу.
\item \textbf{Корисна} (англ. useful power; у навантаженні R): $P = IU = I^2R = \dfrac{U^2}{R}$ (2.43) $= \dfrac{\mathcal E^2R}{(R+r)^2}$ (2.44)
\item \textbf{Загальна} (англ. total power): $P_{\text{заг}} = \mathcal EI = \dfrac{\mathcal E^2}{R+r}$ (2.45)
\item \textbf{ККД} (англ. efficiency): $\eta = \dfrac{P}{P_{\text{заг}}} = \dfrac{R}{R+r}$ (2.46)
\item Максимальна корисна потужність при \textbf{R = r}: $P_{max} = \dfrac{\mathcal E^2}{4r}$, але тоді η = 50 \%.
\item \pen{} Графік P(R): нуль при R = 0, максимум при R = r, прямує до 0 при R → ∞. Поверх — η(R), що зростає до 1.
\end{itemize}

\needspace{5\baselineskip}\section*{Самоперевірка}
\begin{enumerate}
\item Чому не можна одночасно мати максимальну потужність і максимальний ККД?
\item Чому опір металів зростає при нагріванні?
\end{enumerate}

\chapter*{22. Перший та другий закон Кірхгофа}\addcontentsline{toc}{chapter}{22. Перший та другий закон Кірхгофа}\markboth{22. Перший та другий закон Кірхгофа}{22. Перший та другий закон Кірхгофа}

\srcline{\book{} т. 2, \textbf{§ 2.6} «Електричне коло. Правила Кірхгофа» (с. 126–131). \ok{} Формули звірено зі сканом.}
У книзі — \textbf{правила} Кірхгофа, у програмі — «закон». Вживаються обидва.

\needspace{5\baselineskip}\section*{План сторінки}
\begin{layout}
┌ [A] Прості з'єднання: послідовне / паралельне (табличка)    ┐
├ [B] Терміни: розгалужене коло, вузол, ділянка, контур        ┤
├ [C] I правило (2.54) + рис. 2.12 │ [D] II правило (2.55) + рис. 2.13 ┤
└ [E] Як складати рівняння: знаки, незалежні контури           ┘
\end{layout}

\needspace{5\baselineskip}\section*{[A] Послідовне і паралельне з'єднання опорів (с. 127–128)}
{\small
\begin{longtable}{>{\raggedright\arraybackslash}p{0.168\linewidth}>{\raggedright\arraybackslash}p{0.392\linewidth}>{\raggedright\arraybackslash}p{0.392\linewidth}}
\toprule
\textbf{} & \textbf{послідовне (рис. 2.7)} & \textbf{паралельне (рис. 2.9)} \\
\midrule
\endhead
струм & I = I₁ = I₂ = … & I = I₁ + I₂ + … \\
напруга & U = U₁ + U₂ + … & U = U₁ = U₂ = … \\
опір & R = R₁ + R₂ + … & 1/R = 1/R₁ + 1/R₂ + … \\
пропорція & U₁/U₂ = R₁/R₂ & I₁/I₂ = R₂/R₁ \\
\bottomrule
\end{longtable}
}
\begin{itemize}
\item \textbf{Дільник напруги (потенціометр)} (англ. voltage divider), с. 127.
\item Необов'язково, батареї джерел (с. 128–129): послідовно ℰ = Σℰₖ, r = Σrₖ; n однакових паралельно дають r = r₁/n.
\end{itemize}

\needspace{5\baselineskip}\section*{[B] Терміни (с. 129)}
\begin{itemize}
\item \textbf{Розгалужене коло} (англ. branched circuit): мережа з кількох контурів зі спільними ділянками.
\item \textbf{Вузол} (англ. node): «Точку розгалуженого кола, в якій сходяться \textbf{три або більше} провідників із струмами»
\item \textbf{Ділянка} (англ. branch): частина контуру між двома вузлами, по якій тече один струм.
\end{itemize}

\needspace{5\baselineskip}\section*{[C] Перше правило (с. 129)}
\begin{itemize}
\item Виражає \textbf{закон збереження заряду} у вузлі: «алгебрична сума сил струмів, які сходяться у вузлі, дорівнює нулеві» (англ. Kirchhoff's current law)
\item $\sum_{k=1}^n I_k = 0$ (2.54). Струми, що входять у вузол, — зі знаком +, що виходять, — зі знаком −.
\item Чому справджується: інакше у вузлі накопичувався б заряд, змінювалися б потенціали, і струм не був би постійним.
\end{itemize}

\needspace{5\baselineskip}\section*{[D] Друге правило (с. 130)}
\begin{itemize}
\item Узагальнення \textbf{закону Ома} на контур: «у контурі алгебрична сума спадів напруг (добутків сил струмів на опори відповідних ділянок) дорівнює алгебричній сумі електрорушійних сил, які діють у цьому контурі» (англ. Kirchhoff's voltage law)
\item $\sum_{k=1}^n I_kR_k = \sum_{k=1}^m \mathcal E_k$ (2.55)
\item Ідея доведення: записати (2.33) для кожної ділянки контуру й додати; \textbf{потенціали скорочуються}.
\end{itemize}

\needspace{5\baselineskip}\section*{[E] Як користуватися (с. 129–131)}
\begin{enumerate}
\item Напрям кожного струму обирають довільно; від'ємна відповідь означає, що струм тече навпаки.
\item Обирають напрям обходу кожного контуру. IR береться з +, якщо струм збігається з напрямом обходу. ℰ береться з +, якщо при обході джерело проходять від − до +.
\item Беруть лише \textbf{незалежні} контури — ті, які не можна утворити з інших.
\item \pen{} Розв'яжіть самі один приклад за рис. 2.13: три ділянки, два вузли.
\end{enumerate}

\needspace{5\baselineskip}\section*{Самоперевірка}
\begin{enumerate}
\item Який закон збереження виражає кожне правило? (Заряду; енергії — через закон Ома.)
\item Скільки незалежних рівнянь для вузлів при N вузлах? (N − 1.)
\end{enumerate}

\chapter*{23. Закон Джоуля–Ленца}\addcontentsline{toc}{chapter}{23. Закон Джоуля–Ленца}\markboth{23. Закон Джоуля–Ленца}{23. Закон Джоуля–Ленца}

\srcline{\book{} т. 2, \textbf{§ 2.5} «Закон Джоуля — Ленца. Закон збереження енергії електричного поля» (с. 123–126). \ok{} Формули звірено зі сканом.}

\needspace{5\baselineskip}\section*{План сторінки}
\begin{layout}
┌ [A] Формулювання + Q = I²Rt (2.47)            ┐
├ [B] Механізм: зіткнення електронів з ґраткою   ┤
├ [C] Диференціальна форма: w = ρj² = σE² = jE (2.48)–(2.50) ┤
└ [D] Закон збереження енергії в колі (2.51), два наслідки ┘
\end{layout}

\needspace{5\baselineskip}\section*{[A] Закон (с. 123)}
\begin{itemize}
\item Джоуль (1841) і Ленц (1842), встановлено експериментально (англ. Joule–Lenz law).
\item «кількість теплоти, що виділяється в провіднику при проходженні по ньому постійного електричного струму, прямо пропорційна добуткові квадрата сили струму, опору провідника і часу проходження струму»
\item $Q = I^2Rt$ (2.47). Також дорівнює $IUt = \dfrac{U^2}{R}t$ (для нерухомого провідника без хімічних перетворень).
\item Змінний струм: $Q = \int_0^t I^2(t)R\,dt$.
\end{itemize}

\needspace{5\baselineskip}\section*{[B] Механізм (с. 123–124)}
\begin{itemize}
\item Електрони набирають енергію від поля \textbf{між зіткненнями} і віддають її іонам ґратки \textbf{під час зіткнень}. Упорядкований рух перетворюється на хаотичний — це й є теплота.
\item \pen{} Малюнок: «пилкоподібний» графік енергії електрона в часі — зростає й різко падає при кожному зіткненні.
\end{itemize}

\needspace{5\baselineskip}\section*{[C] Диференціальна форма (с. 124)}
\begin{itemize}
\item \textbf{Питома потужність струму} (англ. power density) — теплота на одиницю об'єму за одиницю часу: $w = \dfrac{dQ}{dV\,dt}$
\item $w = \rho j^2$ (2.48), $w = \sigma E^2$ (2.49), $w = \vec j\cdot\vec E$ (2.50)
\item Інтегрування за об'ємом і часом повертає (2.47). Знову локальна vs інтегральна форма, як для закону Ома.
\end{itemize}

\needspace{5\baselineskip}\section*{[D] Закон збереження енергії (с. 125–126)}
\begin{itemize}
\item $dW = \delta Q + dU + \delta A$ (2.51): енергія джерела йде на теплоту, зміну енергії поля та механічну роботу.
\item Наслідки: (1) якщо нічого не рухається, уся робота джерела перетворюється на теплоту; (2) якщо заряди провідників сталі, механічна робота виконується за рахунок зменшення енергії поля.
\item Подробиці про пондеромоторні сили — \textbf{необов'язкові}, їх можна пропустити.
\end{itemize}

\needspace{5\baselineskip}\section*{Пара із законом Ома (добре підходить для низу сторінки)}
{\small
\begin{longtable}{>{\raggedright\arraybackslash}p{0.291\linewidth}>{\raggedright\arraybackslash}p{0.291\linewidth}>{\raggedright\arraybackslash}p{0.370\linewidth}}
\toprule
\textbf{} & \textbf{інтегральна} & \textbf{диференціальна} \\
\midrule
\endhead
Ом & I = U/R & j = σE \\
Джоуль–Ленц & Q = I²Rt & w = σE² = jE \\
\bottomrule
\end{longtable}
}

\needspace{5\baselineskip}\section*{Самоперевірка}
\begin{enumerate}
\item Чому в послідовному колі тонший провідник нагрівається сильніше? (Той самий I, більший R.)
\item Чому нагрівачі роблять з ніхрому, а не з міді?
\end{enumerate}

\part{Молекулярна фізика і термодинаміка}

\chapter*{Програма МКР №2: 30 питань}\addcontentsline{toc}{chapter}{Програма МКР №2: 30 питань}\markboth{Програма МКР №2: 30 питань}{Програма МКР №2: 30 питань}


Попередня програма: Електростатика.

Джерело: \textbf{т. 1, Частина II, розд. 14–19 (с. 298–484)}. Кожне питання зіставлено з параграфом за заголовками з OCR. Питання переписано з фото аркуша.
Статус:  заплановано ·  каркас · \ok{} формули звірено ·  написано від руки.

\needspace{5\baselineskip}\section*{1. Питання → підручник}

{\footnotesize
\begin{longtable}{>{\raggedright\arraybackslash}p{0.024\linewidth}>{\raggedright\arraybackslash}p{0.315\linewidth}>{\raggedright\arraybackslash}p{0.309\linewidth}>{\raggedright\arraybackslash}p{0.095\linewidth}>{\raggedright\arraybackslash}p{0.092\linewidth}>{\raggedright\arraybackslash}p{0.069\linewidth}}
\toprule
\textbf{№} & \textbf{Питання (як у програмі)} & \textbf{§ т. 1} & \textbf{с.} & \textbf{Сторінка} & \textbf{Статус} \\
\midrule
\endhead
1 & Закони ідеальних газів & 14.6, 14.7 & 311–317 & 01 &  \\
2 & Основне рівняння молекулярно-кінетичної теорії газу & 14.2, 14.3 (+14.4 температура) & 299–309 & 02 &  \\
3 & Внутрішня енергія термодинамічної системи & 16.1, 16.2 & 371–377 & 03 &  \\
4 & Теплота, робота, теплоємність & 16.1, 16.2 & 371–383 & 04 &  \\
5 & Теплоємність ідеальних газів & 16.5 & 390–397 & 05 &  \\
6 & Перший закон термодинаміки & 16.2 & 374–383 & 06 &  \\
7 & Ізотермічний процес & 16.4 (+14.7) & 386–390 & 07-09 &  \\
8 & Ізохоричний процес & 16.4 (+14.7) & 386–390 & 07-09 &  \\
9 & Ізобарний процес & 16.4 (+14.7) & 386–390 & 07-09 &  \\
10 & Адіабатичний процес & 16.3 «Адіабатний процес. Рівняння Пуассона» & 383–386 & 10 &  \\
11 & Колові процеси & 16.6 «Оборотні та необоротні процеси. Колові процеси (цикли)» & 397–399 & 11 &  \\
12 & Цикл Карно (необернений) та його ккд & 16.6, 16.7 & 399–403 & 12-13 &  \\
13 & Цикл Карно (обернений) та його ккд & 16.6, 16.7 (холодильна машина) & 399–403 & 12-13 &  \\
14 & Ентропія та її властивості & 16.8, 16.10 & 403–409, 411–414 & 14 &  \\
15 & Теорема Нернста & 16.11 «Третій закон термодинаміки» & 414–415 & 15 &  \\
16 & Другий закон термодинаміки & 16.8 «Другий закон термодинаміки. Теорема Карно» & 403–409 & 16 &  \\
17 & Закон розподілу Больцмана & 14.11 & 331–337 & 17,19 &  \\
18 & Закон розподілу Максвела & 14.8, 14.9, 14.10 & 317–331 & 18 &  \\
19 & Закон розподілу Максвела-Больцмана & 14.11 «Барометрична формула. Розподіл Максвелла — Больцмана» & 331–337 & 17,19 &  \\
20 & Закон рівномірного розподілу енергії по ступенях свободи & 16.2 (с. 376), 16.5 & 376, 390–392 & 20-21 &  \\
21 & Внутрішня енергія ідеальних газів & 16.5 & 390–392 & 20-21 &  \\
22 & Середня довжина вільного пробігу молекули в газах & 15.1 & 343–348 & 22 &  \\
23 & Дифузія в газах & 15.4 (+14.5) & 356–364 & 23-25 &  \\
24 & Внутрішнє тертя в газах & 15.2 «Внутрішнє тертя (в'язкість) у газах» & 348–352 & 23-25 &  \\
25 & Теплопровідність газів & 15.3 & 352–356 & 23-25 &  \\
26 & Відмінність реального газу від ідеального & 17.1 & 419–422 & 26 &  \\
27 & Рівняння Ван-Дер-Ваальса & 17.2, 17.3 (+17.4) & 422–433 & 27 &  \\
28 & Класифікація кристалів & 18.1 «Кристалічні та аморфні тіла» & 444–450 & 28 &  \\
29 & Дефекти в кристалах & 18.3 «Дефекти в реальних кристалах» & 453–456 & 29 &  \\
30 & Рідини (дописано від руки) & 19.1 (+19.2–19.4 необов'язково) & 462–474 & 30 &  \\
\bottomrule
\end{longtable}
}

Разом \textbf{30 питань → 23 сторінки}. Об'єднання (7–9, 12–13, 17+19, 20–21, 23–25) зводять на одну сторінку питання зі спільним методом чи формулою.

\needspace{5\baselineskip}\section*{2. Термінологія: книга vs програма}

На сторінці пишіть термін книги, а варіант програми згадайте один раз.

{\small
\begin{longtable}{>{\raggedright\arraybackslash}p{0.350\linewidth}>{\raggedright\arraybackslash}p{0.618\linewidth}}
\toprule
\textbf{Програма} & \textbf{Книга (т. 1)} \\
\midrule
\endhead
ступені свободи & \textbf{ступені вільності} (і) \\
адіабатичний процес & \textbf{адіабатний процес} \\
Максвела, Ван-Дер-Ваальса & \textbf{Максвелла}, \textbf{Ван-дер-Ваальса} \\
необернений / обернений цикл Карно & прямий / \textbf{обернений} цикл; слово «необернений» у книзі не вживається (с. 399) \\
внутрішнє тертя & внутрішнє тертя \textbf{(в'язкість)} \\
теорема Нернста & третій закон термодинаміки (теорема Нернста) \\
\bottomrule
\end{longtable}
}

\needspace{5\baselineskip}\section*{3. Порядок вивчення (не за номерами програми)}

Програма перескакує з теми на тему. Кожен блок нижче спирається на попередній, тож ідіть у такому порядку:

{\small
\begin{longtable}{>{\raggedright\arraybackslash}p{0.273\linewidth}>{\raggedright\arraybackslash}p{0.177\linewidth}>{\raggedright\arraybackslash}p{0.421\linewidth}>{\raggedright\arraybackslash}p{0.065\linewidth}}
\toprule
\textbf{Блок} & \textbf{Сторінки} & \textbf{Чому тут} & \textbf{Обсяг} \\
\midrule
\endhead
\textbf{A. Молекулярно-кінетична теорія} & 2 → 1 → 20–21 & тиск і температура газу через молекули; i, U = (i/2)νRT & 3 \\
\textbf{B. Статистичні розподіли} & 18 → 17,19 & Максвелл (швидкості) → Больцман (потенціальне поле) → разом & 2 \\
\textbf{C. Явища переносу} & 22 → 23–25 & спершу ⟨l⟩; далі одна спільна схема: D, η, λ ∝ ⅓⟨v⟩⟨l⟩ & 2 \\
\textbf{D. Перший закон} & 3 → 4 → 6 → 5 & U, Q, A → перший закон → $C_{V}$, $C_{p}$, рівняння Майєра & 4 \\
\textbf{E. Процеси} & 7–9 → 10 & таблиця: закон, A, Q, ΔU, C для кожного процесу; потім адіабата & 2 \\
\textbf{F. Цикли і другий закон} & 11 → 12–13 → 16 → 14 → 15 & цикл → Карно → II закон → ентропія → III закон & 5 \\
\textbf{G. Реальні гази} & 26 → 27 & відхилення → рівняння й ізотерми Ван-дер-Ваальса, критичний стан & 2 \\
\textbf{H. Конденсовані стани} & 28 → 29 → 30 & кристали → дефекти → рідини & 3 \\
\bottomrule
\end{longtable}
}

\chapter*{2. Основне рівняння молекулярно-кінетичної теорії газу}\addcontentsline{toc}{chapter}{2. Основне рівняння молекулярно-кінетичної теорії газу}\markboth{2. Основне рівняння молекулярно-кінетичної теорії газу}{2. Основне рівняння молекулярно-кінетичної теорії газу}

\srcline{\book{} т. 1, \textbf{§ 14.2} (с. 299–300), \textbf{§ 14.3} (с. 301–305), \textbf{§ 14.4} (с. 305–309). \ok{} (14.5)–(14.10) звірено зі сканом.}
Блок A, сторінка 1 з 3. Почніть звідси: і газові закони (01), і U = (i/2)RT (20–21) випливають із цієї сторінки.

\needspace{5\baselineskip}\section*{План сторінки}
\begin{layout}
┌ [A] Модель ідеального газу: 6 положень МКТ (стисло)        ┐
├ [B] Виведення тиску: куб l³, молекула б'ється об стінку, рис. 14.1 ┤
├ [C] Основне рівняння (14.6) у рамці + словами               ┤
├ [D] Температура: θ = ⅔⟨ε⟩, θ = kT → ⟨ε⟩ = 3/2 kT (14.10)     ┤
└ [E] p = n₀kT (14.13) → місток до сторінки 01                 ┘
\end{layout}

\needspace{5\baselineskip}\section*{[A] Ідеальний газ (с. 299–300)}
\begin{itemize}
\item Означення книги (с. 300): можна знехтувати розмірами молекул («вважати молекули матеріальними точками») і їхньою взаємодією на відстані, враховуючи лише зіткнення. «Такий газ називають \textbf{ідеальним}» (англ. ideal gas).
\item \textbf{Основні положення МКТ газів} (їх 6, с. 300). Запишіть коротким списком:
\begin{enumerate}
\item Молекули — однакові кульки нехтовно малих розмірів, без притягання й відштовхування.
\item Навіть у малому об'ємі дуже багато молекул, тож діють статистичні закономірності.
\item Зіткнення \textbf{пружні}.
\item Час між зіткненнями ≫ тривалості зіткнення, тож між зіткненнями рух рівномірний і прямолінійний.
\item Без зовнішніх сил молекули розподілені по об'єму рівномірно.
\item Усі напрями швидкостей рівноймовірні — \textbf{«молекулярний хаос»}.
\end{enumerate}
\item \textbf{Тепловий рух} (англ. thermal motion): абсолютно хаотичний рух молекул (с. 300).
\end{itemize}

\needspace{5\baselineskip}\section*{[B] Виведення (с. 301–303)}
\begin{itemize}
\item Куб із ребром l (рис. 14.1). Молекула вдаряється в праву стінку: Δp = 2m₀$v_{x.}$ Повертається через Δt = 2l/$v_{x}$, тож середня сила $f = m_0v_x^2/l$.
\item Сумуємо по всіх n молекулах і ділимо на S = l². Напрями рівноймовірні, отже $p_{x}$ = $p_{y}$ = $p_{z}$, звідки
$p = \dfrac23\dfrac1{l^3}\sum_{i=1}^n\dfrac{m_{0i}v_i^2}{2}$ (14.5)
\item \textbf{Середня квадратична швидкість} (англ. root-mean-square speed): $\langle v_{\text{кв}}^2\rangle = \dfrac{v_1^2+v_2^2+\ldots+v_n^2}{n}$
\item \pen{} Малюнок: куб, молекула відбивається від правої грані, стрілки $v_{x}$ → і −$v_{x}$ ←.
\end{itemize}

\needspace{5\baselineskip}\section*{[C] Основне рівняння МКТ (с. 303) — у рамку}
\[p = \frac23\,n_0\,\frac{m_0\langle v_{\text{кв}}^2\rangle}{2}\qquad(14.6)\]
\begin{itemize}
\item n₀ = n/V — \textbf{кількість молекул в одиниці об'єму}; $m_0\langle v_{\text{кв}}^2\rangle/2$ — \textbf{середня кінетична енергія поступального руху однієї молекули}.
\item Словами (с. 303): «тиск ідеального газу дорівнює двом третинам середньої кінетичної енергії поступального руху молекул одиниці об'єму газу» (англ. basic equation of kinetic theory)
\item Чим важливе: пов'язує \textbf{мікроскопічні} величини (m₀, v) з \textbf{макроскопічною} (p), яку можна виміряти.
\item Тиск — \textbf{статистичний параметр} (англ. statistical parameter): для однієї молекули він не має сенсу (с. 304).
\item Одиниці: Па = Н/м². Також 1 атм = 101 325 Па; 1 тор = 133,322 Па; 1 бар = 10⁵ Па.
\end{itemize}

\needspace{5\baselineskip}\section*{[D] Температура (с. 305–308)}
\begin{itemize}
\item \textbf{Термодинамічна рівновага} (англ. thermodynamic equilibrium): припиняються всі макроскопічні зміни. \textbf{Релаксація} (англ. relaxation) — самочинний перехід до неї; \textbf{час релаксації} — його тривалість. Існування рівноваги — це \textbf{нульове начало термодинаміки} (англ. zeroth law), с. 306.
\item Приклад із поршнем (рис. 14.2): у рівновазі середні кінетичні енергії по обидва боки вирівнюються. Отже, температура «поводиться як» ⟨ε⟩ і є \textbf{інтенсивним параметром} (англ. intensive parameter; не залежить від маси системи).
\item Означення в енергетичних одиницях: $\Theta = \dfrac23\dfrac{m_0\langle v_{\text{кв}}^2\rangle}{2}$, звідки $p = n_0\Theta$ (14.8)
\item У кельвінах: $\Theta = kT$ (14.9), де k — \textbf{стала Больцмана} (англ. Boltzmann constant), 1,38·10⁻²³ Дж/К.
\item \[\frac{m_0\langle v_{\text{кв}}^2\rangle}{2} = \frac32kT\qquad(14.10)\]
\item Зміст k: на скільки змінюється енергія поступального руху однієї молекули при зміні T на 1 К.
\end{itemize}

\needspace{5\baselineskip}\section*{[E] Наслідок (с. 311)}
\begin{itemize}
\item Поєднуючи (14.6) і (14.10): $p = n_0kT$ (14.13). Це крок до 01 — газові закони.
\end{itemize}

\needspace{5\baselineskip}\section*{Самоперевірка}
\begin{enumerate}
\item Де у виведенні використано «пружні зіткнення» і «хаос»?
\item Чому не можна говорити про тиск однієї молекули?
\item Що стане з ⟨v²⟩, якщо T подвоїти? (Зросте вдвічі, отже $v_{\text{кв}}$ — у √2 раза.)
\end{enumerate}

\chapter*{1. Закони ідеальних газів}\addcontentsline{toc}{chapter}{1. Закони ідеальних газів}\markboth{1. Закони ідеальних газів}{1. Закони ідеальних газів}

\srcline{\book{} т. 1, \textbf{§ 14.6} «Рівняння стану ідеального газу» (с. 311–313), \textbf{§ 14.7} «Основні закони ідеального газу» (с. 313–317). \ok{} (14.13)–(14.18) звірено зі сканом; (14.22)–(14.31) — з чіткого OCR.}
Блок A, сторінка 2 з 3. Підхід книги: \textbf{дослідні} закони виводяться \textbf{як наслідки} p = n₀kT.

\needspace{5\baselineskip}\section*{План сторінки}
\begin{layout}
┌ [A] Рівняння стану: pV = nkT → Клапейрон → Клапейрон–Менделєєв ┐
├ [B] Таблиця трьох ізопроцесів (закон, рівняння, графік)         ┤
│     + три маленькі графіки p–V, V–T, p–T (рис. 14.5–14.7)       │
├ [C] Закон Авогадро                                               ┤
└ [D] Межі застосування моделі                                    ┘
\end{layout}

\needspace{5\baselineskip}\section*{[A] Рівняння стану ідеального газу (с. 311–313)}
\begin{itemize}
\item \textbf{Термічне рівняння стану} (англ. thermal equation of state): рівняння, що пов'язує p, V, T у стані теплової рівноваги, $F(p, V, T) = 0$ (14.12).
\item З МКТ: $p = n_0kT$ (14.13) ⟹ $pV = nkT$ (14.14), де n — кількість молекул.
\item Для одного моля: $pV_M = N_AkT$ (14.15). З \textbf{молярною газовою сталою} (англ. molar gas constant) $R = kN_A = 8{,}31$ Дж/(моль·К):
$pV_M = RT$ (14.16) — \textbf{рівняння Клапейрона}.
\item Для довільної кількості речовини, ν = m/M молів:
\[pV = \frac mM RT,\quad\text{\text{або}}\quad pV = \frac{n}{N_A}RT\qquad(14.18)\]
— \textbf{рівняння Клапейрона–Менделєєва} (англ. ideal gas law).
\item Зміст R (с. 313): робота розширення 1 моля ідеального газу при \textbf{ізобарному} нагріванні на 1 К. \pen{} Намалюйте поршень із рис. 14.4.
\end{itemize}

\needspace{5\baselineskip}\section*{[B] Ізопроцеси (m = const) (с. 313–316)}
{\footnotesize
\begin{longtable}{>{\raggedright\arraybackslash}p{0.184\linewidth}>{\raggedright\arraybackslash}p{0.058\linewidth}>{\raggedright\arraybackslash}p{0.207\linewidth}>{\raggedright\arraybackslash}p{0.227\linewidth}>{\raggedright\arraybackslash}p{0.244\linewidth}}
\toprule
\textbf{процес} & \textbf{умова} & \textbf{закон (назва в книзі)} & \textbf{рівняння} & \textbf{графік} \\
\midrule
\endhead
\textbf{ізотермічний} & T = const & закон \textbf{Бойля–Маріотта} & pV = const (14.22) & p–V: гіперболи (рис. 14.5) \\
\textbf{ізобарний} & p = const & закон \textbf{Гей-Люссака} & V/T = const (14.25); V = V₀(1 + αt) (14.26) & V–T: прямі через початок координат (рис. 14.6) \\
\textbf{ізохорний} & V = const & «рівняння ізохорного процесу» & p/T = const (14.28); p = p₀(1 + βt) (14.29) & p–T: прямі через початок координат (рис. 14.7) \\
\bottomrule
\end{longtable}
}

\begin{itemize}
\item α = β = 0,003661 К⁻¹ ≈ 1/273 — \textbf{коефіцієнт об'ємного розширення} і \textbf{термічний коефіцієнт тиску}.
\item Необов'язково: ізотермічний коефіцієнт стисливості $\chi = -\dfrac1V\left(\dfrac{\partial V}{\partial p}\right)_T = \dfrac1p$ (14.24) і зв'язок $\beta = \alpha/(p\chi)$ (14.31).
\item У підручнику ізохорний закон \textbf{не названо законом Шарля}. Якщо лектор так його називає, допишіть самі (англ. Charles's law).
\item \pen{} Намалюйте три маленькі графіки з кількома кривими (T₁ < T₂, p₁ < p₂, V₁ < V₂) і позначте, яка крива відповідає більшому параметру.
\item Тут лише закони. Робота, теплота і ΔU для кожного процесу — у блоці E (07–09).
\end{itemize}

\needspace{5\baselineskip}\section*{[C] Закон Авогадро (с. 316)}
\begin{itemize}
\item З pV = nkT: «в однакових об'ємах ідеальних газів при однакових температурах і тисках міститься однакове число молекул» (англ. Avogadro's law)
\end{itemize}

\needspace{5\baselineskip}\section*{[D] Межі застосування (с. 313)}
\begin{itemize}
\item Модель справджується, коли власним об'ємом молекул і їхньою взаємодією можна знехтувати: \textbf{низький тиск і досить висока температура}. Нормальні атмосферні умови здебільшого підходять. Далі — 26, реальний газ.
\end{itemize}

\needspace{5\baselineskip}\section*{Необов'язково (приклад 14.1, с. 316)}
\begin{itemize}
\item Вивести pV/T = const лише з дослідних законів: ізотерма 1→1′, потім ізохора 1′→2. Гарний прийом для іспиту, займає пів сторінки.
\end{itemize}

\needspace{5\baselineskip}\section*{Самоперевірка}
\begin{enumerate}
\item Для якого рівняння потрібна кількість молекул, а для якого достатньо маси?
\item Чому ізобари на графіку V–T проходять через початок координат (у кельвінах)?
\item Який фізичний зміст R?
\end{enumerate}

\chapter*{20. Закон рівномірного розподілу енергії по ступенях свободи · 21. Внутрішня енергія ідеальних газів}\addcontentsline{toc}{chapter}{20. Закон рівномірного розподілу енергії по ступенях свободи · 21. Внутрішня енергія ідеальних газів}\markboth{20. Закон рівномірного розподілу енергії по ступенях свободи}{20. Закон рівномірного розподілу енергії по ступенях свободи}

\srcline{\book{} т. 1, \textbf{§ 16.2} (с. 375–377, у складі «Перший закон термодинаміки»), \textbf{§ 16.5} «Теплоємність ідеального газу» (с. 390–392). \ok{} (16.3), (16.4), (16.46) і значення $C_{V}$ звірено зі сканом.}
Блок A, сторінка 3 з 3. \warn{} Термінологія: у книзі \textbf{«ступені вільності»}, у програмі — «ступені свободи». Пишіть за книгою, а варіант програми згадайте один раз.

\needspace{5\baselineskip}\section*{План сторінки}
\begin{layout}
┌ [A] Ступені вільності: означення, i = i_пост + i_оберт + i_колив (16.4) ┐
├ [B] Закон рівного розподілу (рамка): ½kT на ступінь; коливальний → kT ┤
├ [C] Таблиця: газ → i → U → C_V, C_p + малюнки молекул (1, 2, 3 атоми)  ┤
├ [D] Внутрішня енергія ідеального газу: U = (i/2)RT; тільки від T      ┤
└ [E] Межі класичної теорії (спирт: теорія 3R, дослід ≈ 8R)            ┘
\end{layout}

\needspace{5\baselineskip}\section*{[A] Ступені вільності (с. 376)}
\begin{itemize}
\item «Число таких незалежних координат називають \textbf{кількістю ступенів вільності} й позначають i» (англ. degrees of freedom). Ідеться про координати, потрібні, щоб визначити положення всіх точок системи.
\item Вільна матеріальна точка: i = 3 (поступальний рух). Вільне тверде тіло: i = 6 (3 поступальні + 3 обертальні).
\item $i = i_{\text{пост}} + i_{\text{оберт}} + i_{\text{колив}}$ (16.4)
\end{itemize}

\needspace{5\baselineskip}\section*{[B] Закон рівного розподілу (с. 377) — у рамку}
\begin{itemize}
\item Назва в книзі: \textbf{«закон рівного розподілу енергії за класичними ступенями вільності (закон рівного розподілу Больцмана)»} (англ. equipartition theorem). На с. 392 його також названо «закон про \textbf{рівномірний} розподіл енергії за ступенями вільності» — це близько до формулювання програми.
\item Формулювання (с. 377): «у стані теплової рівноваги на кожний класичний ступінь вільності руху частинок системи припадає кінетична енергія, що дорівнює в середньому \textbf{½kT}»
\item Обертальний ступінь вільності: ½kT. \textbf{Коливальний: kT} — це ½kT кінетичної плюс ½kT потенціальної енергії.
\item Загальна формула: $\langle w\rangle = (3 + r + 2s)\dfrac{kT}{2}$ (16.45), де r — кількість обертальних, s — коливальних ступенів вільності.
\item Порушується там, де суттєві квантові ефекти (с. 377).
\end{itemize}

\needspace{5\baselineskip}\section*{[C] Таблиця (с. 391–392), жорсткі молекули (s = 0)}
{\footnotesize
\begin{longtable}{>{\raggedright\arraybackslash}p{0.206\linewidth}>{\raggedright\arraybackslash}p{0.238\linewidth}>{\raggedright\arraybackslash}p{0.143\linewidth}>{\raggedright\arraybackslash}p{0.159\linewidth}>{\raggedright\arraybackslash}p{0.079\linewidth}>{\raggedright\arraybackslash}p{0.079\linewidth}}
\toprule
\textbf{газ} & \textbf{приклади} & \textbf{i} & \textbf{U (1 моль)} & \textbf{$C_{V}$} & \textbf{$C_{p}$} \\
\midrule
\endhead
одноатомний & He, Ar, пара Hg & 3 & 3/2 RT & 3/2 R & 5/2 R \\
двохатомний & H₂, O₂, N₂ & 5 (3 + 2) & 5/2 RT & 5/2 R & 7/2 R \\
багатоатомний & пара H₂O, CH₄ & 6 (3 + 3) & 3RT & 3R & 4R \\
\bottomrule
\end{longtable}
}
\begin{itemize}
\item $C_{p}$ = $C_{V}$ + R. Саме \textbf{рівняння Майєра} — у темі 5; тут просто скористайтеся ним.
\item \pen{} Намалюйте три маленькі молекули (точка, «гантель», трикутник) зі стрілками осей обертання. «Гантель» має лише 2 осі обертання: обертання навколо власної осі не враховується.
\item Необов'язково: CO₂ (лінійна) дає ⟨w⟩ = 6,5kT, а H₂O (нелінійна) — 6kT з урахуванням коливань (16.43–16.44).
\end{itemize}

\needspace{5\baselineskip}\section*{[D] Внутрішня енергія ідеального газу (с. 376, 390–392)}
\begin{itemize}
\item Загалом U — це \textbf{кінетична енергія теплового руху + потенціальна енергія взаємодії} (с. 375). Докладно — у темі 3.
\item В ідеальному газі \textbf{взаємодії немає}, тож U — лише кінетична енергія: «внутрішня енергія ідеального газу дорівнює сумарній кінетичній енергії його молекул» (с. 390) (англ. internal energy).
\item Один моль одноатомного газу: $U = \dfrac32RT$ (16.2). Загалом:
\[U = \frac i2RT\qquad(16.3)\]
а для ν молів $U = \nu\dfrac i2RT$.
\item Головний висновок (с. 392): «внутрішня енергія ідеального газу є функцією \textbf{тільки абсолютної температури}». Від V вона не залежить. На цьому ґрунтуються процеси блоку E.
\end{itemize}

\needspace{5\baselineskip}\section*{[E] Межі класичної теорії (с. 392)}
\begin{itemize}
\item Для простих молекул поблизу кімнатної температури теорія узгоджується з дослідом. Для складних, як C₂H₅OH, теорія дає $C_{V}$ = 3R, а дослід — близько 8R.
\item $C_{V}$ = 5/2 R для двохатомних газів справджується лише в обмеженому інтервалі температур, бо при високих T «вмикаються» коливання. Це квантовий ефект.
\end{itemize}

\needspace{5\baselineskip}\section*{Самоперевірка}
\begin{enumerate}
\item Чому молекула-«гантель» має лише 2 обертальні ступені вільності?
\item Чому на коливальний ступінь вільності припадає kT, а не ½kT?
\item Чому U ідеального газу не залежить від об'єму?
\end{enumerate}

\chapter*{18. Закон розподілу Максвела}\addcontentsline{toc}{chapter}{18. Закон розподілу Максвела}\markboth{18. Закон розподілу Максвела}{18. Закон розподілу Максвела}

\srcline{\book{} т. 1, \textbf{§ 14.8} «Швидкості газових молекул та вимірювання їх» (с. 317–320), \textbf{§ 14.9} «Поняття про статистичний розподіл. Функція розподілу. Розподіл молекул за швидкостями (розподіл Максвелла)» (с. 320–330), \textbf{§ 14.10} (с. 331). \ok{} (14.54)–(14.63) звірено зі сканом.}
Блок B, сторінка 1 з 2. Правопис: у книзі \textbf{Максвелла} (з подвоєним «л»), у програмі — «Максвела».

\needspace{5\baselineskip}\section*{План сторінки}
\begin{layout}
┌ [A] Постановка: чому «скільки молекул в інтервалі v…v+dv», а не «з швидкістю v» ┐
├ [B] Функція розподілу f(v): означення (14.42), нормування (14.43)               ┤
├ [C] Формула Максвелла (14.57) у рамці + ГРАФІК f(v) для T₁ < T₂ (рис. 14.13)      ┤
├ [D] Три швидкості: v_н, ⟨v⟩, v_кв (14.58), (14.61), (14.62) + співвідношення     ┤
└ [E] Дослідна перевірка: Штерн (рис. 14.9), Ламмерт (рис. 14.15)                  ┘
\end{layout}

\needspace{5\baselineskip}\section*{[A] Постановка задачі (с. 320)}
\begin{itemize}
\item Швидкість — неперервна величина, тож імовірність \textit{точного} значення дорівнює 0. Правильне запитання: «скільки з усіх наявних молекул мають швидкості, значення яких лежать у певному інтервалі швидкостей в околі заданої швидкості?»
\item \textbf{Простір швидкостей} (англ. velocity space): кожна молекула — точка ($v_{x}$, $v_{y}$, $v_{z}$), рис. 14.11.
\end{itemize}

\needspace{5\baselineskip}\section*{[B] Функція розподілу (с. 322–323)}
\begin{itemize}
\item $dW = \dfrac{dn}{n} = f(\vec v)\,d\omega$ (14.42). f — \textbf{функція розподілу молекул за швидкостями} (англ. velocity distribution function), її також називають \textbf{густиною ймовірності} (англ. probability density).
\item \textbf{Умова нормування} (англ. normalization): $\int f\,d\omega = 1$ (14.43) — молекула напевно має \textit{якусь} швидкість.
\item Напрями рівноймовірні, тож f залежить лише від |v|. Компоненти незалежні, тож $f = f(v_x)f(v_y)f(v_z)$ (14.46). Максвелл (1859).
\item Виведення (с. 323–327) пропустіть; запишіть лише ідею: «незалежність + ізотропність ⟹ гауссіана».
\end{itemize}

\needspace{5\baselineskip}\section*{[C] Закон Максвелла (с. 327) — у рамку}
\[\frac{dn}{n} = \frac{4}{\sqrt\pi}\left(\frac{m_0}{2kT}\right)^{3/2}e^{-\frac{m_0v^2}{2kT}}v^2\,dv\qquad(14.57)\]
\begin{itemize}
\item Словами: \textbf{відносна кількість молекул}, швидкості яких лежать між v і v + dv (англ. Maxwell distribution).
\item Для однієї компоненти: $f(v_x) = \dfrac{1}{\sqrt\pi\,v_\text{н}}e^{-v_x^2/v_\text{н}^2}$ (14.54) — «дзвін» із центром у 0 (рис. 14.14).
\item Форма f(v) (с. 327–328): f = 0 при v = 0 (нерухомих молекул немає) і f → 0 при v → ∞. Крива \textbf{асиметрична}: круто зростає й повільно спадає. Площа під нею дорівнює 1.
\item \textbf{Нагрівання} (с. 328): максимум зміщується вправо (∝ √T), стає нижчим і ширшим, бо площа лишається 1.
\item \pen{} Головний малюнок: дві криві для T₁ < T₂ з однаковою площею; на одній позначте $v_{\text{н}}$, ⟨v⟩ і $v_{\text{кв}.}$ Заштрихуйте одну смужку dv як dn/n.
\item Приклад із книги (табл. 3, O₂ при 0 °C): близько ¾ молекул мають швидкості 200–600 м/с.
\end{itemize}

\needspace{5\baselineskip}\section*{[D] Характерні швидкості (с. 328–330)}
{\small
\begin{longtable}{>{\raggedright\arraybackslash}p{0.565\linewidth}>{\raggedright\arraybackslash}p{0.296\linewidth}>{\raggedright\arraybackslash}p{0.091\linewidth}}
\toprule
\textbf{швидкість} & \textbf{формула} & \textbf{№} \\
\midrule
\endhead
\textbf{найбільш імовірна} $v_{\text{н}}$ (максимум f) (англ. most probable speed) & $\sqrt{2kT/m_0} = \sqrt{2RT/M}$ & (14.58) \\
\textbf{середня (арифметична)} ⟨v⟩ (англ. mean speed) & $\sqrt{8kT/(\pi m_0)}$ & (14.61) \\
\textbf{середня квадратична} $v_{\text{кв}}$ (англ. rms speed) & $\sqrt{3kT/m_0}$ & (14.62) \\
\bottomrule
\end{longtable}
}
\begin{itemize}
\item $v_{\text{кв}} : \langle v\rangle : v_\text{н} \approx 1{,}2 : 1{,}1 : 1$
\item Універсальна форма з u = v/$v_{\text{н}}$: $\dfrac{dn}{n\,du} = \dfrac4{\sqrt\pi}e^{-u^2}u^2$ (14.63) — не залежить ні від газу, ні від T.
\item Порядок величин (с. 318): H₂ ≈ 1,9 км/с, N₂ ≈ 493 м/с, O₂ ≈ 461 м/с при кімнатній температурі.
\end{itemize}

\needspace{5\baselineskip}\section*{[E] Дослідна перевірка (с. 318–320, 331)}
\begin{itemize}
\item \textbf{Дослід Штерна (1920)}: посріблений дріт усередині двох циліндрів, що обертаються. Смужка осаду зміщується й \textbf{розмивається} — отже, атоми мають розкид швидкостей із максимумом. \pen{} Намалюйте рис. 14.9 (вид згори).
\item \textbf{Ламмерт (1929)}: два диски, що обертаються, зі щілинами, зсунутими на кут φ, — селектор швидкостей. Підтвердив закон Максвелла з високою точністю.
\end{itemize}

\needspace{5\baselineskip}\section*{Самоперевірка}
\begin{enumerate}
\item Чому ймовірність точного значення швидкості дорівнює нулю?
\item Чому при нагріванні максимум кривої знижується?
\item Розташуйте $v_{\text{н}}$, ⟨v⟩ і $v_{\text{кв}}$ за зростанням і поясніть чому.
\end{enumerate}

\chapter*{17. Закон розподілу Больцмана · 19. Закон розподілу Максвела-Больцмана}\addcontentsline{toc}{chapter}{17. Закон розподілу Больцмана · 19. Закон розподілу Максвела-Больцмана}\markboth{17. Закон розподілу Больцмана}{17. Закон розподілу Больцмана}

\srcline{\book{} т. 1, \textbf{§ 14.11} «Барометрична формула. Розподіл Максвелла — Больцмана» (с. 331–337). \ok{} (14.64)–(14.75) звірено зі сканом.}
Блок B, сторінка 2 з 2. Логіка: Максвелл вважав, що \textbf{зовнішнього поля немає}. Додаємо тяжіння → барометрична формула → Больцман (будь-яке потенціальне поле) → поєднання з Максвеллом.

\needspace{5\baselineskip}\section*{План сторінки}
\begin{layout}
┌ [A] Барометрична формула: стовп повітря, dp = −ρg dh → (14.68) + рис. 14.16 ┐
├ [B] Розподіл Больцмана (14.71) у рамці + графік n(h) (рис. 14.18)          ┤
├ [C] Розподіл Максвелла–Больцмана (14.74): два незалежні множники            ┤
└ [D] Наслідки: атмосфери планет, флуктуації (необов'язково)                ┘
\end{layout}

\needspace{5\baselineskip}\section*{[A] Барометрична формула (с. 331–333)}
\begin{itemize}
\item Форму атмосфери визначає лише спільна дія \textbf{теплового руху + тяжіння}. Без тяжіння молекули розлетілися б, без теплового руху — впали б на землю.
\item Тонкий шар dh: $dp = -\rho g\,dh$ (14.64), де ρ = m₀n₀ і n₀ = p/kT ⟹ $\dfrac{dp}{p} = -\dfrac{m_0g}{kT}dh$ (14.66)
\item Інтегруємо:
\[p = p_0e^{-\frac{m_0g(h-h_0)}{kT}}\ (14.67),\qquad p = p_0e^{-\frac{Mg(h-h_0)}{RT}}\ (14.68)\]
— \textbf{барометрична формула} (англ. barometric formula; Лаплас, 1821).
\item Припущення: ідеальний газ, \textbf{T = const}, \textbf{g = const}, тож формула точна лише до кількох кілометрів.
\item Тиск спадає експоненціально — швидше для важчих газів (більша M) і холодніших (менша T).
\item \pen{} Намалюйте рис. 14.16: стовп повітря, шар dh, p і p − dp.
\end{itemize}

\needspace{5\baselineskip}\section*{[B] Розподіл Больцмана (с. 335) — у рамку}
\begin{itemize}
\item Оскільки p = nkT: $n = n_0e^{-\frac{m_0g(h-h_0)}{kT}}$ (14.70)
\item Загальна форма:
\[n = n_0e^{-\frac{\Delta U}{kT}}\qquad(14.71)\]
де ΔU — різниця \textbf{потенціальних енергій} молекули (у полі тяжіння ΔU = m₀gh).
\item Больцман показав, що закон \textbf{універсальний}: справджується для будь-яких частинок у тепловій рівновазі в \textbf{будь-якому потенціальному полі}. «Цей статистичний закон називається \textbf{розподілом Больцмана}» (англ. Boltzmann distribution)
\item Зміст: концентрація більша там, де менша U, — це тенденція до стійкої рівноваги.
\item У шарі: $\dfrac{dn}{n} = e^{-\frac{m_0g}{kT}dh}$ (14.73). \warn{} Переписано так, як надруковано на с. 336, але так бути не може: dh стоїть у показнику експоненти. Частка молекул у шарі h…h + dh: $\dfrac{dn}{n} = \dfrac{m_0g}{kT}e^{-\frac{m_0gh}{kT}}dh$.
\item \pen{} Графік рис. 14.18: n від h — експоненціальне спадання від n₀.
\end{itemize}

\needspace{5\baselineskip}\section*{[C] Розподіл Максвелла–Больцмана (с. 336)}
\begin{itemize}
\item Підставляємо n із (14.71) у формулу Максвелла (14.57) (англ. Maxwell–Boltzmann distribution):
\[dn = \frac4{\sqrt\pi}n_0\left(\frac{m_0}{2kT}\right)^{3/2}e^{-\frac{m_0\Delta\varphi}{kT}}e^{-\frac{m_0v^2}{2kT}}v^2\,dv\qquad(14.74)\]
де φ = gh — потенціал.
\item \textbf{Два незалежні множники}: больцманівський $e^{-m_0\Delta\varphi/kT}$ (положення) × максвеллівський $e^{-m_0v^2/2kT}$ (швидкість).
\item Отже, розподіл Больцмана не залежить від розподілу за швидкостями, а розподіл Максвелла — від потенціалу.
\item \pen{} Як запам'ятати: $e^{-(E_\text{п} + E_\text{к})/kT}$ — енергія, поділена на kT, у показнику.
\end{itemize}

\needspace{5\baselineskip}\section*{[D] Необов'язково}
\begin{itemize}
\item Атмосфери планет (с. 333–334): другу космічну швидкість перевищують лише молекули з «хвоста» розподілу Максвелла, тому легкі небесні тіла (Місяць) втрачають атмосферу.
\item Флуктуації (Смолуховський, с. 335): $\delta = \sqrt{\dfrac2\pi\dfrac1n}$ (14.72) — тим менші, чим більше частинок.
\item Дослід Перрена (§ 14.12) використав закон Больцмана, щоб виміряти $N_{A.}$
\end{itemize}

\needspace{5\baselineskip}\section*{Самоперевірка}
\begin{enumerate}
\item Які два «чинники» визначають форму атмосфери?
\item Чому барометрична формула хибна на великих висотах? (Змінюються T і g.)
\item Що означає, що розподіл Максвелла–Больцмана розпадається на два множники?
\end{enumerate}

\chapter*{22. Середня довжина вільного пробігу молекули в газах}\addcontentsline{toc}{chapter}{22. Середня довжина вільного пробігу молекули в газах}\markboth{22. Середня довжина вільного пробігу молекули в газах}{22. Середня довжина вільного пробігу молекули в газах}

\srcline{\book{} т. 1, \textbf{§ 15.1} «Середня довжина вільного пробігу молекул. Число зіткнень» (с. 343–348). \ok{} (15.1), (15.2) звірено зі сканом.}
Блок C, сторінка 1 з 2. \warn{} Позначення: у книзі довжина вільного пробігу — \textbf{⟨l⟩}; λ далі позначає теплопровідність.

\needspace{5\baselineskip}\section*{План сторінки}
\begin{layout}
┌ [A] Чому важливо: зигзаг, зіткнення → рівновага                  ┐
├ [B] Модель: «циліндр» πd²⟨v⟩ (рис. 15.1) → ⟨z⟩ (15.1) → ⟨l⟩ (15.2) ┤
├ [C] Ефективний переріз σ = πd²; ефективний діаметр (Сезерленд)    ┤
├ [D] Залежності: ⟨l⟩ ~ 1/p; числа для N₂                          ┤
└ [E] Послаблення пучка n = n₀e^(−x/⟨l⟩) (15.5); дослід Борна        ┘
\end{layout}

\needspace{5\baselineskip}\section*{[A] Навіщо (с. 343)}
\begin{itemize}
\item Молекули рухаються зі швидкостями в сотні м/с, але \textbf{зигзагом}, тож далеко не просуваються. Зіткнення встановлюють рівновагу й розподіл Максвелла.
\item \textbf{Середня довжина вільного пробігу} ⟨l⟩ (англ. mean free path): «середню відстань, яку проходить молекула між двома послідовними зіткненнями»
\item Модель: однакові абсолютно пружні кульки діаметра d; враховують лише \textbf{парні} зіткнення.
\end{itemize}

\needspace{5\baselineskip}\section*{[B] Виведення (с. 343–344)}
\begin{itemize}
\item «Заморозимо» всі молекули, крім однієї. За 1 с вона описує циліндр з основою πd² і довжиною ⟨v⟩ і зіткнеться з усіма молекулами, центри яких усередині: ⟨z⟩ = πd²n₀⟨v⟩.
\item Інші молекули теж рухаються, тож береться відносна швидкість $\langle v\rangle_{\text{відн}} = \sqrt2\langle v\rangle$ (з розподілу Максвелла):
\[\langle z\rangle = \sqrt2\,\pi d^2n_0\langle v\rangle\qquad(15.1)\]
\[\langle l\rangle = \frac{\langle v\rangle}{\langle z\rangle} = \frac{1}{\sqrt2\,\pi d^2n_0}\qquad(15.2)\]
\item \pen{} Намалюйте рис. 15.1: молекула, «випрямлений» циліндр радіуса d і молекули з центрами всередині.
\end{itemize}

\needspace{5\baselineskip}\section*{[C] Ефективний переріз (с. 344–345)}
\begin{itemize}
\item \textbf{Повний ефективний поперечний переріз} (англ. effective cross-section) σ = πd². Вимірюється в м², але геометричного змісту не має: характеризує \textbf{імовірність} процесу зіткнення.
\item \textbf{Ефективний діаметр} зменшується з ростом T (Сезерленд): $d^2 = d_0^2(1 + C/T)$ (15.3), тож ⟨l⟩ зростає з T.
\end{itemize}

\needspace{5\baselineskip}\section*{[D] Порядки величин і залежності (с. 345)}
\begin{itemize}
\item N₂ за нормальних умов: ⟨l⟩ ≈ 0,6·10⁻⁷ м — \textbf{приблизно в 1000 разів більше за розмір молекули}. Навіть при атмосферному тиску зіткнення рідкісні.
\item При T = const \textbf{⟨l⟩ ∝ 1/p}. Коли ⟨l⟩ сягає розмірів посудини, газ перебуває у стані \textbf{вакууму} (§ 15.5).
\end{itemize}

\needspace{5\baselineskip}\section*{[E] Послаблення пучка (с. 345–347), необов'язково}
\begin{itemize}
\item $-\dfrac{dn}{n} = \dfrac{dx}{\langle l\rangle}$ (15.4) ⟹ $n = n_0e^{-x/\langle l\rangle}$ (15.5). Далі ніж ⟨l⟩ пролітає лише ≈ 1/e ≈ ⅓ молекул.
\item Дослід Борна (1921): пучок атомів срібла крізь чотири діафрагми дає ⟨l⟩ безпосередньо (15.9).
\end{itemize}

\needspace{5\baselineskip}\section*{Самоперевірка}
\begin{enumerate}
\item Звідки √2 у (15.1)?
\item Чому ⟨l⟩ не залежить від T при сталому n₀? (Без поправки Сезерленда.)
\item Що таке «вакуум» мовою ⟨l⟩?
\end{enumerate}

\chapter*{23. Дифузія в газах · 24. Внутрішнє тертя в газах · 25. Теплопровідність газів}\addcontentsline{toc}{chapter}{23. Дифузія в газах · 24. Внутрішнє тертя в газах · 25. Теплопровідність газів}\markboth{23. Дифузія в газах}{23. Дифузія в газах}

\srcline{\book{} т. 1, \textbf{§ 15.2} «Внутрішнє тертя (в'язкість) у газах» (с. 348–352), \textbf{§ 15.3} «Теплопровідність газів» (с. 352–356), \textbf{§ 15.4} «Дифузія газів» (с. 356–363). \ok{} (15.10), (15.14), (15.16)–(15.27) звірено зі сканом.}
Блок C, сторінка 2 з 2. Одна ідея, тричі. Див. сторінку-поняття явища переносу.

\needspace{5\baselineskip}\section*{План сторінки (найкраще — таблиця з трьох колонок зі спільним верхом)}
\begin{layout}
┌ [A] Що таке явища переносу: перенос маси / імпульсу / енергії; спільна модель ┐
│     (площинка ΔS, 1/6 n₀⟨v⟩ в кожен бік, «пам'ять» на відстані ⟨l⟩) + рис. 15.3  │
├────────────────┬────────────────────┬──────────────────────────────────────┤
│ [B] Дифузія    │ [C] В'язкість      │ [D] Теплопровідність                   │
│ Фік (15.22)    │ Ньютон (15.10)     │ Фур'є (15.16)                          │
│ D = ⅓⟨v⟩⟨l⟩    │ η = ⅓ρ⟨v⟩⟨l⟩       │ λ = ⅓ρ⟨v⟩⟨l⟩c_V                        │
├────────────────┴────────────────────┴──────────────────────────────────────┤
└ [E] Зв'язок: η = Dρ, λ = ηc_V; залежність від p і T                          ┘
\end{layout}

\needspace{5\baselineskip}\section*{[A] Явища переносу (с. 348, 349, 353, 357)}
\begin{itemize}
\item Означення книги (с. 348): явища, зумовлені тепловим рухом, що переносять \textbf{масу (дифузія), імпульс (внутрішнє тертя), енергію (теплопровідність)} (англ. transport phenomena). Вони ведуть систему до рівноваги.
\item Спільне виведення:
\begin{itemize}
\item Беремо малу площинку ΔS ⟂ OX, градієнт — уздовж OX.
\item У кожен бік за одиницю часу її перетинає $\tfrac16n_0\langle v\rangle$ молекул.
\item Кожна несе значення величини, яке мала при \textbf{останньому зіткненні}, на відстані ⟨l⟩.
\item Результуючий потік ∝ різниці значень у x − ⟨l⟩ і x + ⟨l⟩, тобто градієнт × 2⟨l⟩.
\end{itemize}
\item \pen{} Один спільний малюнок (рис. 15.3/15.5/15.7): площинка ΔS, ⟨l⟩ ліворуч і праворуч, величина «більша» зліва.
\end{itemize}

\needspace{5\baselineskip}\section*{[B] Дифузія (23) (с. 356–358)}
\begin{itemize}
\item «Дифузія газів — це процес самовільного вирівнювання густин (концентрацій) у суміші, зумовлений тепловим рухом мікрочастинок» (англ. diffusion). В однокомпонентному газі — \textbf{самодифузія} (англ. self-diffusion). Процес \textbf{необоротний}.
\item \textbf{Закон Фіка} (перший) (англ. Fick's law): $j_m = -D\dfrac{d\rho}{dx}$ (15.22), D — \textbf{коефіцієнт дифузії}, м²/с.
\item З кінетичної теорії: $D = \tfrac13\langle v\rangle\langle l\rangle$ (15.27)
\item D ∝ 1/p при T = const і зростає з T.
\item Необов'язково: нестаціонарна дифузія $\Delta n = \Delta n_0e^{-t/\tau}$ (15.29); другий закон Фіка $\partial n/\partial t = D\,\partial^2n/\partial x^2$ (15.30); взаємна дифузія й термодифузія (с. 360–363).
\end{itemize}

\needspace{5\baselineskip}\section*{[C] Внутрішнє тертя (в'язкість) (24) (с. 348–350)}
\begin{itemize}
\item Шари, що рухаються з різними швидкостями, обмінюються \textbf{імпульсом упорядкованого руху}, який переносить тепловий рух. Швидкий шар гальмується, повільний — прискорюється (англ. viscosity).
\item \textbf{Закон Ньютона}: $F = -\eta\dfrac{du}{dx}\Delta S$ (15.10), η — \textbf{коефіцієнт динамічної в'язкості}, Па·с. \textbf{Кінематична в'язкість} ν = η/ρ.
\item З кінетичної теорії: $\eta = \tfrac13\rho\langle v\rangle\langle l\rangle$ (15.14)
\item \spark{} \textbf{Парадокс} (с. 350): η \textbf{не залежить від тиску}, бо ρ ∝ p, а ⟨l⟩ ∝ 1/p. Імпульс переносить більше молекул, але кожна — на меншу відстань. Зменшення p у 500 разів змінює в'язкість повітря лише на 4 \%.
\item η зростає з T (∝ √T, на практиці швидше). Це виміряв Пулюй, с. 350.
\end{itemize}

\needspace{5\baselineskip}\section*{[D] Теплопровідність (25) (с. 352–354)}
\begin{itemize}
\item «Теплопровідність — це молекулярне перенесення енергії в суцільному середовищі, зумовлене градієнтом температури» (англ. thermal conductivity)
\item \textbf{Закон Фур'є} (англ. Fourier's law): $j_q = -\lambda\,\text{grad}\,T$ (15.16), λ — \textbf{коефіцієнт теплопровідності}, Вт/(м·К).
\item Кожна молекула несе свою середню енергію $\tfrac i2kT$. З кінетичної теорії:
$\lambda = \tfrac13\rho\langle v\rangle\langle l\rangle c_V$ (15.21), де $c_{V}$ — питома теплоємність.
\item Також \textbf{не залежить від p}, крім дуже низьких тисків. Легкі гази проводять теплоту краще (H₂ використовують для охолодження).
\end{itemize}

\needspace{5\baselineskip}\section*{[E] Зв'язок між коефіцієнтами (с. 358)}
\[\eta = D\rho,\qquad \lambda = \eta c_V,\qquad \lambda = D\rho c_V\]
{\footnotesize
\begin{longtable}{>{\raggedright\arraybackslash}p{0.204\linewidth}>{\raggedright\arraybackslash}p{0.191\linewidth}>{\raggedright\arraybackslash}p{0.089\linewidth}>{\raggedright\arraybackslash}p{0.191\linewidth}>{\raggedright\arraybackslash}p{0.140\linewidth}>{\raggedright\arraybackslash}p{0.089\linewidth}}
\toprule
\textbf{} & \textbf{що переноситься} & \textbf{закон} & \textbf{коефіцієнт} & \textbf{від p} & \textbf{від T} \\
\midrule
\endhead
дифузія & маса & Фіка & D = ⅓⟨v⟩⟨l⟩ & ∝ 1/p & зростає \\
в'язкість & імпульс & Ньютона & η = ⅓ρ⟨v⟩⟨l⟩ & не залежить & зростає \\
теплопровідність & енергія & Фур'є & λ = ⅓ρ⟨v⟩⟨l⟩$c_{V}$ & не залежить & зростає \\
\bottomrule
\end{longtable}
}

\needspace{5\baselineskip}\section*{Самоперевірка}
\begin{enumerate}
\item Чому D залежить від тиску, а η і λ — ні?
\item Що саме «несе» молекула в кожному з трьох явищ?
\item Звідки множник ⅓ (= 2 × ⅙)?
\end{enumerate}

\chapter*{3. Внутрішня енергія термодинамічної системи}\addcontentsline{toc}{chapter}{3. Внутрішня енергія термодинамічної системи}\markboth{3. Внутрішня енергія термодинамічної системи}{3. Внутрішня енергія термодинамічної системи}

\srcline{\book{} т. 1, \textbf{§ 16.1} «Основні поняття і означення» (с. 371–374), \textbf{§ 16.2} «Перший закон термодинаміки», частина «Внутрішня енергія» (с. 375–377). \ok{} Звірено зі сканом у складі блоку A (с. 375–377).}
Блок D, сторінка 1 з 4. Ця сторінка задає \textbf{мову термодинаміки}, якою користуються блоки D–F. Формула для ідеального газу U = (i/2)RT — на 20–21; тут лише посилання.

\needspace{5\baselineskip}\section*{План сторінки}
\begin{layout}
┌ [A] Термодинаміка: предмет; термодинамічна система; параметри стану   ┐
├ [B] Типи систем (ізольована/закрита/відкрита/адіабатна) + рівноважні процеси ┤
├ [C] Внутрішня енергія: означення, що входить / що НЕ входить              ┤
├ [D] Функція стану: ∮dU = 0 (малюнок двох шляхів A→B на p–V)              ┤
└ [E] Адитивність (16.1); калоричне рівняння стану; ідеальний газ → 20–21  ┘
\end{layout}

\needspace{5\baselineskip}\section*{[A] Основні поняття (с. 371–372)}
\begin{itemize}
\item \textbf{Термодинаміка} (англ. thermodynamics): «наука про найбільш загальні властивості макроскопічних фізичних систем, що знаходяться в стані термодинамічної рівноваги, та про процеси переходу між цими станами». Вона \textbf{аксіоматична}: спирається на закони, узагальнені з досвіду, а не на молекулярні моделі (с. 374).
\item \textbf{Термодинамічна система} (англ. thermodynamic system): макроскопічні тіла й поля, що можуть обмінюватися \textbf{енергією} (у формі роботи й теплоти) та \textbf{речовиною}.
\item \textbf{Термодинамічні параметри стану} (англ. state parameters): p, T, V, ρ, концентрація, напруженості полів… Бувають \textbf{внутрішні} (англ. internal) й \textbf{зовнішні} (англ. external), с. 372.
\end{itemize}

\needspace{5\baselineskip}\section*{[B] Системи і процеси (с. 372–373)}
{\small
\begin{longtable}{>{\raggedright\arraybackslash}p{0.451\linewidth}>{\raggedright\arraybackslash}p{0.251\linewidth}>{\raggedright\arraybackslash}p{0.251\linewidth}}
\toprule
\textbf{система} & \textbf{обмін речовиною} & \textbf{обмін енергією} \\
\midrule
\endhead
ізольована (замкнута) & ні & ні \\
закрита & ні & так \\
відкрита & так & так \\
теплоізольована (адіабатна) & — & без теплообміну \\
\bottomrule
\end{longtable}
}
\begin{itemize}
\item \textbf{Квазістатичний (квазірівноважний) процес} (англ. quasistatic process): неперервна послідовність рівноважних станів — наприклад, на поршень кладуть пісок по піщинці. Лише такі процеси можна зобразити на діаграмі p–V.
\item Необов'язково: \textbf{фаза} (англ. phase) і \textbf{компонента} (англ. component), с. 373.
\end{itemize}

\needspace{5\baselineskip}\section*{[C] Внутрішня енергія (с. 375)}
\begin{itemize}
\item «Енергію, що залежить від внутрішнього стану фізичної системи, називають \textbf{внутрішньою енергією}» (англ. internal energy)
\item \textbf{Входить}: кінетична енергія всіх видів руху частинок і енергія їхньої взаємодії. На практиці — \textbf{кінетична енергія теплового руху + потенціальна енергія взаємодії} (без внутрішньоатомної та ядерної енергії).
\item \textbf{Не входить}: кінетична енергія системи як цілого та її потенціальна енергія в зовнішніх полях — це «зовнішні» енергії.
\item Знак: dU > 0, якщо U зростає.
\end{itemize}

\needspace{5\baselineskip}\section*{[D] Функція стану (с. 375–376) — у рамку}
\begin{itemize}
\item U — \textbf{однозначна функція стану} (англ. state function): ΔU = U₂ − U₁ залежить лише від початкового й кінцевого станів, а не від шляху. Для циклу:
\[\oint dU = 0\]
\item \pen{} Намалюйте два різні шляхи A→B на діаграмі p–V з однаковим ΔU. Це підготовка до контрасту з Q і A (тема 4), які \textbf{залежать} від шляху.
\end{itemize}

\needspace{5\baselineskip}\section*{[E] Ще дві властивості (с. 376)}
\begin{itemize}
\item \textbf{Адитивність} (англ. additivity): $U = \sum U_i$ (16.1)
\item \textbf{Калоричне рівняння стану} (англ. caloric equation of state): U = U(V, T), або U(p, V), або U(p, T). Для рідин і газів U залежить від T (кінетична частина) \textbf{і від V} (взаємодія).
\item \textbf{Ідеальний газ}: взаємодії немає, тож U = U(T). Для одноатомного газу $U = \tfrac32RT$ (16.2); загальний випадок — на 20–21.
\end{itemize}

\needspace{5\baselineskip}\section*{Самоперевірка}
\begin{enumerate}
\item Чому кінетична енергія летючого м'яча не входить у його внутрішню енергію?
\item Що означає «функція стану» для замкненого циклу?
\item Чому U реального газу залежить від V, а ідеального — ні?
\end{enumerate}

\chapter*{4. Теплота, робота, теплоємність}\addcontentsline{toc}{chapter}{4. Теплота, робота, теплоємність}\markboth{4. Теплота, робота, теплоємність}{4. Теплота, робота, теплоємність}

\srcline{\book{} т. 1, \textbf{§ 16.2}, частини «Макроскопічна робота», «Кількість теплоти» (с. 377–381). \ok{} (16.5)–(16.9) звірено зі сканом.}
Блок D, сторінка 2 з 4. Головна думка: \textbf{Q і A — функції процесу}; U — функція стану (тема 3).

\needspace{5\baselineskip}\section*{План сторінки}
\begin{layout}
┌ [A] Робота: поршень (рис. 16.1), d'A = p dV (16.5), A = ∫p dV (16.6)     ┐
├ [B] Робота = площа під кривою; два шляхи A→B (рис. 16.2) → різна A       ┤
├ [C] Кількість теплоти: теплообмін, «мікроскопічна робота»; знаки          ┤
├ [D] Таблиця: робота vs теплота (макро/мікро; обидві — функції процесу)    ┤
└ [E] Теплоємність: C₀ = d'Q/dT (16.7), питома c (16.8), молярна C = Mc (16.9) ┘
\end{layout}

\needspace{5\baselineskip}\section*{[A] Макроскопічна робота (с. 377–378)}
\begin{itemize}
\item Газ під поршнем площею S квазістатично розширюється на dx: $d\text{'}A = F\,dx = pS\,dx = p\,dV$ (16.5) (англ. work)
\item \textbf{Знак} (угода книги): робота \textbf{газу} при розширенні \textbf{додатна}; робота зовнішніх сил над газом — від'ємна.
\item $A = \int_{V_1}^{V_2}p\,dV$ (16.6)
\item Означення книги (с. 378): «робота в термодинаміці — це \textbf{макроскопічна міра зміни внутрішньої енергії} системи у квазістатичних процесах»
\end{itemize}

\needspace{5\baselineskip}\section*{[B] Функція процесу (с. 378–379)}
\begin{itemize}
\item На діаграмі p–V \textbf{A = площа під кривою} процесу (рис. 16.2).
\item Два різні шляхи A→B мають різні площі, тобто різну роботу. Робота — \textbf{функція процесу, а не функція стану} (англ. path function).
\item Звідси позначення \textbf{d'A, а не dA}: це не повний диференціал.
\item У циклі робота може бути ≠ 0 — на цьому й працюють теплові машини.
\item \pen{} Головний малюнок: два шляхи AmB і AnB, площі під ними заштриховано по-різному.
\end{itemize}

\needspace{5\baselineskip}\section*{[C] Кількість теплоти (с. 380)}
\begin{itemize}
\item \textbf{Теплообмін} (англ. heat transfer): «Процес передачі енергії від більш нагрітого тіла до менш нагрітого \textbf{без виконання макроскопічної роботи}». Його міра — \textbf{кількість теплоти Q} (англ. heat).
\item Механізм — зіткнення молекул; іноді це називають «мікроскопічною роботою».
\item Знаки: теплота, \textbf{надана} системі, — +, відібрана — −. Одиниця: Дж.
\item Q теж функція процесу, тому пишуть \textbf{d'Q}.
\end{itemize}

\needspace{5\baselineskip}\section*{[D] Робота vs теплота (с. 380)}
{\small
\begin{longtable}{>{\raggedright\arraybackslash}p{0.179\linewidth}>{\raggedright\arraybackslash}p{0.366\linewidth}>{\raggedright\arraybackslash}p{0.407\linewidth}}
\toprule
\textbf{} & \textbf{робота A} & \textbf{теплота Q} \\
\midrule
\endhead
форма передачі енергії & \textbf{макрофізична}: енергія впорядкованого руху & \textbf{мікрофізична}: сукупність молекулярних процесів \\
пов'язана зі зміною & зовнішніх параметрів (V) & жодних; лише теплообмін \\
функція & процесу (d'A) & процесу (d'Q) \\
спільне & обидві — міри \textbf{передачі енергії}, у Дж &  \\
\bottomrule
\end{longtable}
}

\needspace{5\baselineskip}\section*{[E] Теплоємність (с. 380–381)}
\begin{itemize}
\item $C_0 = \dfrac{d\text{'}Q}{dT}$ (16.7) — теплота, потрібна для нагрівання на 1 К (англ. heat capacity).
\item \textbf{Питома теплоємність} (англ. specific heat capacity) c — на одиницю маси: $C_0 = mc$ (16.8)
\item \textbf{Молярна теплоємність} (англ. molar heat capacity) C — на моль: $C = Mc$ (16.9)
\item Позначення: $C_{p}$ (ізобарна), $C_{V}$ (ізохорна), $C_{T}$ (ізотермічна), $C_{\text{ад}}$ (адіабатна). Оскільки d'Q залежить від процесу, \textbf{теплоємність теж залежить від процесу}. Значення для ідеального газу — на 05.
\end{itemize}

\needspace{5\baselineskip}\section*{Самоперевірка}
\begin{enumerate}
\item Чому пишуть d'A, а не dA?
\item Чому робота за цикл не обов'язково нульова, а ΔU — обов'язково?
\item У чому фізична різниця між теплотою і роботою?
\end{enumerate}

\chapter*{6. Перший закон термодинаміки}\addcontentsline{toc}{chapter}{6. Перший закон термодинаміки}\markboth{6. Перший закон термодинаміки}{6. Перший закон термодинаміки}

\srcline{\book{} т. 1, \textbf{§ 16.2} «Перший закон термодинаміки» (с. 374, 381–383). \ok{} (16.10)–(16.17) звірено зі сканом.}
Блок D, сторінка 3 з 4. Поєднує теми 3 і 4: два способи змінити U.

\needspace{5\baselineskip}\section*{План сторінки}
\begin{layout}
┌ [A] Звідки: механічна енергія не зберігається при терті → теплота = енергія ┐
├ [B] Формулювання + (16.10) dU = d'Q − d'A  ⚠️ знаки!                        ┤
├ [C] Зручна форма (16.11) d'Q = dU + d'A; інтегральна (16.12); цикл (16.13)  ┤
├ [D] Вічний двигун I роду неможливий                                         ┤
└ [E] Для ідеального газу: d'Q = C_V dT + p dV (16.17) → місток до 05, 07–10   ┘
\end{layout}

\needspace{5\baselineskip}\section*{[A] Звідки закон (с. 374)}
\begin{itemize}
\item Механічна енергія зберігається лише без \textbf{дисипативних сил}. При терті вона «зникає», але з'являється теплота. Урахування внутрішньої енергії теплового руху відновлює закон збереження.
\item «Перший закон термодинаміки — це \textbf{узагальнений закон збереження й перетворення енергії} для тих макроскопічних систем, для яких суттєву роль відіграє тепловий рух» (англ. first law of thermodynamics)
\end{itemize}

\needspace{5\baselineskip}\section*{[B] Формулювання (с. 381) — у рамку}
\begin{itemize}
\item «у теплових процесах зміна внутрішньої енергії відбувається за рахунок приросту кількості теплоти і виконаної роботи»
\[dU = d\text{'}Q - d\text{'}A\qquad(16.10)\]
\item \warn{} \textbf{Неточність у книзі.} Під (16.10) сказано «d'A — робота, виконана \textbf{над} системою», але зі знаком мінус і власною угодою книги на с. 378 (d'A = p dV > 0 при розширенні) d'A має бути роботою, виконаною \textbf{самою} системою. Запишіть: \textbf{d'A — робота системи проти зовнішніх сил.} Адіабатний випадок (16.20), d'A = −dU («система виконує роботу внаслідок зменшення внутрішньої енергії»), підтверджує це прочитання.
\end{itemize}

\needspace{5\baselineskip}\section*{[C] Інші форми (с. 382)}
\begin{itemize}
\item Форма через теплоту, зручна для теплотехніки: $d\text{'}Q = dU + d\text{'}A$ (16.11). Словами: надана теплота \textbf{йде на} збільшення U і виконання роботи.
\item Інтегральна форма: $Q = (U_2 - U_1) + A$ (16.12)
\item Цикл (U₂ = U₁): $Q = A$ (16.13) — уся отримана теплота перетворюється на роботу.
\item Якщо d'Q = 0, процес \textbf{адіабатний} (тема 10).
\end{itemize}

\needspace{5\baselineskip}\section*{[D] Вічний двигун першого роду (с. 382)}
\begin{itemize}
\item «неможливо побудувати машину, що діє періодично, яка виконувала б роботу без одержання енергії ззовні або виконувала б більшу роботу, ніж одержана ззовні енергія» (англ. perpetual motion machine of the first kind)
\end{itemize}

\needspace{5\baselineskip}\section*{[E] Ідеальний газ (с. 382–383)}
\begin{itemize}
\item U = U(T, V) (16.14), але для ідеального газу (∂U/∂V)\_T = 0, тож $dU = C_VdT$ (16.16), де $C_V = (\partial U/\partial T)_V$.
\item Один моль:
\[d\text{'}Q = C_VdT + p\,dV\qquad(16.17)\]
Із цього рядка випливає все в блоках D–E: рівняння Майєра (05), ізопроцеси (07–09) і адіабата (10).
\end{itemize}

\needspace{5\baselineskip}\section*{Самоперевірка}
\begin{enumerate}
\item Сформулюйте перший закон словами, без формул.
\item У циклі ΔU = 0; куди дівається теплота?
\item Чому вічний двигун першого роду неможливий?
\end{enumerate}

\chapter*{5. Теплоємність ідеальних газів}\addcontentsline{toc}{chapter}{5. Теплоємність ідеальних газів}\markboth{5. Теплоємність ідеальних газів}{5. Теплоємність ідеальних газів}

\srcline{\book{} т. 1, \textbf{§ 16.2} (рівняння Майєра, с. 382–383), \textbf{§ 16.5} «Теплоємність ідеального газу» (с. 390–397). \ok{} (16.16)–(16.19), (16.46) і таблицю $C_{V}$/$C_{p}$ звірено зі сканом.}
Блок D, сторінка 4 з 4. Саму таблицю $C_{V}$ подано на 20–21. Тут — \textbf{$C_{V}$ vs $C_{p}$}, \textbf{γ} і \textbf{де класична теорія не працює}.

\needspace{5\baselineskip}\section*{План сторінки}
\begin{layout}
┌ [A] C_V = dU/dT; C_p — більше; рівняння Майєра (16.19) у рамці        ┐
├ [B] Класична теорія: C_V = (i/2)R, C_p = ((i+2)/2)R, γ = (i+2)/i — таблиця ┤
├ [C] Графік C_V(T) для двохатомного газу (рис. 16.7): сходинки 3/2 → 5/2 → 7/2 ┤
└ [D] Чому: квантування; характеристичні температури T_оберт, T_колив      ┘
\end{layout}

\needspace{5\baselineskip}\section*{[A] Рівняння Майєра (с. 382–383) — у рамку}
\begin{itemize}
\item З першого закону для одного моля, $d\text{'}Q = C_VdT + p\,dV$ (16.17):
\begin{itemize}
\item V = const: dV = 0, тож $C_V = \left(\dfrac{\partial U}{\partial T}\right)_V$ (16.16)
\item p = const: p dV = R dT, тож $d\text{'}Q = C_VdT + R\,dT$ (16.18)
\[C_p = C_V + R\qquad(16.19)\]
— \textbf{рівняння Майєра} (англ. Mayer's relation).
\end{itemize}
\item Зміст: при сталому p частина теплоти йде на \textbf{роботу розширення}, а не лише на U. Тому $C_{p}$ > $C_{V}$, і саме на R (зміст R — на 01).
\end{itemize}

\needspace{5\baselineskip}\section*{[B] Класична теорія (с. 390–392)}
\begin{itemize}
\item З U = (i/2)RT: $C_V = \dfrac i2R$, $C_p = \dfrac{i+2}2R$, $\gamma = \dfrac{C_p}{C_V} = \dfrac{i+2}i$
\end{itemize}
{\footnotesize
\begin{longtable}{>{\raggedright\arraybackslash}p{0.460\linewidth}>{\raggedright\arraybackslash}p{0.077\linewidth}>{\raggedright\arraybackslash}p{0.096\linewidth}>{\raggedright\arraybackslash}p{0.096\linewidth}>{\raggedright\arraybackslash}p{0.192\linewidth}}
\toprule
\textbf{газ} & \textbf{i} & \textbf{$C_{V}$} & \textbf{$C_{p}$} & \textbf{γ} \\
\midrule
\endhead
одноатомний & 3 & 3/2 R & 5/2 R & 5/3 ≈ 1,67 \\
двохатомний (жорсткий) & 5 & 5/2 R & 7/2 R & 7/5 = 1,4 \\
багатоатомний (жорсткий) & 6 & 3R & 4R & 4/3 ≈ 1,33 \\
\bottomrule
\end{longtable}
}
\begin{itemize}
\item γ — \textbf{показник адіабати} (англ. adiabatic index), знадобиться в темі 10.
\item Загальна формула з коливаннями: $U = (3 + r + 2s)\dfrac{RT}2$ (16.46).
\end{itemize}

\needspace{5\baselineskip}\section*{[C] Невдачі класичної теорії (с. 392–394)}
\begin{enumerate}
\item Складні молекули: для C₂H₅OH жорстка модель дає 3R, пружна — 24R, а дослід — близько 8R. \textbf{Обидві моделі не працюють.}
\item \textbf{Залежність від температури} не пояснюється. Для двохатомного газу, рис. 16.7:
\begin{itemize}
\item низькі T: $C_{V}$ ≈ 3/2 R, ніби є лише поступальний рух (H₂ нижче 60 К поводиться як одноатомний);
\item кімнатні T: 5/2 R («увімкнулося» обертання);
\item тисячі К: прямує до 7/2 R (вмикаються коливання, але молекули раніше дисоціюють).
\end{itemize}
\item \pen{} Намалюйте «сходинки» $C_{V}$(T) з переходами при $T_{\text{оберт}}$ і $T_{\text{колив}.}$
\end{enumerate}

\needspace{5\baselineskip}\section*{[D] Квантове пояснення (с. 394–397), коротко}
\begin{itemize}
\item Енергії обертання й коливань \textbf{квантуються}: $\varepsilon_{\text{пост}} \ll \varepsilon_{\text{оберт}} \ll \varepsilon_{\text{колив}}$, звідси \textbf{характеристичні температури} (англ. characteristic temperatures) $T_{\text{пост}} \ll T_{\text{оберт}} \ll T_{\text{колив}}$.
\item Нижче $T_{\text{оберт}}$ чи $T_{\text{колив}}$ енергії kT замало, щоб збудити відповідний рух, тож ці ступені вільності «заморожені».
\item Значення: $T_{\text{оберт}}$ для H₂ — 85 К, для N₂ — 2,85 К; $T_{\text{колив}}$ для N₂ — 3340 К, для O₂ — 2240 К.
\end{itemize}

\needspace{5\baselineskip}\section*{Самоперевірка}
\begin{enumerate}
\item Чому $C_{p}$ − $C_{V}$ дорівнює саме R?
\item Чому двохатомний H₂ при 60 К поводиться як одноатомний?
\item Яке γ брати для повітря і чому?
\end{enumerate}

\chapter*{7. Ізотермічний процес · 8. Ізохоричний процес · 9. Ізобарний процес}\addcontentsline{toc}{chapter}{7. Ізотермічний процес · 8. Ізохоричний процес · 9. Ізобарний процес}\markboth{7. Ізотермічний процес}{7. Ізотермічний процес}

\srcline{\book{} т. 1, \textbf{§ 16.4} «Робота ідеального газу при ізопроцесах. Політропний процес» (с. 386–390); самі закони — у \textbf{§ 14.7} (с. 313–316, див. 01). \ok{} (16.26)–(16.33) звірено зі сканом.}
Блок E, сторінка 1 з 2. Правопис: у програмі «ізохоричний», у книзі — «\textbf{ізохорний}».
Користуйтеся разом із підсумковою таблицею процесів — найкориснішою окремою сторінкою для МКР.

\needspace{5\baselineskip}\section*{План сторінки (по блоку на процес, у кожному ті самі 5 рядків)}
\begin{layout}
┌ [Ізохорний V = const]  ┬ [Ізобарний p = const] ┬ [Ізотермічний T = const] ┐
│ закон: p/T = const     │ V/T = const (Г-Л)     │ pV = const (Б-М)          │
│ графік p–V: вертикаль  │ горизонталь + площа   │ гіпербола + площа         │
│ A = 0                  │ A = p(V₂−V₁)=νRΔT     │ A = νRT ln(V₂/V₁)          │
│ ΔU = νC_VΔT            │ ΔU = νC_VΔT           │ ΔU = 0                     │
│ Q = ΔU                 │ Q = νC_pΔT            │ Q = A                      │
└────────────────────────┴───────────────────────┴────────────────────────────┘
          + низ: політропний процес pVⁿ = const об'єднує всі (16.40)
\end{layout}
(ν = m/M)

\needspace{5\baselineskip}\section*{Ізохорний процес, V = const (с. 386–387)}
\begin{itemize}
\item Приклад: нагрівання газу в жорсткій герметичній посудині (англ. isochoric process).
\item $d\text{'}A_V = p\,dV = 0$, $A_V = 0$ (16.26)
\item З першого закону: $d\text{'}Q = dU = \dfrac mMC_VdT$ (16.27) і $\Delta U = \dfrac mMC_V(T_2 - T_1)$ (16.28).
\item \textbf{Уся теплота йде на внутрішню енергію.}
\end{itemize}

\needspace{5\baselineskip}\section*{Ізобарний процес, p = const (с. 387)}
\begin{itemize}
\item Приклад: газ під вільно рухомим поршнем або нагрівання води у відкритій посудині (англ. isobaric process).
\item $A_p = p(V_2 - V_1) = \dfrac mMR(T_2 - T_1)$ (16.29). На діаграмі p–V це \textbf{прямокутник} (рис. 16.4).
\item $Q_p = \dfrac mM(C_V + R)(T_2 - T_1) = \dfrac mMC_p(T_2 - T_1)$ (16.30)
\item Теплота йде \textbf{і на ΔU, і на роботу} (Майєр).
\end{itemize}

\needspace{5\baselineskip}\section*{Ізотермічний процес, T = const (с. 387–388)}
\begin{itemize}
\item Потрібен ідеальний тепловий контакт із \textbf{термостатом} (тілом сталої температури) і дуже повільний перебіг (англ. isothermal process). Усі фазові переходи першого роду (наприклад, кипіння) ізотермічні.
\item dU = 0, тож $d\text{'}Q = d\text{'}A$ (16.31): \textbf{уся теплота перетворюється на роботу}.
\item \[A_T = \frac mMRT\ln\frac{V_2}{V_1}\qquad(16.32)\]
\item Також: $Q_T = A_T = \dfrac mMRT\ln\dfrac{p_1}{p_2} = p_1V_1\ln\dfrac{V_2}{V_1}$ (16.33)
\item \pen{} Площа під гіперболою. Пам'ятайте, звідки логарифм: ∫dV/V.
\end{itemize}

\needspace{5\baselineskip}\section*{Низ: політропний процес (с. 388–389), необов'язково, але красиво}
\begin{itemize}
\item Процес зі \textbf{сталою теплоємністю C} (англ. polytropic process): $pV^n = \text{const}$ (16.40), де $n = \dfrac{C_p - C}{C_V - C}$ (16.37).
\end{itemize}
{\small
\begin{longtable}{>{\raggedright\arraybackslash}p{0.571\linewidth}>{\raggedright\arraybackslash}p{0.190\linewidth}>{\raggedright\arraybackslash}p{0.190\linewidth}}
\toprule
\textbf{процес} & \textbf{C} & \textbf{n} \\
\midrule
\endhead
ізобарний & $C_{p}$ & 0 \\
ізотермічний & ∞ & 1 \\
адіабатний & 0 & γ \\
ізохорний & $C_{V}$ & ±∞ \\
\bottomrule
\end{longtable}
}
\begin{itemize}
\item Усі чотири процеси — окремі випадки одного рівняння.
\end{itemize}

\needspace{5\baselineskip}\section*{Самоперевірка}
\begin{enumerate}
\item У якому процесі вся теплота йде на роботу, а в якому — вся на U?
\item Чому $C_{T}$ = ∞? (Теплоту надають, а ΔT = 0.)
\item Намалюйте всі три процеси на одній діаграмі p–V з однієї початкової точки.
\end{enumerate}

\chapter*{10. Адіабатичний процес}\addcontentsline{toc}{chapter}{10. Адіабатичний процес}\markboth{10. Адіабатичний процес}{10. Адіабатичний процес}

\srcline{\book{} т. 1, \textbf{§ 16.3} «Адіабатний процес. Рівняння Пуассона» (с. 383–386). \ok{} (16.20)–(16.25) звірено зі сканом.}
Блок E, сторінка 2 з 2. Термінологія: у програмі «адіабатичний», у книзі — «\textbf{адіабатний}». Пишіть за книгою.

\needspace{5\baselineskip}\section*{План сторінки}
\begin{layout}
┌ [A] Означення + умови (теплоізоляція, «достатня повільність»)        ┐
├ [B] Перший закон: dU + d'A = 0 (16.20) → pdV = −C_V dT (16.21)        ┤
├ [C] Виведення → рівняння Пуассона TV^(γ−1) = const, pV^γ = const      ┤
├ [D] ГРАФІК: адіабата крутіша за ізотерму (рис. 16.3) + чому           ┤
└ [E] Робота (16.24)–(16.25); приклади з життя                          ┘
\end{layout}

\needspace{5\baselineskip}\section*{[A] Означення (с. 383–384)}
\begin{itemize}
\item \textbf{Адіабатний процес} (англ. adiabatic process): «що відбувається \textbf{без теплообміну} з навколишнім середовищем» (d'Q = 0).
\item Дві умови: \textbf{теплоізоляція} (нетеплопровідні стінки) і \textbf{достатня повільність}, щоб процес лишався квазістатичним. На практиці процес роблять швидким, щоб теплота не встигала піти, але повільнішим за швидкість звуку.
\end{itemize}

\needspace{5\baselineskip}\section*{[B] Перший закон (с. 384)}
\begin{itemize}
\item $dU + d\text{'}A = 0$ (16.20) ⟹ $d\text{'}A = -dU$. «система виконує роботу внаслідок \textbf{зменшення внутрішньої енергії}»
\item Один моль: $p\,dV = -C_VdT$ (16.21). \textbf{Розширення охолоджує} газ, стискання — нагріває.
\end{itemize}

\needspace{5\baselineskip}\section*{[C] Рівняння Пуассона (с. 384) — у рамку}
\begin{itemize}
\item Підставляємо p = RT/V: $\dfrac{dT}T = -\dfrac R{C_V}\dfrac{dV}V$ (16.22). Інтегруємо, врахувавши R/$C_{V}$ = γ − 1:
\[TV^{\gamma-1} = \text{const}\ (16.23),\qquad pV^\gamma = \text{const},\quad pT^{\frac\gamma{1-\gamma}} = \text{const}\ (16.23a)\]
\item \textbf{γ = $C_{p}$/$C_{V}$} — \textbf{показник адіабати}, або \textbf{коефіцієнт Пуассона} (англ. adiabatic index; значення — на 05).
\end{itemize}

\needspace{5\baselineskip}\section*{[D] Адіабата vs ізотерма (с. 384–385)}
\begin{itemize}
\item Ізотерма: p ∝ 1/V. Адіабата: p ∝ 1/$V^{γ}$, γ > 1, отже \textbf{адіабата крутіша}.
\item Фізично: при адіабатному розширенні p спадає і тому, що \textbf{V зростає, і тому, що T спадає}.
\item \pen{} Головний малюнок, рис. 16.3: ізотерма й адіабата через одну точку, адіабата крутіша, заштрихована площа роботи V₁→V₂.
\end{itemize}

\needspace{5\baselineskip}\section*{[E] Робота (с. 385)}
\begin{itemize}
\item $A = -\dfrac mMC_V\int_{T_1}^{T_2}dT = \dfrac mMC_V(T_1 - T_2)$ (16.24)
\item \[A = \frac mM\frac{RT_1}{\gamma - 1}\left[1 - \left(\frac{V_1}{V_2}\right)^{\gamma-1}\right]\qquad(16.25)\]
\item Приклади з життя:
\begin{itemize}
\item повітря, що піднімається в тропосфері, адіабатно розширюється й охолоджується — тому вгорі холодніше;
\item двигуни внутрішнього згоряння;
\item одержання низьких температур.
\end{itemize}
\end{itemize}

\needspace{5\baselineskip}\section*{Самоперевірка}
\begin{enumerate}
\item Чому газ охолоджується при адіабатному розширенні, хоча теплоту не відбирають?
\item Чому адіабата крутіша за ізотерму? Дві причини одним реченням.
\item Повітря: γ = 1,4. Як зміниться T, якщо об'єм зменшити вдвічі? (×2\textasciicircum{}0,4 ≈ 1,32.)
\end{enumerate}

\chapter*{11. Колові процеси}\addcontentsline{toc}{chapter}{11. Колові процеси}\markboth{11. Колові процеси}{11. Колові процеси}

\srcline{\book{} т. 1, \textbf{§ 16.6} «Оборотні та необоротні процеси. Колові процеси (цикли)» (с. 397–399). \ok{} Звірено зі сканом (с. 397–399; (16.13) — на с. 382).}
Блок F, сторінка 1 з 5.

\needspace{5\baselineskip}\section*{План сторінки}
\begin{layout}
┌ [A] Оборотні vs необоротні процеси (таблиця + приклади)          ┐
├ [B] Цикл: означення; ΔU = 0 → Q = A (16.13)                       ┤
├ [C] ГРАФІК: цикл 1a2b1 на p–V (рис. 16.8): A = площа всередині     ┤
└ [D] Прямий (теплова машина) vs обернений (холодильна) цикл        ┘
\end{layout}

\needspace{5\baselineskip}\section*{[A] Оборотні та необоротні процеси (с. 397)}
\begin{itemize}
\item \textbf{Оборотний процес} (англ. reversible process): «такий процес, при якому система, проведена в зворотному напрямі, повертається в початковий стан \textbf{через ті самі проміжні стани}, що й у прямому процесі, а стан тіл поза системою залишається незмінним». Усі квазістатичні процеси оборотні.
\item \textbf{Необоротний процес} (англ. irreversible process): такий, що його не можна провести у зворотному напрямі без змін у системі чи навколишніх тілах. \textbf{Усі реальні процеси необоротні.}
\item Типові необоротні процеси: теплообмін за скінченної різниці температур, дифузія, розширення у вакуум, тертя. Головна ознака необоротності — \textbf{теплота дисипації} (с. 398).
\item Навіщо вивчати ідеальні оборотні процеси: вони задають межу, до якої наближаються реальні.
\end{itemize}

\needspace{5\baselineskip}\section*{[B] Цикл (с. 398)}
\begin{itemize}
\item \textbf{Цикл (коловий процес)} (англ. cycle): «термодинамічний процес, у результаті якого система повертається в початковий стан»
\item U — функція стану, тож за цикл ΔU = 0, і $Q = A$ (16.13): отримана теплота перетворюється на роботу (06).
\item Навіщо цикли: одне розширення дає обмежену роботу. Щоб повторити, газ треба стиснути назад \textbf{іншим шляхом}, при меншому тиску, інакше сумарна робота дорівнюватиме 0.
\end{itemize}

\needspace{5\baselineskip}\section*{[C] Робота циклу (с. 398–399)}
\begin{itemize}
\item Робота циклу дорівнює \textbf{площі, обмеженій циклом} на діаграмі p–V (рис. 16.8): робота розширення мінус робота стискання.
\item \pen{} Головний малюнок: цикл 1a2b1, горизонтальне штрихування під верхньою гілкою (розширення), вертикальне — під нижньою (стискання), різниця заштрихована.
\end{itemize}

\needspace{5\baselineskip}\section*{[D] Прямий і обернений цикл (с. 399)}
{\small
\begin{longtable}{>{\raggedright\arraybackslash}p{0.093\linewidth}>{\raggedright\arraybackslash}p{0.464\linewidth}>{\raggedright\arraybackslash}p{0.395\linewidth}}
\toprule
\textbf{} & \textbf{\textbf{прямий}} & \textbf{\textbf{обернений}} \\
\midrule
\endhead
машина & теплова машина (двигун) & холодильна машина \\
обхід на p–V & за годинниковою стрілкою & проти годинникової стрілки \\
робота & $A_{\text{розш}}$ > $A_{\text{стиск}}$, A > 0 & $A_{\text{стиск}}$ > $A_{\text{розш}}$ \\
теплота & частина теплоти → робота, решта → холодильник (англ. cooler) & робота переносить теплоту від холодного до гарячого \\
\bottomrule
\end{longtable}
}
\begin{itemize}
\item \textbf{Оборотний цикл} (англ. reversible cycle): усі його частини оборотні. \textbf{Необоротний}: хоча б одна частина необоротна.
\end{itemize}

\needspace{5\baselineskip}\section*{Самоперевірка}
\begin{enumerate}
\item Чому не можна стиснути газ назад тією самою кривою?
\item У якому напрямі обходить цикл теплова машина на діаграмі p–V?
\item Чому кожен реальний процес необоротний?
\end{enumerate}

\chapter*{12. Цикл Карно (необернений) та його ккд · 13. Цикл Карно (обернений) та його ккд}\addcontentsline{toc}{chapter}{12. Цикл Карно (необернений) та його ккд · 13. Цикл Карно (обернений) та його ккд}\markboth{12. Цикл Карно (необернений) та його ккд}{12. Цикл Карно (необернений) та його ккд}

\srcline{\book{} т. 1, \textbf{§ 16.6} «… Цикл Карно та його коефіцієнт корисної дії» (с. 399–401), \textbf{§ 16.7} «Теплові машини» (с. 402–403). \ok{} (16.53)–(16.68) звірено зі сканом.}
Блок F, сторінка 2 з 5. Термінологія: у програмі «\textbf{необернений}» означає \textbf{прямий} цикл. Книга цього слова не вживає, у ній — \textbf{прямий} / \textbf{обернений} (англ. direct / reverse). Запишіть один раз «прямий (необернений)».

\needspace{5\baselineskip}\section*{План сторінки}
\begin{layout}
┌ [A] Що таке цикл Карно: 2 ізотерми + 2 адіабати; ідеальна машина     ┐
├ [B] ГРАФІК 1→2→3→4→1 (рис. 16.9) з T₁, T₂, Q₁, Q₂ + таблиця 4 стадій ┤
├ [C] ККД прямого циклу: η = 1 − Q₂/Q₁ (16.53) = 1 − T₂/T₁ (16.64) — рамка ┤
├ [D] Обернений цикл: холодильна машина ε (16.68), тепловий насос φ (16.67) ┤
└ [E] Схеми трьох машин (рис. 16.10–16.12) — стрілки потоків енергії     ┘
\end{layout}

\needspace{5\baselineskip}\section*{[A] Цикл Карно (с. 399)}
\begin{itemize}
\item Запитання Карно (1824): чи є межа вдосконалення теплових машин і чи важливе робоче тіло?
\item «Він складається з \textbf{двох адіабатних та двох ізотермічних} процесів» (англ. Carnot cycle). Цикл \textbf{оборотний} — це цикл \textbf{ідеальної} теплової машини, без втрат на теплопровідність, випромінювання чи тертя.
\item Через нього також виходять на кількісне формулювання \textbf{другого закону} (16).
\end{itemize}

\needspace{5\baselineskip}\section*{[B] Прямий цикл з ідеальним газом (с. 400–401), 1 моль, за годинниковою стрілкою}
{\footnotesize
\begin{longtable}{>{\raggedright\arraybackslash}p{0.069\linewidth}>{\raggedright\arraybackslash}p{0.184\linewidth}>{\raggedright\arraybackslash}p{0.214\linewidth}>{\raggedright\arraybackslash}p{0.345\linewidth}>{\raggedright\arraybackslash}p{0.108\linewidth}}
\toprule
\textbf{стадія} & \textbf{процес} & \textbf{контакт} & \textbf{робота} & \textbf{теплота} \\
\midrule
\endhead
1→2 & ізотермічне розширення, T₁ & \textbf{нагрівник} (англ. heater) & $A_{12} = RT_1\ln\frac{V_2}{V_1} = Q_1$ (16.55) & Q₁ > 0 отримано \\
2→3 & адіабатне розширення & ізольовано & $A_{23} = C_V(T_1 - T_2)$ (16.56) & 0 \\
3→4 & ізотермічне стискання, T₂ & \textbf{холодильник} (англ. cooler) & $A_{34} = -RT_2\ln\frac{V_3}{V_4} = -Q_2$ (16.57) & Q₂ віддано \\
4→1 & адіабатне стискання & ізольовано & $A_{41} = C_V(T_2 - T_1)$ (16.58) & 0 \\
\bottomrule
\end{longtable}
}
\begin{itemize}
\item Роботи на адіабатах взаємно знищуються. З рівнянь адіабат: $\dfrac{V_2}{V_1} = \dfrac{V_3}{V_4}$ (16.63) — \textbf{умова замкненості циклу}.
\item \pen{} Головний малюнок, рис. 16.9: дві ізотерми (T₁ вище, T₂ нижче) і дві крутіші адіабати, стрілки 1→2→3→4, Q₁ входить, Q₂ виходить.
\end{itemize}

\needspace{5\baselineskip}\section*{[C] ККД (с. 400–401) — у рамку}
\[\eta = \frac A{Q_1} = \frac{Q_1 - Q_2}{Q_1} = 1 - \frac{Q_2}{Q_1}\ (16.53)\qquad \eta = \frac{T_1 - T_2}{T_1} = 1 - \frac{T_2}{T_1}\ (16.64)\]
\begin{itemize}
\item ККД оборотного циклу Карно залежить \textbf{лише від температур нагрівника й холодильника} (англ. efficiency).
\item \textbf{Завжди η < 1}: для η = 1 треба T₁ → ∞ або T₂ → 0, що неможливо.
\item Якщо ΔT → 0, то η → 0, тож теплова машина \textbf{обов'язково має два термостати}.
\end{itemize}

\needspace{5\baselineskip}\section*{[D] Обернений цикл (с. 399, 402–403)}
\begin{itemize}
\item Обхід проти годинникової стрілки: робота \textbf{зовнішніх сил} переносить теплоту від \textbf{холоднішого} тіла до \textbf{гарячішого}.
\item \textbf{Холодильна машина} (англ. refrigerator): $\varepsilon = \dfrac{Q_2}{A}$ (16.68) — \textbf{холодильний коефіцієнт} (англ. coefficient of performance): теплота, відібрана від холодного тіла, на одиницю роботи.
\begin{itemize}
\item Для оберненого циклу Карно з (16.64) випливає $\varepsilon = \dfrac{T_2}{T_1 - T_2}$. \warn{} Це виведено тут; у книзі цієї формули немає.
\end{itemize}
\item \textbf{Тепловий насос} (англ. heat pump): $\varphi = \dfrac{Q_1}{A} = \dfrac{T_1}{T_1 - T_2}$ (16.67) — \textbf{коефіцієнт перетворення}: теплота, передана теплому тілу, на одиницю роботи. Може бути > 1.
\end{itemize}

\needspace{5\baselineskip}\section*{[E] Три типи теплових машин (с. 402)}
\begin{itemize}
\item \pen{} Три маленькі блок-схеми (рис. 16.10–16.12):
\begin{itemize}
\item двигун: нагрівник → Q₁ → робоче тіло → A назовні, і Q₂ → холодильник;
\item холодильна машина: A ззовні, холодне тіло → Q₂ → робоче тіло → Q₁ → довкілля;
\item тепловий насос: ті самі стрілки, але мета — Q₁ у теплу кімнату.
\end{itemize}
\end{itemize}

\needspace{5\baselineskip}\section*{Самоперевірка}
\begin{enumerate}
\item Чому ККД циклу Карно не залежить від газу?
\item Куди дівається робота на двох адіабатах?
\item Чому φ теплового насоса може бути більшим за 1, не порушуючи закону збереження енергії?
\end{enumerate}

\chapter*{16. Другий закон термодинаміки}\addcontentsline{toc}{chapter}{16. Другий закон термодинаміки}\markboth{16. Другий закон термодинаміки}{16. Другий закон термодинаміки}

\srcline{\book{} т. 1, \textbf{§ 16.8} «Другий закон термодинаміки. Теорема Карно. Поняття про ентропію. Характеристичні функції» (с. 403–409). \ok{} (16.69)–(16.74) звірено зі сканом.}
Блок F, сторінка 3 з 5. Сама ентропія (той самий параграф, від с. 406) — на 14.

\needspace{5\baselineskip}\section*{План сторінки}
\begin{layout}
┌ [A] Що визначає закон: НАПРЯМ процесів + межу перетворення Q → A      ┐
├ [B] Три формулювання: Клаузіус · Томсон · Оствальд (у рамках)          ┤
├ [C] Теорема Карно: два твердження; η ≤ η₀ (16.69)                     ┤
├ [D] Кількісно: нерівність Клаузіуса Q₁/T₁ + Q₂/T₂ ≤ 0 (16.71); зведена теплота ┤
└ [E] Узагальнення на довільний цикл ∮d'Q/T ≤ 0 (16.74) + рис. 16.13     ┘
\end{layout}

\needspace{5\baselineskip}\section*{[A] Що визначає (с. 403)}
\begin{itemize}
\item Перший закон дозволяє будь-який процес, що зберігає енергію. \textbf{Другий закон} (англ. second law of thermodynamics) «визначає \textbf{напрям} термодинамічних процесів та встановлює максимально можливі границі перетворення теплоти в роботу». Це узагальнення досвіду.
\end{itemize}

\needspace{5\baselineskip}\section*{[B] Формулювання (с. 403–404)}
\begin{itemize}
\item \textbf{Клаузіус (1850)}: «теплота не може \textbf{сама собою} переходити від менш нагрітого тіла до більш нагрітого»
\item \textbf{Томсон (Кельвін, 1851)}: «не можна побудувати машину, що діє періодично, яка б неперервно перетворювала теплоту в роботу \textbf{тільки внаслідок охолодження одного тіла}, без того, щоб у навколишніх тілах не відбулося одночасно певних залишкових змін»
\item \textbf{Оствальд}: «\textbf{вічний двигун другого роду} неможливий» (англ. perpetual motion machine of the second kind). Така машина працювала б, охолоджуючи океан; закону збереження енергії вона не порушує, але її забороняє другий закон.
\item \pen{} Порівняйте з вічним двигуном \textbf{першого} роду (06), який заборонено законом збереження енергії.
\end{itemize}

\needspace{5\baselineskip}\section*{[C] Теорема Карно (с. 404)}
\begin{enumerate}
\item ККД \textbf{оборотного} циклу Карно не залежить від робочого тіла й конструкції; він дорівнює ККД циклу з ідеальним газом при тих самих температурах (англ. Carnot's theorem).
\item ККД \textbf{необоротного} циклу завжди менший.
\item $\eta \le \eta_0$ (16.69): рівність — для оборотного циклу, < — для необоротного.
\end{enumerate}

\needspace{5\baselineskip}\section*{[D] Нерівність Клаузіуса (с. 404)}
\begin{itemize}
\item Підставляємо $\eta = \dfrac{Q_1 - |Q_2|}{Q_1}$ і $\eta_0 = \dfrac{T_1 - T_2}{T_1}$ (16.70). При Q₂ < 0:
\[\frac{Q_1}{T_1} + \frac{Q_2}{T_2} \le 0\qquad(16.71)\]
(англ. Clausius inequality)
\item \textbf{Зведена теплота} (англ. reduced heat) — теплота, якою система обмінюється з резервуаром, поділена на абсолютну температуру цього резервуара.
\item Словами: сума зведених теплот в оборотному циклі Карно дорівнює 0, у необоротному — менша за 0.
\end{itemize}

\needspace{5\baselineskip}\section*{[E] Довільний цикл (с. 405)}
\begin{itemize}
\item Будь-який оборотний цикл можна розбити сіткою ізотерм і адіабат на елементарні цикли Карно (рис. 16.13); внутрішні ділянки взаємно знищуються.
\item \[\oint\frac{d\text{'}Q}T = 0\ \text{(\text{оборотний})}\ (16.73),\qquad \oint\frac{d\text{'}Q}T \le 0\ \text{(\text{загалом})}\ (16.74)\]
\item (16.74) — другий закон для довільного циклу. Оскільки (16.73) дорівнює 0, d'Q/T має бути диференціалом \textbf{функції стану} — це \textbf{ентропія}. Продовження — на 14.
\item \pen{} Намалюйте рис. 16.13: овал, розбитий сіткою ізотерм і адіабат.
\end{itemize}

\needspace{5\baselineskip}\section*{Самоперевірка}
\begin{enumerate}
\item Формулювання Клаузіуса й Томсона звучать по-різному. Чому вони рівносильні?
\item Чому вічний двигун другого роду не порушує першого закону?
\item Що таке «зведена теплота» і чому сумують саме її, а не просто теплоту?
\end{enumerate}

\chapter*{14. Ентропія та її властивості}\addcontentsline{toc}{chapter}{14. Ентропія та її властивості}\markboth{14. Ентропія та її властивості}{14. Ентропія та її властивості}

\srcline{\book{} т. 1, \textbf{§ 16.8} (від с. 406, «Поняття про ентропію»), \textbf{§ 16.10} «Зв'язок ентропії з імовірністю стану системи» (с. 411–414); приклад 16.4 (с. 410–411). \ok{} (16.75)–(16.82) звірено зі сканом; (16.84), (16.91) — з чіткого OCR.}
Блок F, сторінка 4 з 5. Два погляди: \textbf{термодинамічний} (dS = d'Q/T) і \textbf{статистичний} (S = k ln w).

\needspace{5\baselineskip}\section*{План сторінки}
\begin{layout}
┌ [A] Означення: dS = d'Q/T (16.75) — рамка; Клаузіус 1854          ┐
├ [B] Властивості: функція стану · адитивна · до сталої (16.76)      ┤
├ [C] Закон зростання ентропії: dS ≥ d'Q/T (16.81), ізольована dS ≥ 0 (16.82) ┤
├ [D] ΔS для процесів ідеального газу (приклад 16.4) — табличка     ┤
├ [E] Статистичний зміст: S = k ln w (16.91); мікро-/макростан; рис. 16.15 ┤
└ низ: основне рівняння TdS = dU + pdV (16.84); «теплова смерть» — чому ні ┘
\end{layout}

\needspace{5\baselineskip}\section*{[A] Означення (с. 406) — у рамку}
\[dS = \frac{d\text{'}Q}T\qquad(16.75)\]
\begin{itemize}
\item З (16.73): d'Q/T — повний диференціал, отже існує функція стану S (англ. entropy). Її ввів Клаузіус у 1854 р.; назва з грецької — «перетворення, зміна».
\item Означення книги: «Ентропія — функція стану системи, яка визначається тим, що її повний диференціал dS при елементарному \textbf{оборотному} процесі дорівнює відношенню нескінченно малої кількості теплоти d'Q, що її дістала система, до термодинамічної температури T системи»
\item $d\text{'}Q = T\,dS$ — та сама структура, що й $d\text{'}A = p\,dV$: T — «узагальнена сила» теплообміну, S — його «узагальнена координата».
\end{itemize}

\needspace{5\baselineskip}\section*{[B] Властивості (с. 406)}
\begin{enumerate}
\item \textbf{Функція стану}: не залежить від шляху.
\item \textbf{Адитивна}: S = ΣSᵢ.
\item Визначена лише \textbf{з точністю до сталої}: $S = S_0 + \int\dfrac{d\text{'}Q}T$ (16.76). Сталу S₀ фіксує третій закон (15).
\end{enumerate}

\needspace{5\baselineskip}\section*{[C] Закон зростання ентропії (с. 406–408)}
\begin{itemize}
\item Необоротний перехід A→B з оборотним поверненням (рис. 16.14) дає $S_2 - S_1 > \int_1^2\dfrac{d\text{'}Q}T$ (16.80). Загалом:
\[dS \ge \frac{d\text{'}Q}T\qquad(16.81)\]
\item \textbf{Ізольована} система (d'Q = 0): $dS \ge 0$ (16.82). Рівність — для оборотних процесів у рівновазі (16.83).
\item «ентропія замкненої системи \textbf{не може зменшуватись}; вона або зростає, якщо в системі відбуваються необоротні процеси, або залишається незмінною …». Ентропія — \textbf{міра необоротності}.
\end{itemize}

\needspace{5\baselineskip}\section*{[D] ΔS для ідеального газу, 1 моль (с. 410–411, приклад 16.4)}
{\small
\begin{longtable}{>{\raggedright\arraybackslash}p{0.343\linewidth}>{\raggedright\arraybackslash}p{0.625\linewidth}}
\toprule
\textbf{процес} & \textbf{ΔS} \\
\midrule
\endhead
адіабатний (оборотний) & 0, \textbf{ізоентропійний} (англ. isentropic) \\
ізотермічний & Q/T = R ln(V₂/V₁) \\
ізохорний & $C_{V}$ ln(T₂/T₁) \\
ізобарний & $C_{p}$ ln(T₂/T₁) \\
\bottomrule
\end{longtable}
}
\begin{itemize}
\item Загалом (приклад 16.5): $\Delta S = \nu C_V\ln\dfrac{T_2}{T_1} + \nu R\ln\dfrac{V_2}{V_1}$
\end{itemize}

\needspace{5\baselineskip}\section*{[E] Статистичний зміст (с. 411–413)}
\begin{itemize}
\item \textbf{Макростан} (англ. macrostate): задані p, T, ρ… \textbf{Мікростан} (англ. microstate): конкретний розподіл частинок за координатами та імпульсами.
\item \textbf{Термодинамічна ймовірність w} (англ. thermodynamic probability): \textbf{кількість мікростанів}, що реалізують даний макростан. Вона ≥ 1, зазвичай величезне число, на відміну від математичної ймовірності. $w = \dfrac{n!}{n_1!\,n_2!\cdots n_N!}$ (16.90)
\item Приклад на с. 412: 6 молекул у 3 комірках. Розподіл 2-2-2 реалізується 90 способами, 6-0-0 — лише одним. \textbf{Рівномірний} розподіл найімовірніший — це і є рівновага.
\item \textbf{Принцип Больцмана} (1886):
\[S = k\ln w\qquad(16.91)\]
\item Зміст: «довільний природний процес протікає в напрямі \textbf{збільшення ймовірності стану}». Ентропія — міра \textbf{невпорядкованості} і близькості системи до рівноваги. Зменшення ентропії не заборонене, лише дуже малоймовірне (флуктуації).
\item \pen{} Намалюйте рис. 16.15: дві молекули у двох комірках з усіма варіантами розміщення.
\end{itemize}

\needspace{5\baselineskip}\section*{Низ}
\begin{itemize}
\item Основне рівняння (16.84): $T\,dS = dU + p\,dV$ — перший і другий закони разом.
\item \textbf{«Теплова смерть Всесвіту»} (с. 413–414): висновок хибний, бо Всесвіт не є термодинамічною системою.
\item Необов'язково: характеристичні функції H = U + pV, F = U − TS, G = U − TS + pV (с. 408–409).
\end{itemize}

\needspace{5\baselineskip}\section*{Самоперевірка}
\begin{enumerate}
\item Чому d'Q — не повний диференціал, а d'Q/T — повний?
\item Чи може ентропія ізольованої системи зменшитися? Дайте строгу (термодинамічну) і статистичну відповідь.
\item Чому рівномірний розподіл молекул відповідає рівновазі?
\end{enumerate}

\chapter*{15. Теорема Нернста}\addcontentsline{toc}{chapter}{15. Теорема Нернста}\markboth{15. Теорема Нернста}{15. Теорема Нернста}

\srcline{\book{} т. 1, \textbf{§ 16.11} «Третій закон термодинаміки. Недосяжність абсолютного нуля температур» (с. 414–415). У книзі це \textbf{третій закон}; «теорема Нернста» — його історична назва. \ok{} (16.93) і наслідки звірено зі сканом (с. 414–415).}
Блок F, сторінка 5 з 5. Коротка тема: пів сторінки.

\needspace{5\baselineskip}\section*{План сторінки}
\begin{layout}
┌ [A] Навіщо: ентропія відома лише до сталої (16.76) → потрібен новий постулат ┐
├ [B] Формулювання Нернста (1906) — рамка                                     ┤
├ [C] Формулювання Планка: lim S = 0 при T → 0 (16.93)                        ┤
└ [D] Наслідки: недосяжність 0 К · C_p, C_V → 0 · коефіцієнти розширення → 0  ┘
\end{layout}

\needspace{5\baselineskip}\section*{[A] Навіщо (с. 414)}
\begin{itemize}
\item Перший і другий закони дають лише \textbf{зміни} ентропії; S₀ у (16.76) лишається невизначеною. Цього досить, доки не змінюється \textbf{хімічний склад} — це проблема хімічної спорідненості.
\item Тож потрібен був \textbf{новий постулат}, що узагальнює нові досліди при низьких температурах.
\end{itemize}

\needspace{5\baselineskip}\section*{[B] Формулювання Нернста (1906) (с. 414–415)}
\begin{itemize}
\item «при наближенні температури до абсолютного нуля ентропія будь-якої рівноважної системи при ізотермічних процесах \textbf{перестає залежати від термодинамічних параметрів стану} і при T → 0 набуває \textbf{ту саму для всіх систем сталу величину}, яку можна вважати такою, що дорівнює нулю» (англ. Nernst heat theorem)
\item Нернст розглядав лише \textbf{конденсовані} системи — тверді тіла й рідини.
\end{itemize}

\needspace{5\baselineskip}\section*{[C] Формулювання Планка (с. 415)}
\begin{itemize}
\item Планк пішов далі: сама ентропія прямує до нуля, для \textbf{будь-якої} системи:
\[\lim_{T\to0}S = 0\qquad(16.93)\]
— найзагальніше формулювання третього закону (англ. third law of thermodynamics).
\end{itemize}

\needspace{5\baselineskip}\section*{[D] Наслідки (с. 415)}
\begin{enumerate}
\item \textbf{Абсолютний нуль недосяжний.} Усі методи охолодження використовують адіабатні процеси. Поблизу 0 К різниця між адіабатним та ізотермічним процесами зникає, тож охолодження стає дедалі менш ефективним. До 0 К можна наближатися лише асимптотично. Закон іноді формулюють як «принцип недосяжності абсолютного нуля».
\item $C_p, C_V \to 0$ при T → 0. Це підтверджено дослідом і квантовими теоріями Ейнштейна та Дебая; порівняйте з 05, де класична теорія не справляється.
\item Термічні коефіцієнти розширення й тиску → 0.
\end{enumerate}

\needspace{5\baselineskip}\section*{Три закони разом (добре підходить для низу сторінки)}
{\small
\begin{longtable}{>{\raggedright\arraybackslash}p{0.099\linewidth}>{\raggedright\arraybackslash}p{0.869\linewidth}}
\toprule
\textbf{закон} & \textbf{суть} \\
\midrule
\endhead
I & енергія зберігається: Q = ΔU + A \\
II & напрям: S ізольованої системи не зменшується \\
III & S → 0 при T → 0; 0 К недосяжний \\
\bottomrule
\end{longtable}
}

\needspace{5\baselineskip}\section*{Самоперевірка}
\begin{enumerate}
\item Чому перші два закони не могли зафіксувати значення ентропії?
\item Чому адіабатне охолодження перестає працювати поблизу 0 К?
\end{enumerate}

\chapter*{26. Відмінність реального газу від ідеального}\addcontentsline{toc}{chapter}{26. Відмінність реального газу від ідеального}\markboth{26. Відмінність реального газу від ідеального}{26. Відмінність реального газу від ідеального}

\srcline{\book{} т. 1, \textbf{§ 17.1} «Реальні гази та відхилення їхніх властивостей від законів ідеального газу. Сили міжмолекулярної взаємодії в газах» (с. 419–422). \ok{} Звірено зі сканом (с. 419–421).}
Блок G, сторінка 1 з 2.

\needspace{5\baselineskip}\section*{План сторінки}
\begin{layout}
┌ [A] Означення реального газу; дві причини відхилень (власний об'єм, притягання) ┐
├ [B] Кількісні vs якісні відхилення; коли модель ідеального газу ще працює    ┤
├ [C] ГРАФІК сили F(r): відштовхування/притягання, r₀ (рис. 17.2)              ┤
├ [D] Сили Ван дер Ваальса: Кеєзом · Дебай · Лондон (таблиця) + водневий зв'язок ┤
└ [E] Ізотерми Амага, крива і температура Бойля (необов'язково)                 ┘
\end{layout}

\needspace{5\baselineskip}\section*{[A] Реальний газ (с. 419)}
\begin{itemize}
\item «\textbf{Реальними} називають гази, для яких суттєву роль відіграють міжмолекулярні сили взаємодії» (англ. real gas)
\item Порівняно з ідеальною моделлю (невзаємодійні точкові молекули) додаються дві речі:
\begin{enumerate}
\item \textbf{власний об'єм молекул}. За нормальних умов він становить ≈ 10⁻⁴ об'єму газу, але при ≈ 10⁸ Па стає суттєвим;
\item \textbf{притягання між молекулами}, що дає додатковий тиск.
\end{enumerate}
\end{itemize}

\needspace{5\baselineskip}\section*{[B] Характер відхилень (с. 419–420)}
\begin{itemize}
\item \textbf{Кількісні} (англ. quantitative): закони ідеального газу, зокрема pV = RT, справджуються лише наближено.
\item \textbf{Якісні} (англ. qualitative): завдяки взаємодії газ може перейти в \textbf{рідкий і твердий} стан, що неможливо для ідеального газу.
\item Закони ідеального газу справджуються при \textbf{низьких тисках і високих температурах}, коли середня потенціальна енергія взаємодії ≪ середньої кінетичної.
\end{itemize}

\needspace{5\baselineskip}\section*{[C] Сили міжмолекулярної взаємодії (с. 420–421)}
\begin{itemize}
\item Відштовхування $F_1 \propto 1/r^n$ (n = 9–15) діє на малих відстанях; притягання $F_2 \propto 1/r^m$ (m = 7) — на більших (англ. intermolecular forces).
\item При r = r₀ F = 0 — \textbf{рівноважна відстань} (англ. equilibrium distance). При r > r₀ переважає притягання, при r < r₀ — відштовхування.
\item \pen{} Головний малюнок, рис. 17.2: F(r) — крива відштовхування згори, притягання знизу, їхня сума і нуль у r₀. Можна додати криву потенціальної енергії з ямою в r₀.
\end{itemize}

\needspace{5\baselineskip}\section*{[D] Сили Ван дер Ваальса (с. 421–422)}
\begin{itemize}
\item Слабкі сили притягання на відстанях ≈ 10⁻⁹ м (англ. van der Waals forces).
\end{itemize}
{\footnotesize
\begin{longtable}{>{\raggedright\arraybackslash}p{0.138\linewidth}>{\raggedright\arraybackslash}p{0.069\linewidth}>{\raggedright\arraybackslash}p{0.207\linewidth}>{\raggedright\arraybackslash}p{0.288\linewidth}>{\raggedright\arraybackslash}p{0.218\linewidth}}
\toprule
\textbf{тип} & \textbf{автор} & \textbf{у яких молекул} & \textbf{природа} & \textbf{особливості} \\
\midrule
\endhead
орієнтаційні & Кеєзом & полярних (H₂O, NH₃) & постійні диполі & залежать від T, неадитивні \\
індукційні & Дебай & полярних + індукованих & індукований диполь & малі, не залежать від T \\
дисперсійні & Лондон & \textbf{неполярних} (основні) & миттєві диполі через рух електронів & адитивні \\
\bottomrule
\end{longtable}
}
\begin{itemize}
\item \textbf{Водневий зв'язок} (англ. hydrogen bond), с. 422: специфічне притягання у воді й спиртах через групи OH.
\end{itemize}

\needspace{5\baselineskip}\section*{[E] Необов'язково: ізотерми Амага (с. 420)}
\begin{itemize}
\item У координатах pV–p для ідеального газу — прямі; криві реального газу спочатку опускаються (притягання), потім піднімаються (власний об'єм). Мінімуми лежать на \textbf{кривій Бойля}; $T_{\text{Б}}$ — \textbf{температура Бойля}.
\end{itemize}

\needspace{5\baselineskip}\section*{Самоперевірка}
\begin{enumerate}
\item Назвіть дві причини відхилень від ідеального газу і скажіть, коли переважає кожна.
\item Чому ідеальний газ не можна зрідити?
\item Які сили Ван дер Ваальса діють між молекулами N₂?
\end{enumerate}

\chapter*{27. Рівняння Ван-Дер-Ваальса}\addcontentsline{toc}{chapter}{27. Рівняння Ван-Дер-Ваальса}\markboth{27. Рівняння Ван-Дер-Ваальса}{27. Рівняння Ван-Дер-Ваальса}

\srcline{\book{} т. 1, \textbf{§ 17.2} «Рівняння Ван-дер-Ваальса» (с. 422–424), \textbf{§ 17.3} «Ізотерми газу Ван-дер-Ваальса. Критичний стан речовини…» (с. 425–430), \textbf{§ 17.4} (с. 430–432). \ok{} (17.4), (17.9)–(17.12) звірено зі сканом.}
Блок G, сторінка 2 з 2. Правопис: у книзі — \textbf{Ван-дер-Ваальса}.

\needspace{5\baselineskip}\section*{План сторінки}
\begin{layout}
┌ [A] Рівняння (17.4) у рамці + що означає кожна поправка          ┐
├ [B] Звідки a/V² (внутрішній тиск) і b (4 × об'єм молекул)        ┤
├ [C] ГРАФІК ізотерм (рис. 17.4/17.5): хвиля, горизонталь, критична точка K ┤
├ [D] Критичні параметри (17.12); критичний стан                  ┤
└ [E] Порівняння з дослідом; зведене рівняння (17.18)              ┘
\end{layout}

\needspace{5\baselineskip}\section*{[A] Рівняння (с. 423) — у рамку}
\[\left(p + \frac a{V_M^2}\right)(V_M - b) = RT\qquad(17.4)\]
\begin{itemize}
\item Для одного моля (англ. van der Waals equation). $V_{M}$ — молярний об'єм; \textbf{a, b} — дослідні сталі газу.
\item \textbf{a/$V_{M}$²} враховує \textbf{притягання} молекул. Має розмірність тиску й називається \textbf{внутрішнім тиском} (англ. internal pressure), $p_{i.}$
\item \textbf{b} — поправка на \textbf{власний об'єм} молекул (відштовхування на малих відстанях).
\item Для маси m (17.7): $\left(p + \dfrac{m^2}{M^2}\dfrac a{V^2}\right)\left(V - \dfrac mMb\right) = \dfrac mMRT$
\item Рівняння наближене: кількісно добре при високих T і низьких p, але \textbf{якісно} справджується скрізь.
\end{itemize}

\needspace{5\baselineskip}\section*{[B] Звідки поправки (с. 423–424)}
\begin{itemize}
\item \textbf{b}: центри двох твердих кульок радіуса r не можуть зблизитися менш ніж на 2r, тож «заборонений» об'єм на одну молекулу ½·(4/3)π(2r)³ = 4·(4/3)πr³. Отже, \textbf{b = 4 × об'єм молекул} у молі.
\item \textbf{a/V²}: молекулу біля стінки тягнуть назад у газ. Ця сила ∝ n₀, і кількість молекул, що б'ються об стінку, теж ∝ n₀, тож $p_{i}$ ∝ n₀² ∝ 1/V².
\item \pen{} Намалюйте молекулу біля стінки із силами від сусідів, що тягнуть її всередину.
\end{itemize}

\needspace{5\baselineskip}\section*{[C] Ізотерми (с. 425–430)}
\begin{itemize}
\item Після розкриття дужок маємо \textbf{кубічне рівняння відносно $V_{M}$} (17.8): при заданих p і T може бути \textbf{три} об'єми.
\item Нижче $T_{\text{к}}$: S-подібна «хвиля». На \textbf{дослідній} ізотермі її заміняє \textbf{горизонтальна} ділянка (рис. 17.5), де рідина й пара співіснують при сталому p.
\begin{itemize}
\item Права гілка — газ; ліва, крута, — рідина (майже нестислива).
\item \textbf{Метастабільні стани} (англ. metastable states): ділянка bc — \textbf{пересичена пара} (англ. supersaturated vapour); ef — \textbf{перегріта рідина} (англ. superheated liquid).
\item Ділянка cde абсолютно нестійка (∂p/∂V > 0) і не реалізується.
\end{itemize}
\item \textbf{Бінодаль} (англ. binodal) обмежує двофазну область; \textbf{спінодаль} (англ. spinodal) — область нестійких станів.
\item Вище $T_{\text{к}}$: монотонні гіперболи, тож \textbf{лише тиском газ не зрідити} (CO₂ вище 31 °C).
\item \pen{} Головний малюнок: сім'я ізотерм T₁ < $T_{\text{к}}$ < T₂, хвиля й горизонталь на T₁, точка перегину K і дзвоноподібна бінодаль з підписами «рідина», «рідина + пара», «газ».
\end{itemize}

\needspace{5\baselineskip}\section*{[D] Критичний стан (с. 426–427)}
\begin{itemize}
\item При T = $T_{\text{к}}$ три корені зливаються, і ізотерма має \textbf{точку перегину K} (англ. inflection point).
\[p_\text{к} = \frac a{27b^2},\qquad V_\text{к} = 3b,\qquad T_\text{к} = \frac{8a}{27Rb}\qquad(17.12)\]
\item У критичному стані різниця між рідиною й газом зникає, а стисливість → ∞.
\item Газ можна перевести в рідину \textbf{без фазового переходу}, обійшовши двофазну область через T > $T_{\text{к}}$ (с. 430).
\end{itemize}

\needspace{5\baselineskip}\section*{[E] Порівняння з дослідом (с. 430–432), коротко}
\begin{itemize}
\item a і b залежать від T. \textbf{Критичний коефіцієнт} $K_\text{к} = \dfrac{RT_\text{к}}{p_\text{к}V_\text{к}} = \dfrac83 \approx 2{,}67$ теоретично, але для H₂ — 3,12, для води — 4,46. Насправді $V_{\text{к}}$ ≈ 2b.
\item \textbf{Зведене рівняння} (англ. reduced equation) з π = p/$p_{\text{к}}$, ω = V/$V_{\text{к}}$, τ = T/$T_{\text{к}}$: $\left(\pi + \dfrac3{\omega^2}\right)(3\omega - 1) = 8\tau$ (17.18) — однакове для всіх речовин. Це \textbf{закон відповідних станів} (англ. law of corresponding states).
\end{itemize}

\needspace{5\baselineskip}\section*{Самоперевірка}
\begin{enumerate}
\item Яка стала відповідає притяганню, а яка — відштовхуванню?
\item Чому b учетверо більша за власний об'єм молекул, а не дорівнює йому?
\item Чому газ вище $T_{\text{к}}$ не можна зрідити, хоч як його стискати?
\end{enumerate}

\chapter*{28. Класифікація кристалів}\addcontentsline{toc}{chapter}{28. Класифікація кристалів}\markboth{28. Класифікація кристалів}{28. Класифікація кристалів}

\srcline{\book{} т. 1, \textbf{§ 18.1} «Кристалічні та аморфні тіла. Основні характеристики кристалів» (с. 444–450); класифікація за типом зв'язку — \textbf{§ 18.2} «Сили зв'язку у твердих тілах» (с. 450–453). \ok{} Звірено зі сканом (с. 446–448).}
Блок H, сторінка 1 з 3. Книга класифікує кристали \textbf{двома способами}: за \textbf{симетрією} (сингонії, гратки Браве) і за \textbf{типом зв'язку}. Розгляньте обидва.

\needspace{5\baselineskip}\section*{План сторінки}
\begin{layout}
┌ [A] Тверде тіло: кристалічне vs аморфне (дальній/ближній порядок, анізотропія) ┐
├ [B] Моно- і полікристали; закони кристалографії; елементи симетрії          ┤
├ [C] Гратка, вузли, елементарна комірка, rₙ = ma₁ + na₂ + pa₃ (18.1)          ┤
├ [D] ТАБЛИЦЯ 7 сингоній + 4 типи граток Браве (14 разом) + малюнки кубічних   ┤
└ [E] Класифікація за типом зв'язку: іонні · ковалентні · металічні · молекулярні ┘
\end{layout}

\needspace{5\baselineskip}\section*{[A] Кристалічні та аморфні тіла (с. 444–445)}
\begin{itemize}
\item \textbf{Тверде тіло} (англ. solid): «агрегатний стан речовини, який характеризується стабільністю форми та об'єму», а також упорядкованим розміщенням атомів, що здійснюють малі коливання біля положень рівноваги.
\end{itemize}
{\small
\begin{longtable}{>{\raggedright\arraybackslash}p{0.143\linewidth}>{\raggedright\arraybackslash}p{0.503\linewidth}>{\raggedright\arraybackslash}p{0.306\linewidth}}
\toprule
\textbf{} & \textbf{кристалічні} & \textbf{аморфні} \\
\midrule
\endhead
порядок & \textbf{дальній} (англ. long-range): строга періодична повторюваність у 3D & лише \textbf{ближній} (англ. short-range) \\
властивості & \textbf{анізотропія} (англ. anisotropy): властивості залежать від напряму & \textbf{ізотропні} (англ. isotropic) \\
приклади & метали, сіль, кварц & скло, бурштин, смола, полімери \\
\bottomrule
\end{longtable}
}
\begin{itemize}
\item \textbf{Монокристал} (англ. single crystal; один великий правильний кристал) vs \textbf{полікристал} (англ. polycrystal; багато дрібних зерен, хаотично зрощених; більшість металів).
\end{itemize}

\needspace{5\baselineskip}\section*{[B] Закони кристалографії і симетрія (с. 445–446)}
\begin{itemize}
\item Плоскі грані й прямі ребра; \textbf{сталість кутів між гранями}; раціональні відношення відрізків на осях.
\item \textbf{Закон симетрії}: можливі лише осі порядку \textbf{1, 2, 3, 4 і 6} (осі 5-го порядку немає).
\item Елементи симетрії: \textbf{вісь симетрії} (англ. symmetry axis; поворот на 2π/n), \textbf{площина симетрії} P (англ. mirror plane; дзеркальна) і \textbf{центр інверсії} C (англ. centre of inversion).
\end{itemize}

\needspace{5\baselineskip}\section*{[C] Просторова гратка (с. 446–447)}
\begin{itemize}
\item Поставивши точку в центр мас кожної структурної одиниці, отримують \textbf{просторові гратки} (англ. space lattice); їхні точки — \textbf{вузли} (англ. nodes).
\item \textbf{Елементарна комірка} (англ. unit cell): паралелепіпед на векторах a₁, a₂, a₃; їхні довжини — \textbf{періоди ідентичності} (англ. lattice periods). Будь-який вузол: $\vec r_n = m\vec a_1 + n\vec a_2 + p\vec a_3$ (18.1).
\item \textbf{Гратки Браве} (англ. Bravais lattices): комірки мінімальної площі, що зберігають повну симетрію. Їх \textbf{14} у \textbf{7 сингоніях} (англ. crystal systems).
\end{itemize}

\needspace{5\baselineskip}\section*{[D] 7 сингоній (с. 447–449, табл. 18.4)}
{\small
\begin{longtable}{>{\raggedright\arraybackslash}p{0.304\linewidth}>{\raggedright\arraybackslash}p{0.105\linewidth}>{\raggedright\arraybackslash}p{0.246\linewidth}>{\raggedright\arraybackslash}p{0.281\linewidth}}
\toprule
\textbf{сингонія} & \textbf{ребра} & \textbf{кути} & \textbf{приклади (з книги)} \\
\midrule
\endhead
триклінна & a ≠ b ≠ c & α ≠ β ≠ γ ≠ 90° & польові шпати \\
моноклінна & a ≠ b ≠ c & α = γ = 90° ≠ β & гіпс, слюда, тальк \\
ромбічна & a ≠ b ≠ c & α = β = γ = 90° & сірка, топаз \\
ромбоедрична (тригональна) & a = b = c & α = β = γ ≠ 90° & кальцит, телур \\
тетрагональна & a = b ≠ c & α = β = γ = 90° & рутил, олов'яний камінь \\
гексагональна & a = b ≠ c & α = β = 90°, γ = 120° & лід, графіт, цинк \\
кубічна & a = b = c & α = β = γ = 90° & більшість металів, алмаз \\
\bottomrule
\end{longtable}
}
\begin{itemize}
\item Типи граток: \textbf{примітивна} (англ. primitive; вузли у вершинах), \textbf{базоцентрована} (англ. base-centred), \textbf{об'ємноцентрована} (англ. body-centred: Fe-α, W, Na), \textbf{гранецентрована} (англ. face-centred: Al, Cu, Ag, Au).
\item \pen{} Намалюйте три кубічні гратки: просту, об'ємноцентровану й гранецентровану.
\item Необов'язково: \textbf{індекси Міллера} (h k l) для площин, с. 449; \textbf{230 просторових груп} (Федоров, 1890); \textbf{32 класи} симетрії.
\end{itemize}

\needspace{5\baselineskip}\section*{[E] За типом зв'язку (§ 18.2, с. 450–453)}
{\footnotesize
\begin{longtable}{>{\raggedright\arraybackslash}p{0.138\linewidth}>{\raggedright\arraybackslash}p{0.115\linewidth}>{\raggedright\arraybackslash}p{0.261\linewidth}>{\raggedright\arraybackslash}p{0.285\linewidth}>{\raggedright\arraybackslash}p{0.121\linewidth}}
\toprule
\textbf{тип кристала} & \textbf{у вузлах} & \textbf{зв'язок} & \textbf{властивості} & \textbf{приклади} \\
\midrule
\endhead
іонний & ± іони & електростатичний & міцний, висока $T_{\text{пл}}$, поганий провідник & NaCl \\
ковалентний (атомний) & нейтральні атоми & спільні електронні пари, \textbf{напрямлений, насичений} & дуже тверді, висока $T_{\text{пл}}$, діелектрики або напівпровідники & алмаз, Si \\
металічний & + іони & \textbf{електронний газ} & добрі електро- й теплопровідники, блиск, ковкість & Cu, Au \\
молекулярний & молекули & ван-дер-ваальсів, водневий & низька $T_{\text{пл}}$, ізолятори & органічні кристали, лід \\
\bottomrule
\end{longtable}
}

\needspace{5\baselineskip}\section*{Самоперевірка}
\begin{enumerate}
\item Чим відрізняються дальній і ближній порядок?
\item Чому кристал не може мати вісь симетрії 5-го порядку? (П'ятикутниками не замостити площину.)
\item Чому метали проводять струм, а ковалентні кристали — ні?
\end{enumerate}

\chapter*{29. Дефекти в кристалах}\addcontentsline{toc}{chapter}{29. Дефекти в кристалах}\markboth{29. Дефекти в кристалах}{29. Дефекти в кристалах}

\srcline{\book{} т. 1, \textbf{§ 18.3} «Дефекти в реальних кристалах» (с. 453–456). \ok{} Звірено зі сканом (с. 453–456).}
Блок H, сторінка 2 з 3. Структура проста: \textbf{класифікація за вимірністю (0D–3D)}.

\needspace{5\baselineskip}\section*{План сторінки}
\begin{layout}
┌ [A] Навіщо: дефекти визначають структурно-чутливі властивості           ┐
├ [B] Класифікація за вимірністю — головна таблиця 0D/1D/2D/3D            ┤
├ [C] Точкові: Шотткі, Френкель, домішки (рис. 18.5 а/б/в) + n_в (18.4)    ┤
├ [D] Лінійні: дислокації — крайова (рис. 18.6) і гвинтова (рис. 18.7); вектор Бюргерса ┤
└ [E] Поверхневі та об'ємні (коротко)                                     ┘
\end{layout}

\needspace{5\baselineskip}\section*{[A] Навіщо (с. 453)}
\begin{itemize}
\item Ідеальна структура існує лише в малих об'ємах. Дефекти значною мірою визначають \textbf{структурно-чутливі властивості} (англ. structure-sensitive properties): електро- й теплопровідність, міцність…
\end{itemize}

\needspace{5\baselineskip}\section*{[B] Класифікація (с. 453) — головна таблиця}
{\small
\begin{longtable}{>{\raggedright\arraybackslash}p{0.169\linewidth}>{\raggedright\arraybackslash}p{0.391\linewidth}>{\raggedright\arraybackslash}p{0.392\linewidth}}
\toprule
\textbf{вимірність} & \textbf{тип} & \textbf{приклади} \\
\midrule
\endhead
0 (нульвимірні) & \textbf{точкові} (англ. point defects): розмір \textasciitilde{} одна або кілька міжатомних відстаней у всіх напрямах & вакансії (Шотткі), пари Френкеля, домішкові атоми \\
1 (одновимірні) & \textbf{лінійні} (англ. line defects) & \textbf{дислокації}: крайові й гвинтові \\
2 (двовимірні) & \textbf{поверхневі} (англ. planar defects) & межі зерен, блоки мозаїки, дефекти упаковки, двійники, поверхня \\
3 (тривимірні) & \textbf{об'ємні} (англ. volume defects) & тріщини, включення, пори \\
\bottomrule
\end{longtable}
}

\needspace{5\baselineskip}\section*{[C] Точкові дефекти (с. 453–455)}
\begin{itemize}
\item \textbf{Атомні}:
\begin{itemize}
\item \textbf{дефект за Шотткі} (англ. Schottky defect): незайнятий вузол, вакансія (рис. 18.5 а);
\item \textbf{дефект за Френкелем} (англ. Frenkel defect): вакансія \textbf{плюс} близько розміщений міжвузловий атом (рис. 18.5 в). Потребує трохи меншої енергії, ніж дефект Шотткі; в обох випадках 2–4 еВ;
\item \textbf{домішкові атоми} (англ. impurity atoms): у міжвузлях або на місці атомів основної речовини (рис. 18.5 б).
\end{itemize}
\item Електронні (вільні електрони, дірки, \textbf{екситон}, \textbf{полярон} — теорія Пекаря) та енергетичні (\textbf{фонони} — кванти коливань ґратки): згадайте одним рядком.
\item Рівноважна концентрація вакансій: $n_{\text{в}} = n_0e^{-W_{\text{в}}/kT}$ (18.4). Швидко зростає з T; поблизу плавлення $n_{\text{в}}$/n₀ ≈ 10⁻⁴. Порівняйте з розподілом Больцмана на 17, 19.
\item \pen{} Намалюйте три маленькі квадратні сітки: з відсутнім атомом, з атомом, зміщеним із вузла в міжвузля, і з чужим атомом.
\end{itemize}

\needspace{5\baselineskip}\section*{[D] Дислокації (с. 455–456)}
\begin{itemize}
\item \textbf{Крайова} (англ. edge dislocation, рис. 18.6): одна атомна площина обривається всередині кристала — як книга, в якій одну сторінку наполовину вирвано.
\item \textbf{Гвинтова} (англ. screw dislocation, рис. 18.7): дві півплощини зсунуто на один період одна відносно одної.
\item Дислокації не можуть обриватися всередині кристала: утворюють петлі або виходять на поверхню.
\item \textbf{Вектор Бюргерса} (англ. Burgers vector): обходять контур цілим числом кроків. Навколо дислокації контур \textbf{не замикається}, і проміжок a′a″ — це вектор Бюргерса, міра спотворення ґратки.
\item \textbf{Густина дислокацій} (англ. dislocation density): кількість дислокаційних ліній, що перетинають одиничну площу.
\item \pen{} Намалюйте «зайву півплощину» з контуром Бюргерса навколо.
\end{itemize}

\needspace{5\baselineskip}\section*{[E] Поверхневі та об'ємні (с. 456)}
\begin{itemize}
\item \textbf{Двійники} (англ. twins): два кристали однієї речовини, що зрослися вздовж кристалографічної площини, яка є площиною дзеркальної симетрії.
\end{itemize}

\needspace{5\baselineskip}\section*{Самоперевірка}
\begin{enumerate}
\item Шотткі vs Френкель: у чому різниця?
\item Чому кількість вакансій зростає з температурою і за яким законом?
\item Як виявити дислокацію за допомогою контуру Бюргерса?
\end{enumerate}

\chapter*{30. Рідини}\addcontentsline{toc}{chapter}{30. Рідини}\markboth{30. Рідини}{30. Рідини}

\srcline{\book{} т. 1, \textbf{§ 19.1} «Загальні властивості та будова рідин» (с. 462–466), \textbf{§ 19.2} «В'язкість рідин» (с. 467). \ok{} Звірено зі сканом (с. 462–467).}
Блок H, сторінка 3 з 3. \warn{} Це питання \textbf{дописане від руки} на аркуші МКР («30. Рідини»), тож його обсяг неясний. Сторінка охоплює будову й властивості рідин (§ 19.1–19.2). Якщо лектор має на увазі поверхневий натяг і капілярність, див. «можливе розширення» наприкінці.

\needspace{5\baselineskip}\section*{План сторінки}
\begin{layout}
┌ [A] Рідина між газом і твердим тілом: ближній порядок без дальнього     ┐
├ [B] ГРАФІК радіальної функції розподілу N(r): тверде/рідина/газ (рис. 19.1) ┤
├ [C] Тепловий рух за Френкелем: коливання + стрибки; час «осілого життя» τ ┤
├ [D] Властивості: мала стисливість β_T, теплове розширення α_p, аномалія води ┤
└ [E] Перенос у рідинах: D, λ, η — і їх залежність від T (протилежна газам!) ┘
\end{layout}

\needspace{5\baselineskip}\section*{[A] Місце рідин (с. 462–463)}
\begin{itemize}
\item Три агрегатні стани відрізняються тим, наскільки \textbf{взаємодія молекул сильна порівняно з їхнім тепловим рухом}.
\item Рідини (англ. liquids) займають \textbf{проміжне положення між газами й твердими тілами}. Молекули взаємодіють сильно, як у твердому тілі, але не розміщені у вузлах періодичної ґратки.
\item Ідея Френкеля: рідина — це «\textbf{розупорядковане тверде тіло}» з \textbf{ближнім порядком} і \textbf{без дальнього порядку} (28).
\end{itemize}

\needspace{5\baselineskip}\section*{[B] Радіальна функція розподілу (с. 462–463)}
\begin{itemize}
\item N(r) — кількість сусідів на відстані r від обраної молекули (англ. radial distribution function).
\begin{itemize}
\item Тверде тіло: гострі дискретні піки аж до великих r (фіксовані положення).
\item Рідина: кілька розмитих максимумів, що затухають до сталої (ближній порядок).
\item Газ: плавна крива без структури.
\end{itemize}
\item \textbf{Модель Бернала}: структуру описують радіуси координаційних сфер, кількість молекул у них і ступінь їхньої розмитості.
\item \pen{} Намалюйте поруч три маленькі графіки N(r), як на рис. 19.1 а/б/в.
\end{itemize}

\needspace{5\baselineskip}\section*{[C] Тепловий рух (Френкель) (с. 463–464)}
\begin{itemize}
\item \textbf{Коливання} біля тимчасового положення рівноваги (період ≈ 10⁻¹² с), потім \textbf{стрибок} у сусіднє положення. Ці стрибки — \textbf{трансляційні рухи}; для них потрібні вакантні місця, «дірки».
\item \textbf{Час осілого життя} (англ. residence time): $\tau = Ae^{\Delta W/kT}$ (19.1), де ΔW — потенціальний бар'єр. Зменшується з ростом T.
\item Стрибки пояснюють \textbf{текучість} і \textbf{дифузію} в рідинах.
\end{itemize}

\needspace{5\baselineskip}\section*{[D] Властивості (с. 464–466)}
\begin{itemize}
\item Зберігають \textbf{об'єм} (але не форму) і мають \textbf{вільну поверхню}.
\item \textbf{Мала стисливість} (англ. low compressibility): $\beta_T = -\dfrac1V\left(\dfrac{\partial V}{\partial p}\right)_T$ (19.2), ≈ 10⁻⁹–10⁻¹⁰ Па⁻¹ — у десятки й сотні тисяч разів менша, ніж у газів. Це прояв сильного відштовхування.
\item \textbf{Теплове розширення} (англ. thermal expansion): $\alpha_p = \dfrac1V\left(\dfrac{\partial V}{\partial T}\right)_p$ (19.3), зазвичай > 0.
\begin{itemize}
\item \spark{} \textbf{Аномалія води}: від 0 до 3,98 °C вода при нагріванні \textbf{стискається}, тож її густина максимальна при 3,98 °C. Причина — ажурна тетраедрична структура з водневими зв'язками.
\end{itemize}
\item Теплоємність близька до теплоємності твердих тіл і значно більша, ніж у газів, бо теплота розподіляється між кінетичною та потенціальною енергіями.
\end{itemize}

\needspace{5\baselineskip}\section*{[E] Перенос у рідинах (с. 466–467) — порівняйте з газами (23–25)}
{\small
\begin{longtable}{>{\raggedright\arraybackslash}p{0.104\linewidth}>{\raggedright\arraybackslash}p{0.230\linewidth}>{\raggedright\arraybackslash}p{0.618\linewidth}}
\toprule
\textbf{} & \textbf{гази} & \textbf{рідини} \\
\midrule
\endhead
механізм & вільні прольоти на ⟨l⟩ & стрибки між положеннями рівноваги \\
D & ⅓⟨v⟩⟨l⟩; 0,1–0,9 см²/с & ⟨δ²⟩/(6τ) (19.4); на 4–5 порядків менший; $D = Be^{-\Delta W/kT}$ (19.5) \\
η від T & \textbf{зростає} з T & \textbf{спадає} з T: $\eta = Ae^{\Delta W/kT}$, \textbf{формула Френкеля} (19.7) \\
λ від T & зростає & спадає приблизно лінійно \\
\bottomrule
\end{longtable}
}
\begin{itemize}
\item В'язкість води від 20 до 30 °C зменшується на 20 \%, гліцерину — у 2,4 раза.
\end{itemize}

\needspace{5\baselineskip}\section*{Можливе розширення (якщо «рідини» охоплюють поверхневі явища)}
\begin{itemize}
\item § 19.3 «Поверхневий шар рідини. Поверхневий натяг» (с. 468–471), § 19.4 «Формула Лапласа. Змочування та капілярні явища» (с. 471–474). У вікі цього ще немає. Уточніть у лектора і, якщо треба, попросіть додати сторінку.
\end{itemize}

\needspace{5\baselineskip}\section*{Самоперевірка}
\begin{enumerate}
\item Чому рідина ближча до твердого тіла, ніж до газу?
\item Чому в'язкість рідини з ростом T спадає, а газу — зростає?
\item Чому вода найгустіша при 4 °C?
\end{enumerate}

\appendix

\part{Додатки}

\chapter*{Інтегральна ↔ диференціальна форма законів}\addcontentsline{toc}{chapter}{Інтегральна ↔ диференціальна форма законів}\markboth{Інтегральна ↔ диференціальна форма законів}{Інтегральна ↔ диференціальна форма законів}

Та сама схема трапляється щонайменше в чотирьох темах програми. Інтегральна форма пов'язує \textbf{різні місця} (поверхню й об'єм або контур). Диференціальна — величини \textbf{в одній точці}, чого й вимагає близькодія. Містками слугують дві математичні теореми (англ. integral vs differential form).

{\footnotesize
\begin{longtable}{>{\raggedright\arraybackslash}p{0.209\linewidth}>{\raggedright\arraybackslash}p{0.172\linewidth}>{\raggedright\arraybackslash}p{0.172\linewidth}>{\raggedright\arraybackslash}p{0.270\linewidth}>{\raggedright\arraybackslash}p{0.098\linewidth}}
\toprule
\textbf{Закон} & \textbf{інтегральна} & \textbf{диференціальна} & \textbf{міст} & \textbf{теми} \\
\midrule
\endhead
Гаусс (джерела E) & ∮E dS = q/ε₀ & div E = ρ/ε₀ & Остроградського–Гаусса & 6, 7, 13 \\
потенціальність E & ∮E dl = 0 & rot E = 0 & Стокса & 4, 9 \\
збереження заряду & ∮j dS = −∂q/∂t & ∂ρ/∂t = −div j & Остроградського–Гаусса & 1, 20 \\
Ом & I = U/R & j = σE & (малий циліндр dl·dS) & 19 \\
Джоуль–Ленц & Q = I²Rt & w = σE² = jE & (малий циліндр) & 23 \\
\bottomrule
\end{longtable}
}

Як це формулює книга (т. 2, с. 41): «Теорема Гаусса дає можливість звести основні рівняння електростатики до диференціальної форми й узгодити їх, таким чином, з теорією близькодії.»

\chapter*{Принцип суперпозиції}\addcontentsline{toc}{chapter}{Принцип суперпозиції}\markboth{Принцип суперпозиції}{Принцип суперпозиції}

Поля від кількох джерел додаються, \textbf{не впливаючи одне на одне} (англ. superposition principle). Це дослідний факт, а не наслідок закону Кулона (т. 2, с. 25).

{\small
\begin{longtable}{>{\raggedright\arraybackslash}p{0.177\linewidth}>{\raggedright\arraybackslash}p{0.538\linewidth}>{\raggedright\arraybackslash}p{0.169\linewidth}>{\raggedright\arraybackslash}p{0.052\linewidth}}
\toprule
\textbf{Величина} & \textbf{як додається} & \textbf{формула} & \textbf{тема} \\
\midrule
\endhead
сила & векторна сума & F = ΣFᵢ (1.10) & 2 \\
напруженість & векторна сума & E = ΣEᵢ (1.13) & 3 \\
потенціал & \textbf{алгебрична} сума & φ = Σφᵢ (1.56) & 4 \\
поле двох площин & між пластинами поля додаються, зовні компенсуються & σ/ε₀ & 11 \\
поле диполя & додаються поля +q і −q & (1.20)–(1.22) & 5 \\
\bottomrule
\end{longtable}
}

Формулювання книги: «суперпозиція — це здатність фізичних полів до накладання, при якому векторні характеристики їх додаються геометрично, а після розходження їх ці поля існують незалежно одне від одного».

\chapter*{Процеси в ідеальному газі: одна таблиця}\addcontentsline{toc}{chapter}{Процеси в ідеальному газі: одна таблиця}\markboth{Процеси в ідеальному газі: одна таблиця}{Процеси в ідеальному газі: одна таблиця}

Теми: 01 закони, 07–09, 10, 05 C

ν = m/M молів. Перший закон: Q = ΔU + A. Для ідеального газу завжди \textbf{ΔU = $νC_{V}$ΔT}. Джерела: т. 1, § 14.7, § 16.3–16.4.

{\scriptsize
\begin{longtable}{>{\raggedright\arraybackslash}p{0.114\linewidth}>{\raggedright\arraybackslash}p{0.058\linewidth}>{\raggedright\arraybackslash}p{0.167\linewidth}>{\raggedright\arraybackslash}p{0.133\linewidth}>{\raggedright\arraybackslash}p{0.106\linewidth}>{\raggedright\arraybackslash}p{0.057\linewidth}>{\raggedright\arraybackslash}p{0.057\linewidth}>{\raggedright\arraybackslash}p{0.071\linewidth}>{\raggedright\arraybackslash}p{0.093\linewidth}}
\toprule
\textbf{процес} & \textbf{умова} & \textbf{рівняння} & \textbf{закон} & \textbf{A (газу)} & \textbf{ΔU} & \textbf{Q} & \textbf{C} & \textbf{n (pVⁿ)} \\
\midrule
\endhead
ізохорний & V = const & p/T = const & (14.28) & 0 & $νC_{V}$ΔT & = ΔU & $C_{V}$ & ±∞ \\
ізобарний & p = const & V/T = const & Гей-Люссака & p(V₂−V₁) = νRΔT & $νC_{V}$ΔT & $νC_{p}$ΔT & $C_{p}$ = $C_{V}$ + R & 0 \\
ізотермічний & T = const & pV = const & Бойля–Маріотта & νRT ln(V₂/V₁) & 0 & = A & ∞ & 1 \\
адіабатний & d'Q = 0 & $pV^{γ}$ = const, $TV^{γ−1}$ = const & Пуассона & $νC_{V}$(T₁−T₂) & −A & 0 & 0 & γ \\
політропний & C = const & pVⁿ = const & (16.40) & — & $νC_{V}$ΔT & νCΔT & C & ($C_{p−C}$)/($C_{V−C}$) \\
\bottomrule
\end{longtable}
}

\pen{} Одна діаграма p–V зі спільної точки з усіма чотирма кривими: ізохора (вертикаль), ізобара (горизонталь), ізотерма (гіпербола) й адіабата (крутіша гіпербола).

Значення: одноатомний газ γ = 5/3, двохатомний 7/5, багатоатомний 4/3 (див. 05).

\chapter*{Явища переносу: один закон, три величини}\addcontentsline{toc}{chapter}{Явища переносу: один закон, три величини}\markboth{Явища переносу: один закон, три величини}{Явища переносу: один закон, три величини}

Усі три явища переносу мають вигляд \textbf{потік = −коефіцієнт × градієнт} (англ. flux = −coefficient × gradient). У газі кожен коефіцієнт дорівнює \textbf{⅓ × (що несе одна молекула) × ⟨v⟩ × ⟨l⟩}.

{\footnotesize
\begin{longtable}{>{\raggedright\arraybackslash}p{0.184\linewidth}>{\raggedright\arraybackslash}p{0.104\linewidth}>{\raggedright\arraybackslash}p{0.239\linewidth}>{\raggedright\arraybackslash}p{0.155\linewidth}>{\raggedright\arraybackslash}p{0.165\linewidth}>{\raggedright\arraybackslash}p{0.057\linewidth}}
\toprule
\textbf{} & \textbf{потік чого} & \textbf{закон} & \textbf{градієнт} & \textbf{коефіцієнт (газ)} & \textbf{тема} \\
\midrule
\endhead
дифузія & маси & $j_{m}$ = −D dρ/dx (Фік) & густини & D = ⅓⟨v⟩⟨l⟩ & 23–25 \\
в'язкість & імпульсу & F/ΔS = −η du/dx (Ньютон) & швидкості течії & η = ⅓ρ⟨v⟩⟨l⟩ & 23–25 \\
теплопровідність & енергії & $j_{q}$ = −λ dT/dx (Фур'є) & температури & λ = ⅓ρ⟨v⟩⟨l⟩$c_{V}$ & 23–25 \\
\bottomrule
\end{longtable}
}

Той самий вигляд «потік ∝ −градієнт» має й закон з електрики: \textbf{закон Ома j = σE = −σ grad φ} (електростатика, 19). Усі чотири закони локальні, диференціальні; див. інтегральна ↔ диференціальна форма.

Спільний мікроскопічний складник — середня довжина вільного пробігу ⟨l⟩ (22).

\chapter*{Глосарій}\addcontentsline{toc}{chapter}{Глосарій}\markboth{Глосарій}{Глосарій}
\needspace{5\baselineskip}\section*{1. Електростатика}

\needspace{5\baselineskip}\subsection*{Тема 1. Електричний заряд та його властивості (§ 1.1)}

\begin{itemize}
\item \textbf{електричний заряд} — внутрішня характеристика деяких елементарних матеріальних частинок, яка проявляється в їхній взаємодії на відстані або під час дії на них зовнішніх електромагнітних полів; є як джерелом, так і об'єктом дії електромагнітного поля · $F_1/F_2 = q_1/q_2$ (1.1) · [Кл] · (англ.) electric charge · \book{} т. 2, с. 14–15
\item \textbf{електростатика} — розділ фізики, в якому вивчають властивості й взаємодію нерухомих електричних зарядів та властивості їхніх електричних полів · (англ.) electrostatics · \book{} т. 2, с. 14
\item \textbf{електростатичне поле} — поле нерухомих зарядів, характеристики якого в кожній точці не змінюються з часом (функції лише координат) · (англ.) electrostatic field · \book{} т. 2, с. 14–15
\item \textbf{позитивні / негативні заряди} — два види зарядів; позитивні виникають на скляній паличці при натиранні шовком, негативні — на ебонітовій (янтарній) при натиранні хутром; однойменні відштовхуються, різнойменні притягуються · (англ.) positive / negative charges · \book{} т. 2, с. 15
\item \textbf{електризація тертям} — перерозподіл електронів між тілами при щільному контакті; тіла набувають однакових за модулем і протилежних за знаком зарядів, сумарний заряд не змінюється · (англ.) triboelectric charging · \book{} т. 2, с. 16
\item \textbf{елементарний заряд} — найменша порція електрики, що дорівнює за модулем заряду електрона · $|e| = 1{,}602\,189\,2(46)\cdot10^{-19}$ Кл · (англ.) elementary charge · \book{} т. 2, с. 16, 18
\item \textbf{закон квантування електричного заряду} — усі заряди, які існують у природі у вільному стані, складаються з окремих дискретних зарядів, однакових за абсолютним значенням і таких, що дорівнюють зарядові електрона; приймається як постулат · $q = n|e|$, $n = \pm1, \pm2, \ldots$ · (англ.) charge quantization · \book{} т. 2, с. 18–19
\item \textbf{питомий заряд електрона} — відношення заряду частинки до її маси · $e/m = 1{,}7588\cdot10^{11}$ Кл/кг · (англ.) specific charge · \book{} т. 2, с. 17
\item \textbf{маса спокою електрона} · $m = 9{,}109\,534(47)\cdot10^{-31}$ кг · (англ.) electron rest mass · \book{} т. 2, с. 18
\item \textbf{інваріантність заряду} — незалежність числового значення заряду від стану інерціальних систем відліку (від швидкості руху носія) · (англ.) charge invariance · \book{} т. 2, с. 20
\item \textbf{замкнена система} — система, в яку з оточуючого середовища крізь її межі не проникають електрично заряджені частинки · (англ.) closed (isolated) system · \book{} т. 2, с. 20
\item \textbf{закон збереження заряду} — повний електричний заряд замкненої системи є величиною постійною і дорівнює алгебричній сумі її позитивних та негативних зарядів · $q = \sum_{i=1}^{n} q_i = \text{const}$ (1.2) · (англ.) conservation of charge · \book{} т. 2, с. 20–21
\item \textbf{об'ємна густина заряду} · $\rho = \lim_{\Delta V\to0} \Delta q/\Delta V = dq/dV$ (1.3) · [Кл/м³] · (англ.) volume charge density · \book{} т. 2, с. 21
\item \textbf{поверхнева густина заряду} · $\sigma = \lim_{\Delta S\to0} \Delta q/\Delta S = dq/dS$ (1.4) · [Кл/м²] · (англ.) surface charge density · \book{} т. 2, с. 21
\item \textbf{лінійна густина заряду} · $\tau = \lim_{\Delta l\to0} \Delta q/\Delta l = dq/dl$ (1.5) · [Кл/м] · (англ.) linear charge density · \book{} т. 2, с. 21
\end{itemize}

\needspace{5\baselineskip}\section*{2. Постійний струм}

\needspace{5\baselineskip}\section*{3. Молекулярна фізика та термодинаміка}

\needspace{5\baselineskip}\subsection*{Блок A. Основне рівняння МКТ · Закони ідеальних газів · Ступені вільності, U}

\begin{itemize}
\item \textbf{ідеальний газ} — газ, у якому можна знехтувати розмірами молекул (вважати їх матеріальними точками) та їхньою взаємодією на відстані, враховуючи лише взаємодію під час зіткнень · (англ.) ideal gas · \book{} т. 1, с. 300
\item \textbf{тепловий рух} — абсолютно хаотичний рух молекул, коли всі напрями швидкостей рівноймовірні · (англ.) thermal motion · \book{} т. 1, с. 300
\item \textbf{середня квадратична швидкість} · $\langle v_{\text{кв}}^2\rangle = (v_1^2+\ldots+v_n^2)/n$ · (англ.) root-mean-square speed · \book{} т. 1, с. 303
\item \textbf{основне рівняння молекулярно-кінетичної теорії газів} — тиск ідеального газу дорівнює двом третинам середньої кінетичної енергії поступального руху молекул одиниці об'єму газу · $p = \tfrac23 n_0 \tfrac{m_0\langle v_{\text{кв}}^2\rangle}{2}$ (14.6) · (англ.) basic equation of kinetic theory · \book{} т. 1, с. 303
\item \textbf{тиск} (газу) — усереднена за часом результуюча сила дії молекул на одиницю площі стінки; статистичний параметр · [Па = Н/м²] · (англ.) pressure · \book{} т. 1, с. 301, 304
\item \textbf{термодинамічна рівновага} — стан, у якому припиняються всі макроскопічні зміни; самовільний перехід до неї — \textbf{релаксація}, час — \textbf{час релаксації} · (англ.) thermodynamic equilibrium, relaxation · \book{} т. 1, с. 305
\item \textbf{нульове начало термодинаміки} — кожна система при фіксованих зовнішніх умовах з часом самовільно сягає стану термодинамічної рівноваги · (англ.) zeroth law · \book{} т. 1, с. 306
\item \textbf{температура} — інтенсивний параметр, що характеризує стан внутрішнього руху частинок рівноважної системи; для ідеального газу $\Theta = \tfrac23\langle\varepsilon\rangle$, $\Theta = kT$ (14.9) · [К] · (англ.) temperature · \book{} т. 1, с. 307–308
\item \textbf{стала Больцмана} — множник, що виражає співвідношення між одиницями енергії і термодинамічної температури · $k = 1{,}38\cdot10^{-23}$ Дж/К; $\tfrac{m_0\langle v_{\text{кв}}^2\rangle}{2} = \tfrac32kT$ (14.10) · (англ.) Boltzmann constant · \book{} т. 1, с. 308
\item \textbf{термічне рівняння стану} — рівняння, яке в стані теплової рівноваги пов'язує p, V, T · $F(p,V,T)=0$ (14.12); для ідеального газу $pV = nkT$ (14.14) · (англ.) thermal equation of state · \book{} т. 1, с. 311–312
\item \textbf{молярна (універсальна) газова стала} — $R = kN_A = 8{,}31$ Дж/(моль·К); чисельно дорівнює роботі розширення одного моля ідеального газу при ізобарному нагріванні на 1 К · (англ.) molar gas constant · \book{} т. 1, с. 312–313
\item \textbf{рівняння Клапейрона} · $pV_M = RT$ (14.16) · (англ.) Clapeyron equation · \book{} т. 1, с. 312
\item \textbf{рівняння Клапейрона–Менделєєва} · $pV = \tfrac mM RT$ (14.18) · (англ.) ideal gas law · \book{} т. 1, с. 312–313
\item \textbf{закон Бойля–Маріотта} (ізотермічний процес, T = const) · $pV = \text{const}$ (14.22) · (англ.) Boyle–Mariotte law · \book{} т. 1, с. 314
\item \textbf{закон Гей-Люссака} (ізобарний процес, p = const) · $V/T = \text{const}$ (14.25), $V = V_0(1+\alpha t)$ (14.26), α = 0,003661 К⁻¹ · (англ.) Gay-Lussac's law · \book{} т. 1, с. 315
\item \textbf{рівняння ізохорного процесу} (V = const) · $p/T = \text{const}$ (14.28), $p = p_0(1+\beta t)$ (14.29) · (англ.) isochoric law (Charles's law) · \book{} т. 1, с. 315–316
\item \textbf{закон Авогадро} — в однакових об'ємах ідеальних газів при однакових температурах і тисках міститься однакове число молекул · (англ.) Avogadro's law · \book{} т. 1, с. 316
\item \textbf{внутрішня енергія} — енергія, що залежить від внутрішнього стану фізичної системи; однозначна функція стану; адитивна · $U = \sum U_i$ (16.1) · (англ.) internal energy · \book{} т. 1, с. 375–376
\item \textbf{ступені вільності} (у програмі: «ступені свободи») — кількість незалежних координат, потрібних для визначення положення всіх точок системи · $i = i_{\text{пост}}+i_{\text{оберт}}+i_{\text{колив}}$ (16.4) · (англ.) degrees of freedom · \book{} т. 1, с. 376–377
\item \textbf{закон рівного розподілу енергії за класичними ступенями вільності} (закон рівного розподілу Больцмана) — у стані теплової рівноваги на кожний класичний ступінь вільності припадає в середньому кінетична енергія ½kT (коливальний — kT) · (англ.) equipartition theorem · \book{} т. 1, с. 377
\item \textbf{внутрішня енергія ідеального газу} — сумарна кінетична енергія молекул; функція тільки абсолютної температури · $U = \tfrac i2RT$ (16.3) на 1 моль · (англ.) internal energy of an ideal gas · \book{} т. 1, с. 377, 390–392
\end{itemize}

\needspace{5\baselineskip}\subsection*{Блок B. Розподіл Максвелла · Больцман, Максвелл–Больцман}

\begin{itemize}
\item \textbf{функція розподілу молекул за швидкостями} (густина ймовірності) — $dW = dn/n = f(\vec v)\,d\omega$ (14.42); умова нормування $\int f\,d\omega = 1$ (14.43) · (англ.) velocity distribution function · \book{} т. 1, с. 322
\item \textbf{закон розподілу Максвелла} (у програмі: «Максвела») — відносна кількість молекул зі швидкостями від v до v + dv · $\frac{dn}{n} = \frac4{\sqrt\pi}\left(\frac{m_0}{2kT}\right)^{3/2}e^{-m_0v^2/2kT}v^2dv$ (14.57) · (англ.) Maxwell distribution · \book{} т. 1, с. 327
\item \textbf{найбільш імовірна швидкість} · $v_\text{н} = \sqrt{2kT/m_0}$ (14.58) · (англ.) most probable speed · \book{} т. 1, с. 328
\item \textbf{середня (арифметична) швидкість} · $\langle v\rangle = \sqrt{8kT/\pi m_0}$ (14.61) · (англ.) mean speed · \book{} т. 1, с. 330
\item \textbf{середня квадратична (теплова) швидкість} · $v_{\text{кв}} = \sqrt{3kT/m_0}$ (14.62); $v_{\text{кв}} : \langle v\rangle : v_\text{н} \approx 1{,}2 : 1{,}1 : 1$ · (англ.) rms speed · \book{} т. 1, с. 318, 330
\item \textbf{барометрична формула} · $p = p_0e^{-Mg(h-h_0)/RT}$ (14.68); T = const, g = const · (англ.) barometric formula · \book{} т. 1, с. 332
\item \textbf{розподіл Больцмана} — універсальний закон розподілу частинок у будь-якому потенціальному полі в стані теплової рівноваги · $n = n_0e^{-\Delta U/kT}$ (14.71) · (англ.) Boltzmann distribution · \book{} т. 1, с. 335
\item \textbf{розподіл Максвелла–Больцмана} — добуток больцманівського й максвеллівського множників · (14.74) · (англ.) Maxwell–Boltzmann distribution · \book{} т. 1, с. 336
\end{itemize}

\needspace{5\baselineskip}\subsection*{Блок C. Довжина вільного пробігу · Явища переносу}

\begin{itemize}
\item \textbf{середня довжина вільного пробігу} — середня відстань, яку проходить молекула між двома послідовними зіткненнями · $\langle l\rangle = 1/(\sqrt2\pi d^2n_0)$ (15.2) · (англ.) mean free path · \book{} т. 1, с. 343–344
\item \textbf{середнє число зіткнень} · $\langle z\rangle = \sqrt2\pi d^2n_0\langle v\rangle$ (15.1) · (англ.) collision frequency · \book{} т. 1, с. 344
\item \textbf{ефективний поперечний переріз} (повний) · $\sigma = \pi d^2$; характеризує ймовірність процесу, а не геометрію · (англ.) effective cross-section · \book{} т. 1, с. 344
\item \textbf{явища переносу} — зумовлені тепловим рухом молекул явища перенесення маси (дифузія), імпульсу (внутрішнє тертя) і енергії (теплопровідність) · (англ.) transport phenomena · \book{} т. 1, с. 348
\item \textbf{дифузія газів} — процес самовільного вирівнювання густин (концентрацій) у суміші, зумовлений тепловим рухом мікрочастинок; в однокомпонентному середовищі — \textbf{самодифузія} · (англ.) diffusion, self-diffusion · \book{} т. 1, с. 356
\item \textbf{закон Фіка} · $j_m = -D\,d\rho/dx$ (15.22); $D = \tfrac13\langle v\rangle\langle l\rangle$ (15.27) · [м²/с] · (англ.) Fick's law · \book{} т. 1, с. 356–357
\item \textbf{внутрішнє тертя (в'язкість)} — гальмівні сили між шарами газу, зумовлені обміном імпульсом впорядкованого руху · закон Ньютона $F = -\eta\frac{du}{dx}\Delta S$ (15.10); $\eta = \tfrac13\rho\langle v\rangle\langle l\rangle$ (15.14) · [Па·с] · (англ.) viscosity · \book{} т. 1, с. 348–350
\item \textbf{теплопровідність} — молекулярне перенесення енергії в суцільному середовищі, зумовлене градієнтом температури · закон Фур'є $j_q = -\lambda\,\text{grad}\,T$ (15.16); $\lambda = \tfrac13\rho\langle v\rangle\langle l\rangle c_V$ (15.21) · [Вт/(м·К)] · (англ.) thermal conductivity · \book{} т. 1, с. 352–354
\end{itemize}

\needspace{5\baselineskip}\subsection*{Блок D–E. U · Q, A, C · I закон · C газів · Ізопроцеси · Адіабата}

\begin{itemize}
\item \textbf{термодинаміка} — наука про найбільш загальні властивості макроскопічних систем у стані термодинамічної рівноваги та про процеси переходу між цими станами · (англ.) thermodynamics · \book{} т. 1, с. 371
\item \textbf{термодинамічна система}; ізольована / закрита / відкрита / адіабатна · (англ.) thermodynamic system · \book{} т. 1, с. 371, 373
\item \textbf{квазістатичний (квазірівноважний) процес} — ідеалізований процес із неперервної послідовності рівноважних станів · (англ.) quasistatic process · \book{} т. 1, с. 372
\item \textbf{функція стану / функція процесу} — U, S — функції стану (dU, dS); A, Q — функції процесу (d'A, d'Q) · (англ.) state function / path function · \book{} т. 1, с. 375, 378–380
\item \textbf{калоричне рівняння стану} — залежність внутрішньої енергії від параметрів стану · U = U(V, T) · (англ.) caloric equation of state · \book{} т. 1, с. 376
\item \textbf{робота газу} — макроскопічна міра зміни внутрішньої енергії у квазістатичних процесах · $d\text{'}A = p\,dV$ (16.5), $A = \int p\,dV$ (16.6) · (англ.) work · \book{} т. 1, с. 378
\item \textbf{теплообмін, кількість теплоти} — передача енергії від більш нагрітого тіла до менш нагрітого без виконання макроскопічної роботи; Q — її міра · [Дж] · (англ.) heat transfer, heat · \book{} т. 1, с. 380
\item \textbf{теплоємність} · $C_0 = d\text{'}Q/dT$ (16.7); \textbf{питома} c: $C_0 = mc$ (16.8); \textbf{молярна} C = Mc (16.9) · (англ.) heat capacity · \book{} т. 1, с. 380–381
\item \textbf{перший закон термодинаміки} — у теплових процесах зміна внутрішньої енергії відбувається за рахунок приросту кількості теплоти і виконаної роботи · $dU = d\text{'}Q - d\text{'}A$ (16.10), $Q = \Delta U + A$ (16.12) · (англ.) first law of thermodynamics · \book{} т. 1, с. 381–382
\item \textbf{вічний двигун першого роду} — машина, що виконувала б роботу без одержання енергії ззовні; неможливий · (англ.) perpetual motion machine of the first kind · \book{} т. 1, с. 382
\item \textbf{рівняння Майєра} · $C_p = C_V + R$ (16.19) · (англ.) Mayer's relation · \book{} т. 1, с. 383
\item \textbf{показник адіабати (коефіцієнт Пуассона)} · $\gamma = C_p/C_V$ · (англ.) adiabatic index · \book{} т. 1, с. 384
\item \textbf{ізохорний процес} (у програмі: «ізохоричний») · $A = 0$, $Q = \Delta U$ (16.26)–(16.28) · (англ.) isochoric process · \book{} т. 1, с. 386–387
\item \textbf{ізобарний процес} · $A = p(V_2-V_1) = \nu R\Delta T$ (16.29), $Q = \nu C_p\Delta T$ (16.30) · (англ.) isobaric process · \book{} т. 1, с. 387
\item \textbf{ізотермічний процес} · $\Delta U = 0$, $Q = A = \nu RT\ln(V_2/V_1)$ (16.32) · (англ.) isothermal process · \book{} т. 1, с. 387–388
\item \textbf{адіабатний процес} (у програмі: «адіабатичний») — процес без теплообміну з навколишнім середовищем · $d\text{'}A = -dU$ (16.20) · (англ.) adiabatic process · \book{} т. 1, с. 383–384
\item \textbf{рівняння Пуассона} · $TV^{\gamma-1} = \text{const}$ (16.23), $pV^\gamma = \text{const}$ (16.23a) · (англ.) Poisson's (adiabat) equation · \book{} т. 1, с. 384
\item \textbf{політропний процес} — процес в ідеальному газі зі сталою теплоємністю · $pV^n = \text{const}$ (16.40), $n = (C_p - C)/(C_V - C)$ (16.37) · (англ.) polytropic process · \book{} т. 1, с. 388–389
\end{itemize}

\needspace{5\baselineskip}\subsection*{Блок F. Колові процеси · Карно · II закон · Ентропія · Нернст}

\begin{itemize}
\item \textbf{оборотний процес} — процес, при якому система, проведена в зворотному напрямі, повертається в початковий стан через ті самі проміжні стани, а стан тіл поза системою залишається незмінним · (англ.) reversible process · \book{} т. 1, с. 397
\item \textbf{необоротний процес} — процес, який не можна здійснити у зворотному напрямі без відповідних змін у системі або оточенні; всі реальні процеси — необоротні · (англ.) irreversible process · \book{} т. 1, с. 397
\item \textbf{цикл (коловий процес)} — процес, у результаті якого система повертається в початковий стан; робота = площа циклу на p–V · (англ.) cycle · \book{} т. 1, с. 398
\item \textbf{прямий / обернений цикл} — цикл теплової машини (за годинниковою стрілкою) / холодильної машини (проти) · (англ.) direct / reverse cycle · \book{} т. 1, с. 399
\item \textbf{цикл Карно} — оборотний цикл з двох ізотерм і двох адіабат; цикл ідеальної теплової машини · $\eta = 1 - Q_2/Q_1$ (16.53) $= 1 - T_2/T_1$ (16.64) · (англ.) Carnot cycle · \book{} т. 1, с. 399–401
\item \textbf{холодильний коефіцієнт} · $\varepsilon = Q_2/A$ (16.68) · (англ.) coefficient of performance (refrigerator) · \book{} т. 1, с. 403
\item \textbf{коефіцієнт перетворення (теплового насоса)} · $\varphi = Q_1/A = T_1/(T_1-T_2)$ (16.67) · (англ.) heat pump COP · \book{} т. 1, с. 402–403
\item \textbf{другий закон термодинаміки} — Клаузіус: теплота не може сама собою переходити від менш нагрітого тіла до більш нагрітого; Томсон: неможлива періодична машина, що перетворює теплоту в роботу лише охолодженням одного тіла · (англ.) second law of thermodynamics · \book{} т. 1, с. 403–404
\item \textbf{вічний двигун другого роду} — неможливий (Оствальд) · (англ.) perpetual motion machine of the second kind · \book{} т. 1, с. 404
\item \textbf{теорема Карно} — ККД оборотного циклу Карно не залежить від робочого тіла; ККД необоротного циклу завжди менший · $\eta \le \eta_0$ (16.69) · (англ.) Carnot's theorem · \book{} т. 1, с. 404
\item \textbf{зведена теплота}; \textbf{нерівність Клаузіуса} · $Q_1/T_1 + Q_2/T_2 \le 0$ (16.71), $\oint d\text{'}Q/T \le 0$ (16.74) · (англ.) reduced heat, Clausius inequality · \book{} т. 1, с. 404–405
\item \textbf{ентропія} — функція стану, повний диференціал якої при оборотному процесі $dS = d\text{'}Q/T$ (16.75); адитивна; закон зростання $dS \ge 0$ для ізольованої системи (16.82) · [Дж/К] · (англ.) entropy · \book{} т. 1, с. 406–408
\item \textbf{термодинамічна ймовірність} — число мікростанів, які реалізують даний макростан · (16.90) · (англ.) thermodynamic probability (statistical weight) · \book{} т. 1, с. 411–412
\item \textbf{принцип Больцмана} · $S = k\ln w$ (16.91) · (англ.) Boltzmann's entropy formula · \book{} т. 1, с. 413
\item \textbf{третій закон термодинаміки (теорема Нернста)} — при T → 0 ентропія рівноважної системи перестає залежати від параметрів і прямує до сталої (нуля); $\lim_{T\to0}S = 0$ (16.93, Планк) · (англ.) third law, Nernst theorem · \book{} т. 1, с. 414–415
\end{itemize}

\needspace{5\baselineskip}\subsection*{Блок G–H. Реальний газ · Ван-дер-Ваальс · Кристали · Дефекти · Рідини}

\begin{itemize}
\item \textbf{реальний газ} — газ, для якого суттєву роль відіграють міжмолекулярні сили взаємодії · (англ.) real gas · \book{} т. 1, с. 419
\item \textbf{сили Ван дер Ваальса} — орієнтаційні (Кеєзом), індукційні (Дебай), дисперсійні (Лондон) · (англ.) van der Waals forces · \book{} т. 1, с. 421–422
\item \textbf{рівняння Ван-дер-Ваальса} · $(p + a/V_M^2)(V_M - b) = RT$ (17.4); a/V² — \textbf{внутрішній тиск}, b — поправка на власний об'єм · (англ.) van der Waals equation · \book{} т. 1, с. 423
\item \textbf{критичний стан}; \textbf{критичні параметри} · $p_\text{к} = a/27b^2$, $V_\text{к} = 3b$, $T_\text{к} = 8a/27Rb$ (17.12) · (англ.) critical state · \book{} т. 1, с. 426
\item \textbf{метастабільні стани} — пересичена пара, перегріта рідина · (англ.) metastable states · \book{} т. 1, с. 429
\item \textbf{бінодаль / спінодаль} — межа двофазної області / межа нестійких станів · (англ.) binodal / spinodal · \book{} т. 1, с. 429–430
\item \textbf{тверде тіло} — агрегатний стан зі стабільністю форми та об'єму і впорядкованим розміщенням атомів · (англ.) solid · \book{} т. 1, с. 444
\item \textbf{дальній / ближній порядок} — строга повторюваність структури на великих відстанях (кристали) / узгодження лише сусідніх частинок (аморфні тіла, рідини) · (англ.) long-range / short-range order · \book{} т. 1, с. 444, 463
\item \textbf{анізотропія} — залежність фізичних властивостей від напряму · (англ.) anisotropy · \book{} т. 1, с. 445
\item \textbf{елементарна комірка; гратки Браве; сингонія} — 14 граток Браве, 7 сингоній · $\vec r_n = m\vec a_1 + n\vec a_2 + p\vec a_3$ (18.1) · (англ.) unit cell, Bravais lattices, crystal system · \book{} т. 1, с. 446–447
\item \textbf{точкові дефекти} — за Шотткі (вакансія), за Френкелем (вакансія + міжвузловий атом), домішки · $n_{\text{в}} = n_0e^{-W_{\text{в}}/kT}$ (18.4) · (англ.) point defects · \book{} т. 1, с. 453–455
\item \textbf{дислокація} (крайова, гвинтова); \textbf{вектор Бюргерса} · (англ.) dislocation, Burgers vector · \book{} т. 1, с. 455–456
\item \textbf{рідина} — «розупорядковане тверде тіло» з ближнім порядком; тепловий рух за Френкелем: коливання + стрибки · $\tau = Ae^{\Delta W/kT}$ (19.1) · (англ.) liquid · \book{} т. 1, с. 462–464
\item \textbf{формула Френкеля (в'язкість рідин)} · $\eta = Ae^{\Delta W/kT}$ (19.7), спадає з T · (англ.) Frenkel viscosity formula · \book{} т. 1, с. 467
\end{itemize}


\end{document}
